\documentclass[UTF8,fontset=fandol,zihao=-4,oneside,openany]{ctexbook}
\usepackage[a4paper,inner=26mm,outer=24mm,top=25mm,bottom=25mm]{geometry}
\usepackage{amsmath,amssymb,booktabs,longtable,array,xcolor,graphicx}
\usepackage{tikz,pgfplots}
\pgfplotsset{compat=1.18}
\usepackage{hyperref}
\hypersetup{unicode=true,colorlinks=true,linkcolor=teal,citecolor=teal,urlcolor=teal,pdftitle={新时代资本论：2026年中国大陆的生产、分配与风险},pdfauthor={姚林翔宇}}
\setlength{\emergencystretch}{3em}
\linespread{1.28}
\newcommand{\Chart}[3]{\begin{figure}[htbp]\centering\csname chart@#1\endcsname\caption{#2}\par\small #3\end{figure}}
% CHART_DEFINITIONS_BEGIN
\pgfplotsset{bookplot/.style={width=.88\textwidth,height=55mm,axis lines=left,ymajorgrids=true,grid style={gray!18},tick label style={font=\small},label style={font=\small},xtick=data,enlarge x limits=.22,scaled y ticks=false,yticklabel style={/pgf/number format/1000 sep={}},bar width=15pt,nodes near coords,nodes near coords style={font=\footnotesize},every axis plot/.append style={fill=teal!65,draw=teal!85!black}}}
\expandafter\def\csname chart@sector\endcsname{\begin{tikzpicture}\begin{axis}[bookplot,ybar,ymin=0,ymax=70,ylabel={占比（\%）},xticklabels={第一产业,第二产业,第三产业}]\addplot coordinates {(1,6.7) (2,35.6) (3,57.7)};\end{axis}\end{tikzpicture}}

\expandafter\def\csname chart@prices\endcsname{\begin{tikzpicture}\begin{axis}[bookplot,ybar,ymin=-1.5,ymax=9,ylabel={同比（\%）},xticklabels={CPI,PPI,购进价格,生产资料,生活资料}]\addplot coordinates {(1,0.9) (2,5.0) (3,6.9) (4,6.5) (5,-0.1)};\end{axis}\end{tikzpicture}}

\expandafter\def\csname chart@migrantpay\endcsname{\begin{tikzpicture}\begin{axis}[bookplot,ybar,ymin=0,ymax=7000,ylabel={元／从业月},xticklabels={外出农民工,本地农民工}]\addplot coordinates {(1,5774) (2,4376)};\end{axis}\end{tikzpicture}}

\expandafter\def\csname chart@credit\endcsname{\begin{tikzpicture}\begin{axis}[bookplot,ybar,ymin=-3,ymax=14,ylabel={万亿元},xticklabels={住户,企事业单位}]\addplot coordinates {(1,-1.03) (2,11.26)};\end{axis}\end{tikzpicture}}

\expandafter\def\csname chart@fiscal\endcsname{\begin{tikzpicture}\begin{axis}[bookplot,ybar,ymin=-36,ymax=12,ylabel={同比（\%）},xticklabels={地方一般预算,土地出让,全国付息}]\addplot coordinates {(1,3.1) (2,-28.6) (3,5.4)};\end{axis}\end{tikzpicture}}

\expandafter\def\csname chart@care\endcsname{\begin{tikzpicture}\begin{axis}[bookplot,ybar,ymin=0,ymax=250,ylabel={分钟／日},xticklabels={男性参与者,女性参与者}]\addplot coordinates {(1,112) (2,209)};\end{axis}\end{tikzpicture}}

\expandafter\def\csname chart@income\endcsname{\begin{tikzpicture}\begin{axis}[bookplot,ybar,ymin=0,ymax=36500,ylabel={元／人（半年）},xticklabels={全国平均,全国中位,城镇平均,农村平均}]\addplot coordinates {(1,22981) (2,19036) (3,30126) (4,12699)};\end{axis}\end{tikzpicture}}

\expandafter\def\csname chart@classflow\endcsname{\begin{tikzpicture}[every node/.style={draw=teal!65,rounded corners,align=center,text width=32mm,minimum height=23mm,font=\small}]
\node (a) at (0,0) {资源与权利\\公共权限、资产\\劳动与技能};
\node (b) at (4.4,0) {制度传递\\合同、保障\\家庭与公共服务};
\node (c) at (8.8,0) {实际处境\\收入、选择空间\\损失与恢复能力};
\draw[->,thick] (a)--(b);\draw[->,thick] (b)--(c);
\end{tikzpicture}}

\expandafter\def\csname chart@trade\endcsname{\begin{tikzpicture}\begin{axis}[bookplot,ybar,ymin=0,ymax=28,ylabel={万亿元人民币},xticklabels={货物出口,货物进口,货物差额}]\addplot coordinates {(1,23.2) (2,16.74) (3,6.46)};\end{axis}\end{tikzpicture}}

\expandafter\def\csname chart@rights\endcsname{\begin{tikzpicture}[every node/.style={draw=teal!65,rounded corners,align=center,text width=29mm,minimum height=15mm,font=\small}]
\node (a) at (0,0) {法律归属\\公共资产};\node (b) at (3.5,0) {治理与用途\\投资、服务};\node (c) at (7,0) {分配制度\\预算、定价};\node (d) at (10.5,0) {居民实际获得\\现金、质量、保障};
\draw[->,thick] (a)--(b);\draw[->,thick] (b)--(c);\draw[->,thick] (c)--(d);
\end{tikzpicture}}

\expandafter\def\csname chart@risk\endcsname{\begin{tikzpicture}\begin{axis}[bookplot,ybar stacked,nodes near coords={},ymin=0,ymax=12,ylabel={损失单位},xtick={1,2},xticklabels={初始分担,最终分担},legend style={at={(.5,-.2)},anchor=north,legend columns=4,font=\small,draw=none},enlarge x limits=.7]
\addplot+[fill=teal!80] coordinates {(1,4)(2,4)};\addplot+[fill=blue!45] coordinates {(1,3)(2,1)};\addplot+[fill=orange!60] coordinates {(1,2)(2,2)};\addplot+[fill=gray!50] coordinates {(1,1)(2,3)};\legend{股东,银行,供应商,财政}
\end{axis}\end{tikzpicture}}

\expandafter\def\csname chart@debt\endcsname{\begin{tikzpicture}\begin{axis}[bookplot,ymin=85,ymax=120,xlabel={期数},ylabel={债务／收入（\%）},xtick={0,1,2,3,4,5},enlarge x limits=.05,nodes near coords={},legend style={at={(.02,.97)},anchor=north west,legend columns=1,font=\small,draw=none}]
\addplot+[fill=none,thick,mark=*] coordinates {(0,100.000000) (1,98.095238) (2,96.226757) (3,94.393867) (4,92.595888) (5,90.832157)};\addplot+[fill=none,draw=orange!80!black,thick,mark=square*] coordinates {(0,100.000000) (1,102.941176) (2,105.968858) (3,109.085589) (4,112.293989) (5,115.596753)};\legend{名义增长5\%、利率3\%,名义增长2\%、利率5\%}
\end{axis}\end{tikzpicture}}
% CHART_DEFINITIONS_END
\begin{document}
\frontmatter
\hypersetup{pageanchor=false}
\begin{titlepage}
\centering\vspace*{28mm}
{\Huge\bfseries 新时代资本论\par}
\vspace{12mm}
{\LARGE 2026年中国大陆的\par 生产、分配与风险\par}
\vspace{22mm}
{\large 原创研究专论\quad 扩充修订版\par}
\vspace{8mm}
{\large 作者：姚林翔宇\par}
\vfill
{\large 资料截止：2026年10月14日\par}
\vspace{8mm}
{\small 以历史与制度分析为起点，以可追溯证据和可反驳命题为约束\par}
\end{titlepage}
\hypersetup{pageanchor=true}
\chapter*{阅读说明}
\addcontentsline{toc}{chapter}{阅读说明}
本书是围绕当代中国大陆经济的原创研究专论，不是马克思原著的新版译本，也不代表某种已经得到学界公认的理论替代。沿用“资本论”这一问题意识，是要把生产、所有权、分配、金融和社会生活重新放在同一个分析框架内，而不是为任何既定结论寻找权威。

本书区分五种陈述：统计机构发布的观察值、研究者的估计、机构的预测、本文的机制推论，以及用来说明逻辑的假设例子。前三者分别注明来源、样本或预测时点；后两者不冒充调查结果。所有“2026年”的现实描述均受资料截止日约束，全年结果尚未形成。正文中的历史研究只在其覆盖范围内使用，不能因为出版物名称带有“2026”，就把它引用的2024年数据改成2026年事实。

研究边界为中国大陆。国际资料中的China系列须检查是否单列香港、澳门和台湾；无法确认同口径的序列不拼接。官方统计有不可替代的覆盖优势，也有口径修订和观察盲区；国外资料有方法上的补充价值，也可能沿用同一批国内基础数据。两种材料同时出现，不自动构成两份独立证据。

本文没有开展新问卷调查，没有取得受限的家庭微观数据库，也没有运行新的全国因果回归。可复核计算包括公开数字的比例、差额及明确标注的情景算例。理论贡献以定义澄清、部门关系的一致核算和可检验命题为主。证据不足的地方保留空白，避免用一套过于整齐的故事覆盖复杂现实。

\tableofcontents
\mainmatter
\chapter*{导论：为什么需要重新研究资本}
\addcontentsline{toc}{chapter}{导论：为什么需要重新研究资本}
\section{从一个普通家庭的问题出发}
一个家庭的生活改善，通常不取决于它能否背出国内生产总值增长率。它关心的是工作能持续多久，工资能购买多少东西，房贷是否压缩生活，老人失能后由谁照料，孩子读书和迁移是否需要重复支付成本。一个企业则关心订单、回款、融资、技术和用工。地方政府面对财政收入、公共服务、产业发展和存量债务。银行需要在贷款安全、盈利、流动性和政策要求之间选择。这些问题相互连接，却经常在不同学科和不同统计表中被分开叙述。

% ARGUMENT P001
\phantomsection\label{arg:P001}
把这些问题放在一起，有一个实际理由：同一笔收入在不同组织中承担不同作用。工资在企业账上是成本，在家庭账上是生活来源，在地方商业中又成为需求。企业为了渡过订单不足而减少用工，可能改善自身现金流，却使周边商户失去顾客。商户随后缩减采购，又会成为别家企业的订单压力。这个过程并非任何一个人有意安排，而是相互依赖的交易共同造成的。

分析的起点因此是可以追踪的关系。假设一个家庭遇到老人需要长期照护的情况，它可以购买服务、减少工作时间或者请亲属帮忙。选择取决于服务供给、工资、保险和亲属距离。不同选择分别改变服务企业的收入、雇主的劳动力供给和亲属家庭的时间安排。只有记录这些去向，才能解释同样一次健康冲击为什么在有些家庭中只是额外开销，在另一些家庭中却会中断整个职业轨迹。这是贯穿本书的情景分析，不是对某一真实家庭的调查记录。

重新研究资本，要从这种连接开始。资本不是只存在于富人的账户里。购买设备的贷款、地方基础设施的债券、居民为住房承担的按揭、平台对商家的结算规则，都影响当期的生产组织和未来的收入要求。一个部门认为安全的资产，可能是另一个部门需要兑现的负债。一个组织提高了利润，也可能同时增加了家庭的无酬劳动、供应商的垫资或未来财政的负担。把这些联系看清楚，才知道经济增长改变了什么。

% ARGUMENT P002
\phantomsection\label{arg:P002}
这里还需要区分资本关系与所有金融交易。家庭购买自住住房，并不因此成为以雇佣劳动和积累为主要目标的企业；它却进入了由土地权利、开发融资、按揭和资产定价构成的关系网络。金融合同把不同主体的现在和未来连接起来，使家庭的劳动收入成为偿债承诺的来源。研究资本怎样影响家庭，并不要求把家庭的全部生活都重新命名为资本活动。

反过来，不能因为某种成本没有进入企业会计，就断定它没有发生。员工下班后恢复体力所需的时间、照料孩子的安排、污染对附近居民健康的影响，都可能影响生产能否持续。但指出账外成本，仅仅完成了提出问题的一步。还必须辨认成本是否由该企业活动引起、原先由谁承担、改变安排后是否真的减少，以及能否用合理的制度使承担者得到补偿。把所有不幸都归于资本，与只看企业报表一样，会使具体机制消失。

马克思的一个持久贡献，是要求人们把市场交易放回生产和社会关系中理解。劳动合同并不只是一项价格交换；它也安排指挥、监督、剩余索取和退出。商品之间的交换，也依赖看不见的组织、信用和制度。这个研究方向仍然有力量。但是，研究方向的有效，不等于原著中的每一个具体命题都无须重新证明。劳动投入与价格之间的关系、利润来源的分解、阶级结构的表达和危机发生的条件，都必须接受现代证据的检验。\cite{Marx1,Marx2,Marx3}

% ARGUMENT P003
\phantomsection\label{arg:P003}
历史分析首先要求承认研究对象已经变化。十九世纪工厂中的业主、工人与产品市场，并不能完整代表当代企业集团、公共资本、银行体系和家庭保障。与此同时，劳动者必须在收入中断之前找到新工作、管理者可以规定工作任务、资产持有者可以要求未来现金流，这些关系并没有因新技术出现而自动消失。历史性的含义，是辨认哪些关系延续、哪些条件改变，而不是以年代远近替代理论判断。

本书对原典的重审也需要公平。马克思并没有把每一次成交价格都等同于劳动价值；第三卷讨论平均利润和生产价格，就是对竞争中价值分配形式的进一步说明。真正需要检验的是各层次之间怎样连接，连接需要哪些假设，现代资料能观察到其中什么。把复杂理论压缩成一句容易否定的话，无法构成更新。第二章将先重建问题，再讨论可以保留与需要修订的解释。\cite{MarxPrice}

\section{2026年不是一张颜色单一的照片}
观察当代中国，不能只选一种符合预期的指标。国家统计局初步核算的2025年GDP为1401879亿元，实际增长5.0\%；三次产业增加值占比分别为6.7\%、35.6\%和57.7\%。这说明生产已形成农业、工业和广泛服务活动共同构成的结构。企业生产、家庭生活和公共服务之间的联系，远超十九世纪工厂的典型图景。\cite{NBS2025}

% ARGUMENT P004
\phantomsection\label{arg:P004}
总产出增长说明经济活动的规模变化，并不直接给出一个家庭在其中的位置。生产可以更多发生在资本密集的行业，也可以发生在就业吸纳较强的服务部门；新增收入可以进入工资、企业留存、财政或者金融索取权。即使这些去向都合法，居民当期可用的收入和生活压力也可能很不相同。因此，从总量进入分配，必须经过行业结构、要素报酬和家庭获得三个层次。

产业份额还受相对价格影响。某种工业品价格下降时，即使产量增加，它按现价计算的增加值份额也可能下降。服务业份额上升同样不能单独证明生产已经脱离实体活动，运输、维修、护理和软件都可能直接维持生产与生活。本书使用产业构成来描述结构，而不把某一产业占比当作经济先进或落后的独立尺度。真正的比较应考察质量、效率、就业条件和资源消耗，并在可能时分开价格变化与数量变化。

2026年上半年居民人均可支配收入的平均数为22981元，中位数为19036元。它们是全体居民口径，不是在职员工的平均工资，也不是典型家庭的年收入。两个数的差别首先提醒我们观察分布，不能从全国平均增速推断每个家庭的境况。相同工资还可能对应不同家庭人数、住房债务和保障条件。\cite{Income2026}

% ARGUMENT P005
\phantomsection\label{arg:P005}
平均数由总收入除以人数得到，中位数则描述排序中间的位置。两者之间的距离可以提示分布偏斜，但不能告诉我们最高收入群体占有多少，也不能显示收入最低者的生活缺口。若想从这两个数字推导出全国阶层结构，就会把统计没有提供的信息擅自补入。第八章因此将收入排序与社会位置分开处理，再研究两者怎样相互影响。

家庭人数与保障条件也会改变比较。两个收入总额相同的家庭，如果一个只有两名健康成年人，另一个需要供养孩子和失能老人，其可支配时间与必要开支就不同。简单按人数平均又不能完全反映共同居住的规模经济。本文不自行规定一个全国统一的生活等价系数，而是在讨论中说明家庭构成，必要时同时比较人均现金收入、不可避免支出和公共服务获得。这样可以避免把照护负担误读为消费偏好，也避免把资产支持误读为个人工资能力。

时间变化同样不能忽略。IMF在2026年2月发布的中国评估强调房地产调整、内需和价格压力；世界银行7月报告则讨论高技术投资和出口支持下的增长与再平衡。两份报告都是有时点的判断，不是全年不变的说明书。到10月14日发布的9月数据，CPI同比上涨0.9\%，PPI同比上涨5.0\%，工业购进价格同比上涨6.9\%。因此，本书不把当下写成“所有价格都在下跌”。上游涨价、部分终端需求较弱与特定资产价格调整可以同时发生。\cite{IMF2026,WB2026,CPI2026,PPI2026}

% ARGUMENT P006
\phantomsection\label{arg:P006}
不同发布日期的资料承担不同任务。年度公报适合描述较完整的结构，月度指标适合捕捉变化方向，专题调查适合观察总量表看不到的家庭活动。它们不能因为都标有某个年份，就被当作同一时点的照片。例如一份报告在七月讨论增长条件，其中部分基础数据可能来自上一年度；九月的价格变化也不能追溯改变那份报告在七月能够知道的信息。

在机制分析中，上游投入涨价可能压缩下游企业利润，也可能通过提价转给消费者，或者通过减少用工转给劳动者。哪条路径占主导，取决于库存、合同期限、替代材料和市场竞争。本文不从两个价格指数的差额直接计算全国企业利润损失，因为采购篮子与销售篮子并不相同。可以据此提出成本传递问题，却必须进一步观察行业和企业资料。资料时点与价格口径的限制，决定了推论能够走多远。

这种并存不是逻辑矛盾。经济由不同部门构成，价格指数覆盖不同商品，名义收入与真实生活水平也不是一回事。如果分析只允许繁荣或衰退两个词，就会漏掉最值得解释的结构：谁的收入增加，谁的成本增加，谁拥有调整合同的能力，谁只能接受结果。

% ARGUMENT P007
\phantomsection\label{arg:P007}
这种结构性观察可以具体化为一组连续问题：冲击最先发生在哪个市场，哪些合同允许及时调价，哪些支付必须固定履行，谁能等待，谁必须立即变现。假设原材料涨价，拥有长期销售合同的生产商可能暂时无法提高售价，而拥有稳定库存的竞争者受到的冲击较小；同一行业也会出现不同结果。若企业融资便利，它可能吸收一段时期的成本；若必须靠当月回款发工资，调整可能更快落到劳动者身上。

因此，阶层后果不是宏观指标旁边附加的一段评论，而是传导机制的一部分。拥有资产缓冲的人，可以把价格波动视作等待机会；缺少储蓄的家庭，可能在相同波动中被迫接受更差的劳动条件。这不意味着前者永远获益、后者永远受损。资产集中在下跌行业的持有者也会受损。需要解释的，是风险暴露与应对能力的组合，而不是预先为每个群体安排固定结局。

\section{四个问题与三本账}
本书围绕四个问题展开：财富怎样生产，资源与决策由谁控制，收益怎样分配，风险最终由谁承担。为避免概念滑动，至少要同时记录三本账。

% ARGUMENT P008
\phantomsection\label{arg:P008}
这四个问题的次序也有意义。没有真实生产与服务，收益只能在既有财富和索取权之间重新分配；没有决策规则，生产可能无法协调；没有分配规则，生产结果难以进入生活；没有风险安排，今天的承诺可能在明天失效。不过，四个环节并不是只向前运行的直线。分配影响需求和劳动选择，风险预期影响投资，资源控制又受到过去收益积累的改变。

由此形成的是一个动态研究对象。某项制度在起步时期帮助集中资源，可能随着规模扩大而产生监督困难；某种激励在需求旺盛时期推动创新，也可能在市场接近饱和时鼓励重复投资。本文不会只问一种制度在抽象意义上好不好，而要问它在什么阶段、面对什么约束、由什么组织执行，以及改变后谁的行为会首先变化。这使制度评价能够与历史过程连接，而不依赖一个跨越所有情境的排名。

第一本是生产账。它记录劳动、设备、知识、自然条件和组织如何形成商品与服务。生产账不能把收入转移当作新产出，也不能因为一项活动没有市场价格就断言它没有社会作用。第二本是权利与收入账。它记录谁可以决定使用资产、改变合同、收取费用、分享利润和要求偿债。第三本是再生产与风险账。它记录家庭、公共部门和生态条件怎样被维持，以及失败时损失如何分摊。

% ARGUMENT P009
\phantomsection\label{arg:P009}
三本账的作用，在于避免一种改善掩盖另一种恶化。生产账显示新建了多少服务能力，权利与收入账显示使用资格、价格和收入归属，再生产与风险账显示这些能力在冲击中能否维持。一个托育机构增加床位，若收费超出目标家庭承受能力，或者营业时间与父母工作错开，生产能力的账面增加就没有充分转化为家庭可用服务。

这里没有一个可以把所有结果无争议地加总的货币单位。健康损失、照护时间和财务亏损可以分别测量，但货币化需要额外假设。本书优先保留并列指标，再说明它们之间的关系。若必须做成本收益比较，应公开时间价值、折现和风险权重，并说明改变这些假设是否改变政策排序。让不同账目相互约束，比把所有东西折成一个看似精确的数更能暴露真正的取舍。

三本账可以相互解释，不能随意相加。工厂支付给工人的工资属于生产形成的收入分配，不应在已经核算的增加值之外再加一遍；工人消费住房服务时，该项服务按自身生产边界核算，而不是按钱流转的次数增加产出。政府提供教育需要资源，但把教育支出和居民获得的同一项服务再算一次，仍然是重复计算。家庭照料可以另设时间账和卫星账户，不能把估算的家务价值无条件叠加到官方GDP之上。

% ARGUMENT P010
\phantomsection\label{arg:P010}
这里需要进一步澄清租赁服务。支付房租本身是支出与收入的对应，并不按支付次数创造价值；但实际提供当期住房使用服务，属于国民核算中的服务产出，相关维护和管理也可能创造增加值。不能因为租金具有收入转移和权利索取的性质，就把全部住房服务从生产账户中抹去。关键在于分开服务的生产、资产所有权产生的索取以及旧资产的价格重估。

同理，政府购买服务与居民享受服务，是同一供给从支付端和使用端观察的两面，不应重复计算。补贴只是改变支付责任，是否增加真实供给要看资源有没有投入。当家庭照护转为市场服务时，统计产出可能增加，即使总照护时间变化不大；这既不是统计必然虚假，也不代表福利必然同比增加。评价还需比较照护质量、家庭释放的时间和新增财政或现金负担。

资产账还存在价格变化。旧房价格上涨，可以增加持有者按市价衡量的财富，却不意味着当期多生产了一套房。价格上涨带来的抵押能力，可能进一步影响投资和消费，这才是资产账连接生产账的机制。只有把“发生了重估”与“发生了生产”区分开，财富研究才不会被账面数字迷惑。

% ARGUMENT P011
\phantomsection\label{arg:P011}
旧资产升值之所以能影响现实生产，是因为金融机构和家庭可能根据其估值改变行为。抵押品价值上升时，借款人可能获得更大融资额度；卖出资产的人可以增加消费，仍持有资产的人也可能降低储蓄。但这些行为都不是重估本身，而是重估经过融资规则、心理预期和交易发生以后形成的结果。若资产无法出售、银行不接受抵押，账面财富就未必带来可用资金。

对于自住家庭，房价上涨还可能同时提高下一次置换住房的成本。只有卖出后不再购买类似住房，或者搬到价格更低的地区，部分升值才能真正转成其他消费。对长期租房者，周边房价上涨则可能意味着租金压力或未来购房门槛提高。把所有居民视为同一个资产持有部门，会掩盖这些方向相反的影响，也会错误地把资产价格上涨当作普遍的生活改善。

\section{方法的保留与命题的重审}
本文保留三个方法原则。第一，制度具有历史性，不能把某一时期的财产权利、劳动组织和财政安排当作永恒自然事实。第二，个人选择嵌入约束条件；解释劳动者接受低工资时，应同时研究他能否迁移、获得保障和拒绝不利合同。第三，整体分析必须经过部门和机制，不能用“资本”“国家”或“人民”这样的总称代替具体关系。

% ARGUMENT P012
\phantomsection\label{arg:P012}
制度的历史性，可以通过进入与退出成本具体理解。一家企业在过去形成的供应网络、一个地方已经建成的基础设施、一户家庭已经承担的按揭，都使未来选择带有路径依赖。即便新的安排在长期更有效，立即转换也可能造成设备闲置、技能失配和生活中断。历史分析的任务，是解释这些沉没投入怎样约束当下，而不是把既有安排都视为不可改变。

个人选择也具有真实能动性。约束并不意味着每个人必然采取同一行动，而是某些选择的代价对不同人并不相等。对储蓄充足的专业人员而言，拒绝一份不合适的工作可能只是延后收入；对没有保障的家庭而言，同样拒绝可能意味着下月无法支付房租。研究两者选择差异时，不能只根据最终行为评判耐心、能力或努力。应先比较他们面对的机会集合，再观察个人偏好和判断怎样起作用。

同时，本文拒绝三种捷径。不能用劳动创造物质产品，直接证明所有市场价格都与劳动时间成比例；不能用收入差距存在，直接证明其中每一部分都是同一种剥夺；也不能用一个理论能够讲通过去，就宣布它已经准确预测未来。理论需要解释不符合直觉的事实，并说明什么观察会使自己失败。

% ARGUMENT P013
\phantomsection\label{arg:P013}
所谓反证，并不是要求每个社会现象都能用一次实验解决，而是要求理论承担出错的风险。如果把利润增加解释为资本力量增强，把利润减少又解释为资本为了未来布局，那么无论出现什么事实，理论都不会失败。这类解释缺少可区分的内容。更可用的命题应说明，在其他条件相近时，哪些控制权变化应当影响哪类结果，出现什么相反变化就需要修正。

宏观研究还必须区分恒等式和因果关系。债务余额由过去余额、借款和偿还共同变化，这是账户关系；借款增加会不会提高收入，则取决于资金用途、供给条件与预期。把恒等式中的某一项移到等号另一边，不能自动产生政策结论。本文因此将“按定义成立”“在模型条件下成立”和“在现实资料中得到支持”分别使用，不以数学符号增加未经证实的确定性。

这里所谓创新，有一个较严格的标准：新增概念必须减少歧义，新增机制必须产生可以与其他解释区分的结论，新增命题必须提出可观察的检验对象。如果把财产权理论、社会再生产理论和金融不稳定理论放到中国制度中联合使用，但没有证明这种组合超出已有成果，就称为理论综合。把已有思想换一个名称，不算理论发现。

% ARGUMENT P014
\phantomsection\label{arg:P014}
三个候选命题的额外价值，要与已有解释比较。产权与收益兑现的分析，与财产权、公共财政和治理研究有直接联系；家庭再生产与生产扩张的分析，与劳动供给、社会再生产和消费理论相通；风险转移的分析，则连接软预算约束、金融不稳定和保险研究。这些思想来源都应承认。本文尝试增加的是跨账户追踪的要求，而不是宣称此前没有人研究过其中任何一个问题。

例如，公共企业分红增加以后，若只观察财政收入，就可能漏掉企业通过新增债务支付分红的情况；若再看服务价格，又可能发现居民现金负担同步上升。只有把这些结果合起来，才知道公共收益是否真正增加。但联合观察并不天然构成新理论。只有它能识别单独指标遗漏的重要机制，并产生可复核的不同预测，才有理由认为这种综合获得了额外解释力。

\section{判断经济制度时公开价值标准}
纯粹没有价值立场的政治经济学并不存在，因为“改善”必须有评价对象。本文采用可讨论的标准：生产效率、基本生活保障、机会与议价条件、风险分担和生态可持续性。它们并不总是同时改善。提高某项保障可能增加财政支出，限制垄断可能改变创新融资，房地产调整可能降低新购房者成本同时损害原有业主净资产。

% ARGUMENT P015
\phantomsection\label{arg:P015}
价值标准之间的冲突可以通过政策比较公开呈现。假设一种改革提高总体产出，但使某些家庭承担集中的失业风险。强调平均效率的人可能支持它，强调基本保障的人会要求补偿和过渡。二者争论的部分内容来自不同权重，另一些内容则来自对就业效果和补偿可行性的不同判断。研究应把这两类分歧分开，事实可以进一步调查，价值权重则需要明确讨论。

本书尤其重视“不能承受的损失”与“可以补偿的损失”的区别。同样金额的损失，对储蓄充足者可能只是投资回报下降，对脆弱家庭可能导致失学、延误治疗或永久退出劳动市场。仅比较货币总额，会低估后者的不可逆影响。但这也不意味着任何既有收入都必须保护。政策需要先确定基本生活底线，再讨论超出底线的风险由谁承担，以及保护会不会鼓励新的风险积累。

研究的责任是把这种冲突公开，而不是把政策偏好隐藏在看似中性的术语里。本文不预先认定某一种所有制总是有效，也不把增长率最高当作唯一目标。应检验同样的资源能否带来更可靠的生活、更有用的产出和更少的不可逆损害，并交代转型中的成本由谁承担。

% ARGUMENT P016
\phantomsection\label{arg:P016}
评价所有制时，至少需要比较相似任务。承担偏远地区基础服务的企业，与只服务高利润客户的企业，其收入与成本结构原本不同。若不调整服务义务就比较利润率，可能把公共任务误判为经营低效；反过来，若只要声称承担公共任务就不再检查成本，也会为低效提供无限借口。合理评价需要将任务、补偿、质量与财务结果同时公开。

对中国经济的分析同样应允许收益与代价在同一制度中并存。组织协调可能缩短基础设施建设的时间，也可能增加错误决策被复制的范围；分散竞争可能带来创新，也可能让个别企业通过拖欠和降质把成本推给别人。问题不是在两种抽象制度之间选一个永远正确的答案，而是辨认协调、竞争、监督和保障的具体组合。这些组合的效果，应由它们实际改变了谁的能力和负担来衡量。

读者因此可以不同意本文的政策排序，同时仍然使用它的数据和机制分析。一本有用的政治经济学著作，应允许这种分离。它提供可讨论的框架，不要求读者接受作者的全部价值权重。

% ARGUMENT P017
\phantomsection\label{arg:P017}
阅读全书可以沿着两条线进行。一条线从商品和劳动开始，经过企业、信用与土地，观察生产与分配怎样组织起来；另一条线从家庭和社会位置出发，追踪工作、住房、保障和照护条件怎样受到这些组织关系影响。第八章集中讨论阶层，但阶层关系并不只存在于那一章。任何涉及合同、价格和风险的机制，都有相应的社会位置。

因此，读者看到一个判断时，可以依次追问：它描述的是实际观察还是假设情景，成立需要什么条件，哪类人能够改变规则，哪类人只能调整自己的生活，结论在什么情况下不再成立。这样的阅读方式，也构成作者必须接受的检验。一本研究专著的完整性，不在于所有问题都有肯定答案，而在于已知、推论与未知之间的边界清楚，并且留下可以继续核查的路径。

\chapter{中国经济的制度结构}
\section{从部门关系代替抽象标签}
中国大陆经济包含市场交易、公共预算、国有资产、私人产权、家庭劳动和国际交换。把它命名为某种经济类型，可以指出历史定位，却不足以解释一家企业为什么容易融资、一名劳动者为什么难以获得当地服务，或者一个地方为什么依赖土地收入。可操作的分析从部门及其相互要求开始。

% ARGUMENT P018
\phantomsection\label{arg:P018}
制度结构的分析可以从一笔普通投资开始。企业决定购买设备，需要预期订单、获得融资、雇用人员并取得相应经营条件。设备供应商希望及时收款，银行希望未来现金流覆盖贷款，地方可能希望项目带来就业和税收，家庭则关心岗位是否稳定。这些诉求既有共同部分，也存在不同时间尺度：企业需要当期现金，地方可能重视未来税源，劳动者需要每月收入。

市场与政府因此并非两块彼此隔绝的空间。合同的执行需要法律和公共组织，公共项目也通过采购与企业发生关系。将某项结果归因于“市场”或者“政府”，必须进一步说明是哪一种规则、哪一个组织、哪类激励产生了作用。否则，同样的失败可以被随意归到任何一方，理论也就失去检验对象。本书采用部门关系，是为了让这些笼统标签落到可以观察的交易与责任上。

家庭提供劳动，进行消费与储蓄，也可能拥有住房、企业和金融资产。企业组织生产，支付工资和税费，向银行借款或发行证券。政府制定和执行规则，提供公共服务，征收税费，管理公共资产，并承担部分最终风险。银行与其他金融机构连接不同期限的资金要求，同时在自身资产负债表上承担信用和流动性风险。对外部门通过贸易、投资和资产交易与国内联系。部门划分是核算工具，不能把同一人的所有活动固定在一个身份里。

% ARGUMENT P019
\phantomsection\label{arg:P019}
核算部门与社会阶层之间并非一一对应。家庭部门中既有主要依靠工资生活的人，也有大型企业的最终所有者；企业部门中既有经营自主权较强的集团，也有依赖单一客户的小供应商。把家庭部门收入增加理解为普通劳动者收入增加，或者把企业部门亏损理解为所有资本持有者同等受损，都会越过资料能够支持的范围。

因此，本书在总量账户之下再加入关系位置。研究谁取得收益时，要追踪企业利润最终成为谁的工资、分红、留存或税款；研究谁承担风险时，要追踪企业损失是否由股东资本、债权减值、欠款或公共支持吸收。某个人可以处在不同账户的交点上，但他的实际决策权和可用资产并不因此自动相等。部门分析提供总量约束，阶层分析解释约束内部的分配，二者相互补充。

一个经营小餐馆的家庭，同时是劳动者、经营者、租户和借款人。一名国企员工既领取工资，也可能有住房和基金份额，但不能因此认为他实际控制了企业重大投资。一家民营企业可能依赖公共基础设施与政策性融资，一家国企也要面对产品市场和债权人要求。现实关系在多条轴线上交叉，所有制是其中重要的一条，而不是唯一的一条。

% ARGUMENT P020
\phantomsection\label{arg:P020}
以小餐馆为例，营业收入需要先支付食材、房租、水电、工资、税费及设备维护，余额中还包含经营者本人的劳动报酬。若直接把余额全部视为资本利润，就会把厨师、采购和管理劳动混入资本收益；若因为老板亲自劳动就忽略其雇佣关系，又会看不到他对员工工时和工资的决定权。这是一种兼有劳动依赖与局部控制的位置。

企业之间的权力关系也会改变这种位置。餐馆可能决定几名员工的工作，却难以改变商场租约、平台规则或银行条件；大型连锁企业则可能拥有更强议价能力。分析不能只问“是否有营业执照”，而应问经营者收入中有多少依赖本人持续劳动，能否替换关键交易对象，经营失败时是否会牵连家庭住房。这些问题解释了为什么名义上同属经营者的人，对经济变化会有截然不同的承受能力。

\Chart{sector}{2025年增加值的产业构成}{数据图；单位：GDP占比，\%。来源：国家统计局2025年统计公报。第三产业包含市场服务与公共服务，不能直接等同于“金融化”。\cite{NBS2025}}

\section{国家不是一个只有一种目标的行动者}
中央政府、地方政府、预算部门、国有企业和金融机构具有不同职责、考核和约束。中央可能关注全国稳定、跨地区均衡和长期能力，地方则同时面对本地就业、财政支出和招商竞争。部门之间既能合作，也可能转移成本。分析这些差异，并不等于否认统一制度，而是说明统一制度需要通过组织链条发挥作用。

% ARGUMENT P021
\phantomsection\label{arg:P021}
多层次政府关系的困难，来自信息与责任分布的不一致。基层更容易知道本地道路、学校和企业的具体情况，上级更容易协调跨地区资源和控制系统性风险。如果所有决定都集中，信息可能在上报中丢失；如果所有决定都分散，地方可能忽略对邻近地区和未来财政的影响。这里需要比较的是不同事项适合怎样配置决策权，而不是预先肯定集中或分散本身。

财政关系尤其能显示这种张力。流入地承担新增公共服务成本，部分劳动者家庭成员却仍留在流出地；项目的税收收益可能集中在某地，生态影响却跨越边界。若责任与财力不能对应，地方就可能减少服务或寻找非稳定收入。公共资源配置者、普通执行人员和服务使用者因此处于不同位置。前者选择预算，后者可能只能执行既定额度，居民则承受服务不足。把三者合并为一个群体，不能解释责任怎样传递。

周黎安对地方官员晋升激励的研究，把增长表现与组织内的激励联系起来。它为地方为何积极发展经济提供了一个机制，也指出目标和政府职能之间可能发生偏离。这个2007年的研究不是2026年考核制度的实时测量，不能用它断言所有地方仍只追求GDP。更合理的使用方式，是把考核目标、财政回报、任期和问责条件作为需要重新观察的变量。\cite{Zhou2007}

% ARGUMENT P022
\phantomsection\label{arg:P022}
晋升激励的分析要经过若干中间环节：上级选择什么指标，地方能否影响这些指标，指标如何进入考核，干部是否预期考核会影响职业结果，最后才是项目选择。任何一个环节改变，行为都可能改变。例如，指标若从建设规模转向长期使用效果，能够迅速显现但后续收益不足的项目就未必仍有优势。若公共服务质量缺乏可靠信息，形式上的多目标考核也可能继续偏向容易计量的指标。

这一机制还需要与其他解释比较。地方扩张投资可能来自居民实际需求、产业协作机会、上级专项任务或金融条件改善，而非晋升竞争单独造成。研究应观察相似地区在激励变化前后的不同表现，并检查资金来源和需求条件是否同时变化。历史论文提供了可用的问题结构，但不能代替当前考核材料和项目证据。本文因此采用其分析思路，不把所有地方行为还原为同一种动机。

熊伟的“官员治理增长模型”将中央与地方之间的委托代理问题纳入增长框架，提示增长激励可能伴随短期行为和杠杆积累。本文采用的是这种条件化思路：当收益在任期内显现、成本在未来兑现，而且损失可转由更大范围承担时，组织更容易偏好扩张。反过来，如果信息公开、长期责任与公共服务评价得到加强，这种倾向可能减弱。不能把模型中的假设直接当成每一个地方的事实。\cite{Xiong2018}

% ARGUMENT P023
\phantomsection\label{arg:P023}
短期收益与长期成本分离，会产生一种跨期责任问题。一个建设项目在开工时形成投资和就业，在运行后才显现维护费用与需求不足。如果决策者只在前一阶段接受评价，未来成本即使可以预见，也可能没有充分进入当期选择。这一问题同样存在于私人企业：管理者可以获得短期奖金，而项目风险在其离任以后才暴露。

但是，延长责任期限并不总是越长越好。若对任何未来损失都追究最初决策者，管理者可能不愿承担具有合理不确定性的创新项目。更合适的区分是，决策时是否披露了风险、是否遵循程序、是否隐瞒信息，以及后续损失究竟来自不可预见冲击还是明知不可持续的安排。制度应约束可避免的风险转移，同时保留合理试错空间。否则，纠正冒险可能转化为普遍不作为，生产能力也会受到损害。

\section{公共所有权与公众得到什么}
法律上的国家所有，是财产归属的规定；管理者拥有的配置权限、企业保留的利润、公共财政收到的收入和居民实际享受的服务，是不同层次的关系。现行公司法对国家出资公司设专章，并规定履行出资人职责以及党组织参与重大经营管理事项的制度安排。研究必须把这些正式权利与具体经营结果联系，而不是用私人股权结构机械套用。\cite{CompanyLaw}

% ARGUMENT P024
\phantomsection\label{arg:P024}
公共所有权包含一个较长的委托链条。公众通过制度拥有公共资产，履行出资人职责的机构行使股东权利，企业治理结构决定经营事项，具体管理者安排生产与用工。每一层都可能掌握上层难以完全观察的信息，也可能面对不同考核。法律明确归属和权限，是分析的起点；收益怎样进入预算和服务，需要沿着这条链继续追踪。

公共资源配置者因此不能简单等同于私人所有者。配置权通常受到法定职责、组织程序与任期限制，不必然形成可出售、可继承的私人权利；但它仍能影响进入机会、融资和公共服务，对社会位置具有实质意义。普通公共部门雇员通常不掌握这类权限，他们的稳定工资也不等于控制公共资本。区分这些层次，既可以分析公共权力的经济作用，也能避免根据职业单位就给所有人赋予相同权力。

公共资产可能通过利润上缴、税收、基础设施服务、公共目标和风险承接影响居民。它的效果不能只看分红，也不能把所有可能的社会贡献都想象成已经兑现。若某项政策要求企业提供低价服务，需要说明受益人是谁、成本如何核算、财政是否补偿、质量是否下降。若企业承担战略储备，应把保险价值与日常经营回报分开评价。

% ARGUMENT P025
\phantomsection\label{arg:P025}
低价公共服务的分配效果，要看服务对象和使用强度。普遍降低某项费用，可能使使用更多的家庭取得更大绝对收益；针对基本需求设置优惠，则可能更集中地帮助保障不足者。两种安排都有管理成本和识别困难。评价不能仅看价格下降，而应观察谁实际使用、质量是否稳定，以及费用减少是否由其他税费或未来债务承担。

战略储备则需要采用不同尺度。某些设施在正常时期利用率较低，却在供应中断时避免严重损失，其作用类似保险。合理评价应比较建设维护成本、风险发生概率和替代方案，而不只比较当期营业利润。但保险价值也不能无限推定。若长期没有明确风险对象、储备无法使用或重复建设过多，就应重新评估。公共任务需要被具体说明，才能同时防止对必要服务的低估和对低效项目的过度保护。

本文提出“公众可兑现收益”的分析对象：公共资产在扣除真实维护、必要再投资和可识别风险成本后，经过财政和服务制度转化为居民实际利益的部分。这不是一个已公布的全国统计量，也不是把公共资产机械平均到个人名下。它要求逐个记录传递环节，从资产产生收益，到预算分配，再到服务的实际可获得性。

% ARGUMENT P026
\phantomsection\label{arg:P026}
“可兑现”至少包含资格、可达性和持续性三个条件。一个家庭在名义上享有某项服务，如果附近没有供给、办理成本过高或者等待时间过长，收益就尚未充分实现；若服务只能维持很短时间，家庭也难以据此安排就业和照护。收益不只是预算拨付，更包括资源在使用者一端形成的实际效果。

这一概念不能直接计算成“公共利润减去风险等于居民福利”的单一公式。维护与投资有跨期作用，服务质量不总能货币化，风险成本还存在不确定性。实际研究可以先采用并列指标：可持续上缴、服务覆盖与质量、居民费用、资本维护和或有负债。若某项改革提高分红却明显降低维护，便不能只凭现金流增加宣布成功。本文提出的是评价对象与追踪路径，尚未提供可以比较全国所有公共企业的综合得分。

这一视角也适用于私人产权。拥有股份不保证当期分红，购买理财不意味着直接决定贷款，拥有住房不意味着有足够现金支付医疗费用。产权形式规定资格，收益兑现还受管理、合同、市场和时间影响。公共所有与私人所有都需要研究中间机制，只是两者的问责对象和制度目标不同。

% ARGUMENT P027
\phantomsection\label{arg:P027}
私人产权的兑现问题可以用股权解释。一份股票赋予一定剩余索取权，但股东能否影响决策取决于持股结构、投票与治理；企业盈利是否变成个人现金，还取决于分红、再投资和市场交易。小额基金持有人通常更依赖中介治理，不能像控制股东那样选择企业战略。因此，“拥有资本资产”与“作为资本控制者行动”需要分开。

同一原则也适用于住房。一套自住房提供居住服务和一定风险缓冲，却可能没有稳定现金收益；多套出租物业的持有者则具有不同收入来源与调整能力。如果两者都被简单归为有产者，阶层分析就会失去对劳动依赖的辨别。研究实际收益与退出能力，并不是否认产权重要，而是说明产权必须通过具体使用、合同和市场条件，才转化为生活机会。

\section{三种预算约束及其混合}
企业、地方政府和家庭面对的约束并不相同。家庭通常难以长期维持现金流赤字，失业和疾病可迅速迫使其调整。企业可能依靠融资和股东投入跨越亏损阶段，但能否继续融资取决于信用、抵押、前景及救助预期。政府可以通过征税、债务和跨期安排进行调整，却也受真实资源、制度权限和公众承受能力约束。

% ARGUMENT P028
\phantomsection\label{arg:P028}
家庭、企业和政府的跨期能力不同，还与收入来源的稳定性和借款权有关。家庭通常不能以未来全部劳动收入自由抵押，企业可以发行股权但需要让渡收益与控制，政府则需要在法定权限和财政制度内安排收入。信用较强的主体能够等待项目成熟，信用较弱者可能在尚有长期价值时就被迫出售资产或退出。

这种差异既有合理部分，也可能强化既有优势。贷款者需要防范无法偿还的风险，所以要求担保并非天然不合理；但如果只有已经拥有资产的人能够获得创业融资，新进入者的能力就可能难以转化为项目。公共政策需要比较信息改善、担保、风险共担和破产安排的效果，而不能只要求银行对所有人提供相同贷款。平等机会与完全相同的融资条件不是一回事，前者需要降低与真实风险无关的障碍。

软预算约束研究强调的不是“有亏损”，而是组织预期在财务困难时会获得追加支持，这种预期会反过来改变事前行为。它既可能出现在国企，也可能出现在具有系统重要性的私人金融机构或大型企业。若把软预算约束简单翻译为某一种所有制的专属缺陷，就会丢失最重要的动态激励。\cite{Kornai2003}

% ARGUMENT P029
\phantomsection\label{arg:P029}
软预算约束的关键证据应当出现在救助之前。研究者需要观察组织是否在融资条件、投资选择或谈判中体现了对支持的预期。事后得到救助，可能因为突发冲击危及必要公共功能，并不说明此前已经存在过度冒险；事后没有救助，也不证明此前不存在错误预期。预期与实际结果应分别记录。

风险可能通过多条渠道被转移：债权人续贷，供应商被迫延长账期，员工推迟取得报酬，财政承担项目支出，或者服务使用者支付更高费用。某些转移在法律账户中十分明显，另一些则表现为时间和服务质量的变化。若只看债务是否违约，就可能把后几种成本遗漏。对不同阶层而言，获得支持的机会也不均等，小经营者和缺少抵押的家庭通常难以采用与大型组织相同的跨期策略；具体差异仍需合同和案例验证。

救助也不天然无效。在危机中维持支付、关键公共服务和有前景的生产能力，可以减少连锁损失。问题在于救助是否保护了必要功能，是否要求原有股东和管理者承担责任，是否形成新的无条件扩张预期。企业被保留与原有所有者不受损失，是两件不同的事。严格识别这一点，才可能同时讨论稳定和纪律。

% ARGUMENT P030
\phantomsection\label{arg:P030}
保护功能与保护原有权利的区别，可以用一项必要服务说明。若服务企业陷入财务困难，公共部门可以支持其继续运营，同时调整股权、债权期限或管理安排。用户获得服务连续性，并不要求原有投资者维持全部收益。相反，若支持只补足所有既有索取权，却没有改变低效经营，公共资金可能只是延迟下一次损失。

实施中仍有困难。迅速让债权人承担损失，可能影响其他机构的资产和融资；拖延确认又会掩盖真实成本。选择需要比较不同方案的总损失和分配，而不是只追求某一账户表面稳定。对普通家庭而言，维持支付和基本服务通常比维护特定投资产品的原有回报更直接；对供应商和劳动者而言，明确清偿顺序与信息公开也可能比含糊的“整体化解”更重要。制度设计必须把这些对象分别说明。

\section{一个部门的盈余必须放进整体}
在简化开放经济核算中，可以用支出恒等式说明部门联系。这里的消费、投资和政府购买均是最终支出，不包括同一产品在生产链中的重复销售：
\[
Y=C+I+G+X-M.
\]
其中$C$为居民最终消费，$I$为全社会资本形成（含存货变动），$G$为政府最终消费；政府投资计入$I$，不再计入$G$。这是一项事后核算关系，不说明增长一定由哪项支出造成。增加政府购买可能带动闲置产能，也可能在供给受限时推高价格；出口增长可能带动工资，也可能集中在就业扩张有限的行业。要解释作用，需要额外行为假设和现实证据。

% ARGUMENT P031
\phantomsection\label{arg:P031}
支出恒等式中的每一项都有对应产出，但增加支出能否增加实际产量，需要生产端有可用资源。若存在闲置设备、可匹配劳动和原料，新增需求可能带动产出；若关键环节已经受限，更多名义支出可能首先提高价格。供给约束还可能只存在于某些行业，所以总量刺激与结构调整需要结合。

阶层分配会改变需求的组成。增加低收入家庭可用收入，可能较快转化为基本消费；增加高资产家庭财富，可能更多进入储蓄或资产交易。这里是需要微观资料检验的行为差异，不能由恒等式直接证明。即使消费倾向不同，实际生产效果仍取决于国内供给、进口和价格反应。因此，政策评价要同时追踪钱由谁获得、花在什么地方、谁提供相应商品，以及新增收入是否再次集中，而不只是观察某项支出的总额。

同样，全国储蓄与投资的差额对应经常账户余额，但一个家庭节省开支不能保证所有家庭的收入同时增加。若许多家庭因风险同时削减消费，而企业又因需求不足减少投资，总收入可能下降。这不是批评节俭的道德价值，而是指出个人策略汇总以后会改变其原先赖以成立的条件。

% ARGUMENT P032
\phantomsection\label{arg:P032}
储蓄与投资关系还需要区分国内生产总值和国民可支配收入。若把储蓄定义为国民可支配收入减去最终消费，储蓄投资差额才与经常账户相对应；在只包含国内生产、消费和货物服务贸易的简化模型中，得到的则是净出口关系。跨境要素收入和经常转移不能在精确核算中省略。本文在使用简化式时说明边界，避免把不同体系的余额拼成一张表。

个人节俭形成总量收缩，是一种有条件的反馈。如果减少消费的资金被有效转化为有用投资，或者其他部门扩大需求，收入未必下降；若储蓄愿望增加但投资需求不足，收入调整就可能使原先计划的储蓄无法实现。因而，“增加消费”和“提高储蓄”都不能脱离环境独立判断。应观察生产能力、未来需求和金融传导，而不是用对某种个人习惯的评价代替宏观机制。

讨论地方债务时也要考虑合并边界。将一家融资平台的债务转为地方政府债券，会改变法律责任、利率和期限，但在足够宽的公共部门合并账户中，不一定减少相同数额的最终负担。债权人同时持有资产，不能把债务本金与资产损失、后续财政支出不加区分地累加为“总成本”。

% ARGUMENT P033
\phantomsection\label{arg:P033}
合并账户首先要确定哪些机构属于同一经济责任范围。法律上独立的企业是否存在可信担保，担保覆盖本金还是特定项目，公共支持是否有明确上限，都会改变风险判断。不能因为一家企业名称中带有地方或国有，就把全部债务自动视为财政债务；也不能只因法律形式独立，就忽略反复支持形成的实际预期。

债务置换主要改变偿付时间、利率和责任安排。较低利率可以减少未来现金支出，较长期限可以缓解当期压力，这些都具有实际价值；但如果项目缺少持续收入，延期并不自动填补长期缺口。真正的损失需要结合项目资产价值、未来收入和各方让步判断。把旧债本金、新债本金和未来付息全部相加，会严重夸大成本；只报告利率降低而忽略新增担保，则可能低估成本。两种错误都妨碍公开分担。

\section{本章的修订与可检验问题}
对经典国家与资本关系的修订，是从一个总括性的主体关系，推进到可观察的组织与账户关系。中国制度的重要特征，在于公共权力、公共产权、地方激励、金融体系与市场竞争相互嵌入。这个特征本身既不是效率证明，也不是低效判决。

% ARGUMENT P034
\phantomsection\label{arg:P034}
从组织关系出发，并不意味着国家只是彼此分离的部门之和。全国性规则、统一市场、中央财政能力和公共目标，可以改变各部门的行动边界。问题在于这些总体能力怎样通过具体制度落实。例如统一的服务资格，需要地方供给、结算和人员配置相配合，才能成为居民实际获得的保障。

反过来，部门行为也会反馈到整体制度。地方为争取项目设置的优惠，可能影响全国税基和产业布局；大型企业形成的融资需求，可能改变银行风险集中度。整体与局部之间存在双向作用。本文因此既避免把所有结果归于一个全知全能的主体，也避免把系统性结果解释为完全偶然的个体选择。制度结构提供稳定的激励与约束，具体组织则在其中根据自身信息和利益行动，二者共同形成可观察的经济过程。

可检验的问题是：在控制行业、地区和风险之后，哪些权利安排改善生产与公共服务，哪些安排主要改变收益归属和损失转移？公共资产增加是否转化为可观察的服务改善，还是更多留在企业扩张与债务承担中？地方获得新的融资空间后，新增支出流向公共服务还是再次形成低回报项目？这些问题需要项目和部门层面的证据，不能由全国总量直接回答。

% ARGUMENT P035
\phantomsection\label{arg:P035}
这些检验不能只依靠“控制了若干变量”的回归表。行业、地区和风险本身可能受到制度安排影响，过度控制有时会把政策真正通过的渠道消除；遗漏重要差异又会把原有优势误当作改革效果。因此，应先画清因果顺序，再决定比较什么。最基本的要求，是知道政策发生前各组是否已经呈现不同趋势。

公共资产改革还可能有外溢影响。某地改善服务，吸引更多人口流入，居民平均获得变化就可能同时反映服务改善与人口构成变化。若只比较改革后的平均数，会混淆两种效果。研究可以同时追踪原有使用者、新进入者和供给总量，并观察财政负担如何调整。本书没有执行这样的新研究，而是明确提出数据需求。对识别条件的说明，能够使未来检验真正针对命题，而不是事后选择有利数字。

本书据此把制度看成一套分配决策权、现金流权和风险责任的规则。它不取代历史分析，而是使历史分析能够落在具体对象上。此后的各章，都将追问这三种权利是否一致，以及不一致时发生什么。

% ARGUMENT P036
\phantomsection\label{arg:P036}
决策权、现金流权与风险责任的一致，也不是要求三者永远由同一个人承担。现代组织本来就依赖分工：专业管理者决策，投资者提供资金，保险和公共部门分担部分风险。这种分离可以提高效率。需要关注的是分离以后是否存在信息披露、约束和补偿，使决策者不能持续把可避免损失推给缺少选择的人。

本章由此留下一个贯穿性问题：当一个主体拥有改变规则的能力，却不承担相应后果时，谁能纠正这种安排；当一个主体承担重要后果，却缺少参与决定的渠道时，制度如何表达其利益。劳动合同、地方财政、平台治理和家庭照护都会以不同形式遇到这个问题。各章将把它落实到具体合同和资源条件，逐步建立生产、分配与风险之间能够相互校验的解释。


\chapter{商品、价值与价格}
\section{先问一种东西改善了什么生活}
一杯饮用水、一项医疗服务和一个搜索结果，都可能满足需要，却以不同方式形成和交换。水依赖自然条件、处理设施和输送网络；医疗依赖专业劳动、设备以及患者难以完全判断的质量；搜索结果可以低成本复制，却依赖此前投入的算力、数据组织和入口。把它们称为商品，说明它们进入交换关系，并没有回答每项收入究竟来自劳动、稀缺条件还是排他权利。

% ARGUMENT P037
\phantomsection\label{arg:P037}
商品的两面之所以需要分开，是因为满足需要与取得付款不是同一件事。某些急迫需要缺少支付能力，市场就可能无法形成足够需求；某些购买只是为了未来转售，成交价格却很高。这里不能简单推断市场价格无用，而是要说明它反映的是受收入分配和制度约束的支付意愿，并不直接等于需要的强度。

生产本身也有不同结构。可储存的标准产品可以在不同时点销售，医疗和照护常要求生产与使用同时发生，数字产品则可能具有较高前期成本和较低复制成本。不同结构改变库存、定价和质量监督。若用制造标准商品的模型解释所有服务，便容易把低复制成本误解为零生产成本，或者把服务难以加速误解为劳动者效率低。价值研究需要从实际活动开始，再决定哪些抽象能够用于比较。

经典分析区分使用价值与交换价值，并以社会必要劳动时间解释商品价值。这个区分仍有用：价格高不保证有用，价格低也不保证社会贡献小。疾病救治、基础教育和照护的价值，不能仅由购买者当期支付能力表示。但从这一判断走向“全部价格最终都由劳动时间唯一决定”，还需要额外证明。\cite{Marx1}

% ARGUMENT P038
\phantomsection\label{arg:P038}
社会必要劳动时间并非任何劳动者实际花费的时间。若技术落后的生产者需要更多工时，其全部耗费不因此都被社会承认为同等价值；产品还需要满足一定社会需求，私人劳动才能在交换中得到确认。因此，经典概念本身已经包含技术条件、平均效率和社会承认，不能简化为“花的时间越长就越值钱”。\cite{Marx1}

但加入“社会必要”也带来检验难题。如果必要性完全由最后成交价格倒推，理论就容易变成用结果定义原因。为了使解释具有经验内容，研究者应尽可能在价格之外测量技术、工时、质量与需求条件，再观察其关系。这里要批评的是不可检验的循环，而不是否认社会劳动具有客观约束。劳动价值论能否解释某一对象，应通过其明确版本讨论，不能把所有版本混为一个命题。

本文把社会财富理解为在一定分配和生态条件下，可以持续满足人的需要的物品、服务与能力。它不宜压缩成一个没有争议的总分。桥梁连接了哪里，教育改善了谁的能力，电力供应是否可靠，污染是否损害健康，都是财富评价的一部分。市场价格提供稀缺性和支付意愿的信息，公共评价则补充市场难以完整表达的需要与外部影响。两者都可能失真，应相互校正。

% ARGUMENT P039
\phantomsection\label{arg:P039}
财富还包含维持能力的条件。一所学校的建筑具有物质形态，教师的知识和学生获得的能力却不能全部表现为可交易资产；一条河流的清洁状态可能支持农业和健康，却不必然进入任何人的私人财产表。将财富只理解为市场资产，会低估这些条件；将所有社会需要都计入一个任意估价的总财富，又会造成不可复核的扩张。

本文因此采用层次化描述：先记录可以直接计量的物品与服务，再记录质量、可获得性和持续条件，最后讨论分配。相同总量的医疗资源，若只集中在少数人能够进入的机构，其社会效果与广泛可用时不同。这里的差异并非纯粹生产数量差异，而是生产与权利安排共同作用的结果。对各阶层而言，财富的意义在于能否实际使用、是否需要承担过高代价，以及这种使用能否在风险发生时继续维持。

\section{劳动投入解释什么，又不能解释什么}
劳动价值论的一个强版本认为，在可重复生产的商品之间，交换比例应当围绕直接和间接劳动投入形成。所谓间接劳动，是制造机器、原料及其上游投入所需的劳动，不能只数最终车间的工时。如果忽略产业链，一台几乎无人操作的机器就会被误认为完全脱离劳动。

% ARGUMENT P040
\phantomsection\label{arg:P040}
完全劳动投入可通过产业间依赖理解。制作一件最终产品需要直接劳动，也需要材料、能源和设备，而这些投入在上游又消耗劳动。只把所有企业工时直接相加，会把中间产品与最终产品混在一起；正确方法要沿着投入关系递推，在一致的生产边界内计算为最终需求所需的劳动。若存在进口，还要决定采用国内劳动还是包括国外劳动，二者回答不同问题。

即使技术系数已知，劳动也未必完全同质。训练、经验与工作强度如何折算，会影响结果。可以分别保留不同劳动类型，也可以采用公开的标准化方案，但不能把争议藏在一个总工时数字里。对中国制造业的检验尤其需要处理服务投入、外包和进口零部件。最终工厂的自动化，并不说明产业链上的劳动消失；然而发现上游劳动存在，也不足以证明它唯一决定了最终价格。

但严格比例关系遇到几个困难。同样劳动时间可以产出不同质量；消费者支付能力不平等；自然资源和区位不能自由复制；不同产业需要垫付的设备和周转资金不同；风险、垄断和税收影响成交价格。即使把技能折算为简单劳动，也需要说明折算依据。如果用工资或价格反推技能权重，再用加权劳动解释价格，论证会发生循环。

% ARGUMENT P041
\phantomsection\label{arg:P041}
利润率与资本周转是另一项必要条件。如果两个部门每年使用相同劳动，却需要垫付不同规模、不同期限的资金，在资本可流动并要求回报的条件下，它们的长期供给价格就未必按劳动量同比例形成。这是生产价格问题的核心之一。马克思第三卷已经讨论竞争如何使不同部门的剩余分配发生变化，因此现实价格偏离劳动量本身，不能单独构成对整个理论体系的否定。\cite{MarxPrice}

真正的争议在于怎样一致处理投入价格、工资、利润和总量关系。不同理论方案采用同时估价、历史成本或货币表达等不同方法，所回答的问题也不同。本书不把其中任何一种写成已经结束争论的结论，而是要求提出可复算账户。若一个方案能在定义上维持某个总量等式，还需要进一步说明它是否提供了对观察结果的独立解释；定义上的一致与经验上的有效是两项不同要求。

一种较弱的命题是，劳动时间约束一个社会在既定技术下能够生产多少东西。这比价格比例命题更容易成立，却不意味着劳动是唯一物质来源。自然提供条件，设备和知识提高生产能力，组织协调减少浪费。马克思在《哥达纲领批判》中也明确讨论自然对使用价值的作用。因此，修订劳动价值理论不必从“劳动之外没有任何贡献”这样过强的前提出发。\cite{Gotha}

% ARGUMENT P042
\phantomsection\label{arg:P042}
劳动并非财富的唯一物质条件，并不意味着劳动关系不重要。自然资源、设备和知识在生产中共同作用，但谁能够使用它们、谁有权决定用途、谁可以要求分配，仍然是社会制度问题。把设备具有生产作用直接转化为设备所有者理应取得某一确定份额，是从技术事实跳到了规范结论；反过来，把劳动的社会重要性直接转化为全部收入只能按工时分配，也忽略了维护、风险与公共需要。

有用的分析应分三步：先辨认生产所需条件，再解释这些条件怎样被组织，最后讨论收益分配规则。技术层面可以存在互补，分配层面仍可能存在冲突。一台机器使工人产出增加，工人与机器的技术关系是互补的；机器投资者与工人如何分享增加部分，则取决于合同、竞争与议价。把两种关系分开，才能既承认技术合作，又分析社会权力。

这里应保留的是对社会劳动配置的追问：谁决定劳动用于哪些需要，多少劳动被重复竞争和低质量生产耗散，谁有能力把别人的劳动转化为自己的收入要求。应放弃的，是未经实证检验就把劳动时间与每项实际价格直接画等号。弱化一个不能普遍成立的价格命题，可以使关于劳动组织的研究更精确。

% ARGUMENT P043
\phantomsection\label{arg:P043}
劳动配置问题还包含无效劳动的识别。重复建设、返工、错误信息和不必要的等待，会占用时间却未必增加最终可用产出。但“无效”不能仅由研究者对产品的个人喜好决定，也不能因为企业失败就倒推当初所有试验都没有社会价值。创新需要承担部分失败，关键是决策是否基于合理信息、失败能否产生可共享知识，以及成本是否由有能力承担的人负责。

在中国的具体研究中，可以比较相同服务的完成时间、质量和资源投入，观察组织改进是否真正减少浪费。若减少企业内部工时却增加消费者自行办理或家庭照护时间，总社会耗费未必下降。劳动价值问题由此扩展为跨组织的时间配置问题。但这个扩展需要保持账户边界，不把所有生活时间都当作应被压缩的成本。休息、交往和照护也可能直接构成生活的目的。

\section{一个可以算清楚的价格例子}
考虑一个完全假设的两部门经济。生产一单位甲产品需要0.2单位甲投入和0.5小时劳动；生产一单位乙产品需要0.4单位甲投入和0.25小时劳动。这里忽略固定资产折旧、联合生产、税收和资源稀缺。完全劳动投入满足
\[
\lambda_{\text{甲}}=0.2\lambda_{\text{甲}}+0.5=0.625,
\qquad
\lambda_{\text{乙}}=0.4\lambda_{\text{甲}}+0.25=0.5.
\]
甲乙劳动投入之比为1.25。现在另作一个定价假设：两部门工资每小时为1，原料成本取得20\%的加成，工资在销售时支付、不计入加成基数。价格方程变为
\[
p_{\text{甲}}=1.2(0.2p_{\text{甲}})+0.5,
\qquad
p_{\text{乙}}=1.2(0.4p_{\text{甲}})+0.25.
\]
求得甲价格约0.658，乙价格约0.566，比值约1.163，与劳动投入比例不同。全部数字都是示例，不是中国投入产出表的估计。改变工资支付时间和利润率条件，会再次改变结果。

% ARGUMENT P044
\phantomsection\label{arg:P044}
这个算例的求解顺序可以进一步说明。甲部门需要的甲投入，也必须由甲部门自身生产，所以一单位最终甲产品要求的劳动不是直接工时，而是直接劳动与上游循环投入劳动之和。把这一递推写成方程，才得到完全劳动系数。乙部门则使用甲产品，因此其劳动要求包含甲部门形成的间接劳动。两个方程展示的是给定技术下的物量关系。

价格方程另外加入了一个分配假设：垫付原料要求回报，而工资在期末支付。若工资也预付，或者两个部门资金周转速度不同，价格体系会改变。这里称为加成，是为了明确模型的设定；它不证明现实企业总能取得这一比例，也不意味着已经测定了均衡利润率。读者可以把加成设为零，此时价格与劳动系数在这组特殊假设下重新成比例。这个变化说明，分配与时间条件正是连接投入和价格时不能省略的环节。

这个小例子不能证明现实价格遵循统一利润率，也不能证明劳动价值毫无解释力。它仅说明：即使没有垄断、欺骗和资源租金，生产结构与成本回收条件也足以使价格比例偏离劳动比例。要解释现实，必须提出可观察的生产关系和定价条件，不能用一个恒等定义替代这些条件。

% ARGUMENT P045
\phantomsection\label{arg:P045}
算例的用途，是排除一个过强推论：既然每件商品都涉及劳动，就能仅凭劳动数量确定所有交换比例。它没有排除劳动投入对成本和长期供给的约束，也没有证明市场价格总能准确表达社会资源耗费。两种极端结论都超出了模型。一个反例可以说明某项普遍断言不成立，却不能自动建立另一套普遍理论。

把模型放回中国经济，需要加入实际生产组织。不同企业融资成本、产能利用、进口投入、税费和品牌条件各不相同，观察价格时还要区分出厂、批发和零售环节。如果研究者只比较最终售价与工厂工时，就可能把运输、销售、售后与金融安排全部误记为不明剩余。解释这些环节并不意味着认可所有收费，而是先把收入归属定位清楚，再判断哪些来自必要服务、哪些来自市场力量。

经典的价值向生产价格转化问题，因此应被视为模型一致性问题。对投入与产出采取什么估价、是否比较同一时期、总利润与总剩余如何对应，都会影响结论。本文不宣称通过重新命名解决了长期争论，而采用双重记录：劳动投入说明社会劳动要求，价格与收入账户说明货币分配。二者联系需要估计，不能预先强制一致。

% ARGUMENT P046
\phantomsection\label{arg:P046}
一致核算要求投入与产出的时点相容。历史成本账户记录企业实际支付了多少，重置成本账户则询问今天重新取得同样生产条件需要多少。物价变动时，两种利润会不同，但都可能有用：前者接近企业现金与财务经历，后者更接近持续再生产的条件。把历史成本投入与当前重置价值产出随意混合，会制造并不存在的收益或损失。

劳动与货币的双重记录因此是一种研究纪律。它允许观察单位最终需求需要多少劳动，以及这些劳动通过货币交易获得了什么收入；同时保留资产重估、地租和金融收益的独立记录。某些理论方案试图通过特定转换关系使总量一致，本文可以讨论这些方案，但不以预先设定总量相等代替检验。对于无法取得完整投入产出与工时资料的部分，最诚实的结论是尚不能判断量化关系，而不是以概念推演填补数据。

\section{市场价格中的四类机制}
理解中国价格，可以先区分生产成本、供需与质量、市场力量和制度性权利四类机制。制造商提高良品率，可能在降低单位劳动投入的同时扩大销量；医院面对信息不对称，价格未必充分表达疗效；核心区位的土地使用权具有难以复制的条件；数字平台可以通过入口、排序和转移成本影响收费。它们可以同时作用，不能把观察到的全部利润都放进同一个解释。

% ARGUMENT P047
\phantomsection\label{arg:P047}
这四类机制之间可能互相作用。例如，质量提高可以增加消费者愿意支付的价格，同时需要更高技术投入；专利既可能保护研发回报，也可能构成暂时的进入限制；规模扩大可能降低成本，又增强企业对供应商的议价能力。现实收入通常由多种机制共同形成，不能仅凭企业所属行业就一次性分解。

辨认机制应寻找条件变化。若原料成本下降而售价长期不降，需要考察需求、质量与竞争是否变化；若新的竞争者进入以后价格下降而质量稳定，市场力量解释就获得更多支持。即使如此，价格下降也可能来自新技术而非竞争本身，仍需比较成本。阶层分析则继续追踪变化的后果：消费者节省费用是否以供应商利润或劳动条件恶化为代价，或者生产效率确实使多方同时受益。

判断一种收益是否主要来自限制竞争，至少需要考察进入条件、替代选择、成本变化与收费变化。如果一家企业先承担研发失败风险，成功后获得暂时超额收益，这可能是创新回报；如果它进一步通过排他规则阻止没有效率劣势的竞争者进入，收益中可能增加租金成分。创新与租金不是企业永久不变的身份，而是不同阶段、不同业务中的机制。

% ARGUMENT P048
\phantomsection\label{arg:P048}
经济租金是相对于反事实的概念，即超过维持某项资源供给所必需回报的收益。这个必要回报通常不能直接观察，因为研发失败概率、机会成本和进入条件都会影响判断。一家成功企业的高利润，不应脱离失败项目和投资风险单独评价；但不能因为创新有风险，就认为成功以后任何排他收费都合理。

应把创新激励与持续限制竞争的权利分开讨论。某段保护期可能支持前期投入，保护过宽或无法退出则可能阻碍后续改进。对于资本所有者，临时超额收益是积累来源；对后来进入的中小企业，过高门槛可能限制竞争；对劳动者，企业优势有时提高工资，也可能主要增加股东收入。政策的对象应是具体权利与行为，而不是永久给某一企业贴上创新者或寻租者的标签。

2026年9月工业出厂价格的结构也说明，不能用一个总指数代替各环节境况。生产资料价格同比上涨6.5\%，生活资料价格同比下降0.1\%。同一条产业链上，原料生产者收入改善和终端厂商利润受挤压可以并存。指数的上涨既可能来自成本冲击，也可能来自需求，单凭同比数字无法识别原因。\cite{PPI2026}

% ARGUMENT P049
\phantomsection\label{arg:P049}
价格冲击的分担还受到合同期限影响。原料可以按日调价，终端产品却可能按季度签约，企业利润就会在过渡期变化；等合同更新后，成本又可能转给客户。只看某一个月，容易把暂时的时点差异理解为永久分配变化。库存会进一步改变影响，因为企业出售的产品可能使用较早购入的原料。

不同家庭面临的消费篮子也不同。食品、能源、住房和服务在开支中的比重不一样，全国平均价格指数不能完整表示每户的生活成本变化。收入较低、基本开支占比较高的家庭，通常更难通过延迟消费调整；拥有较多可选支出的家庭，则可能更容易改变购买组合。这里需要家庭支出资料才能量化。本文只从公开指数提出价格传递与分配问题，不把单一指数的变化直接当成各阶层实际购买力的同幅变化。

\Chart{prices}{2026年9月价格变化的分层}{数据图；同比，\%。CPI、PPI及其分项覆盖对象不同，柱高不可相加。来源：国家统计局2026年10月14日发布。\cite{CPI2026,PPI2026}}

\section{公共服务、免费产品与无酬活动}
没有直接收费，不意味着没有成本。公共教育通过税费和预算筹集资源；免费应用可能通过广告、交易抽成或其他付费用户取得收入；家庭照护消耗时间并影响照护者的就业机会。因此，“免费”首先描述付款环节，而不是社会资源耗费为零。

% ARGUMENT P050
\phantomsection\label{arg:P050}
免费产品的融资来源决定了谁最终承担成本。广告支持的服务由广告主支付，但广告费用可能进入商品价格，也可能减少企业利润；公共服务由税收支持，其负担取决于税制和实际经济转嫁。不能因为使用者没有在入口付款，就把成本完全归于某个看不见的外部主体。

同时，公共融资可以解决个体支付能力不足和外部收益难以收费的问题。例如教育的部分收益由受教育者取得，另一些收益可能通过合作、公共卫生或社会交流扩散。只按个人即时支付意愿供给，可能低估这些作用。但公共支付也需要监督质量与成本，否则“有社会收益”容易成为无上限投入的理由。社会财富的评价应同时看到免费使用带来的机会改善，以及为此使用了哪些稀缺资源。

也不能把所有免费活动都认定为企业雇佣劳动。用户点击、社群交流、业余分享与接受企业指挥的工作，在合同、控制和目的上不同。它们可能给商业组织创造数据价值，但价值形成并不自动给出应有报酬的唯一标准。是否需要报酬、隐私权、可携带权或公共补偿，要结合权利安排和具体损害讨论。

% ARGUMENT P051
\phantomsection\label{arg:P051}
数据活动是否构成劳动，要具体观察目的与控制。受雇标注数据的人按照任务取得报酬，平台用户在交流中留下数据，研究者开放分享材料，这三种活动的时间、动机和合同关系不同。企业能够从中获益，并不足以使三者自动成为同一种雇佣关系；但形式上没有雇佣合同，也不能排除实际存在持续指挥和劳动依赖。

权利问题因此需要分别提出。对隐私的保护涉及个人信息使用边界，对收益分享的讨论涉及贡献、议价与公共规则，对劳动保障的讨论涉及实际控制和风险。把这些问题都归入一个“数字劳动价值”口号，会使具体纠纷难以解决。更可操作的路径是确认谁决定数据用途、谁能撤回授权、谁承担泄露损害，以及某项活动是否实际上构成持续谋生的受控工作。

同样，无酬家务可以通过替代成本或机会成本估值，但两种估计回答不同问题。聘请保姆的市场价格可能低于照护者离职前的工资，也可能无法复制亲属照护的质量。把其中一个估值直接称为家务的“真实价值”，会制造虚假的精确性。时间、任务、受益人和机会约束应先记录，货币估值再作为敏感性分析。

% ARGUMENT P052
\phantomsection\label{arg:P052}
两种估值方法的差异，可以用一个照护者说明。若他原本有较高工资，离职照护的机会成本会较高；市场替代服务的价格可能较低，但服务时间、信任和质量未必相同。对原本没有就业机会的人，机会成本估计可能很低，却不能据此认为其照护不重要。这说明货币估价会受到劳动力市场原有不平等的影响。

因此，家庭再生产的第一步应记录任务和时间，并观察照护是否妨碍必要休息、就业与学习。政策比较可以询问：同样质量的服务采用家庭、公共或市场安排，各自需要多少资源，成本怎样分担。无酬劳动进入可见范围，不等于要求所有亲属关系都变成市场交易。研究的目的，是使长期被隐藏的负担能够被讨论，并让家庭拥有更实际的选择，而不是给每一次关心标价。

\section{价值理论的可检验版本}
本文建议把一个包揽一切的价值概念拆成三个可检验问题。第一，完全劳动投入能在多大程度上解释可重复生产部门的相对价格？第二，成本之外的收益在多大程度上与进入壁垒、产权和稀缺条件有关？第三，市场收入与满足社会需要的贡献在哪些方面系统偏离？三者使用不同数据，也允许不同答案。

% ARGUMENT P053
\phantomsection\label{arg:P053}
三个问题也对应三种可能不同的结论。劳动投入可能对某些标准化制造品的价格具有较强解释力，而对稀缺土地和独特服务较弱；进入限制可能解释部分超额收益，却无法解释所有利润；市场价格可能在普通消费品配置上有效，却在基本医疗和生态损害方面存在明显局限。允许这些结论并存，比要求一个理论在所有对象上给出相同答案更符合研究逻辑。

检验还应避免只选择有利样本。如果只分析价格容易获得的上市企业，就可能漏掉大量小经营者和非市场服务；只分析能够存续的企业，会忽略失败风险。需要说明样本覆盖、退出情况和测量误差。本文并未完成全国价格与劳动系数的实证比较，因此不报告相关系数或显著性。这里提供的是可执行的研究问题，以及不同结果应怎样改变理论用途。

检验第一问需要投入产出表、就业与工时、价格及质量调整，不能只用工资总额代替劳动时间。检验第二问需要企业、市场和制度变化的证据，不能把回归残差直接叫作剥削。第三问涉及福利标准，应公开服务质量、保障和生态等评价维度。若劳动投入在合理调整后缺乏额外解释力，应缩小其价格解释用途；若特定制度下联系稳定，也不应因理论偏好拒绝这一结果。

% ARGUMENT P054
\phantomsection\label{arg:P054}
投入产出表中的金额系数本身已经包含价格。如果直接用它构造劳动指标，再与同一套价格高度相关，就必须检查这种相关是否部分来自共同的金额基础。较严格的研究需要处理物量、价格基期和部门聚合，并在可能时使用不同资料交叉验证。部门划得越粗，内部质量和技术差异越可能被平均掩盖。

租金识别也需要反事实。企业利润高，可能因为竞争受限，也可能因为尚未被模仿的生产能力；政策放松准入以后利润下降，若同时发生需求收缩，就不能全归于竞争。社会需要的评价则更需要公开标准：减少等待时间、提高服务可达性与降低费用，可能比收入增加更接近目标。统计残差、理论概念和伦理判断必须分别说明，才能防止一个无法解释的数字被直接改名为某种确定的社会关系。

这一修订保留商品背后的社会关系分析，拒绝把一个概念同时当作投入计量、市场价格、道德应得与社会福利。概念分开之后，它们之间的张力才能成为研究对象，而不是被一个含混词语提前消除。

% ARGUMENT P055
\phantomsection\label{arg:P055}
在这一基础上，商品分析与阶层分析形成连接。劳动者依靠出售劳动能力取得生活资料，资本控制者通过组织生产和索取剩余扩大积累，资源权利持有者可能通过稀缺条件获得收入，公共部门则参与规则、服务和再分配。同一市场价格对这些主体具有不同含义：它可能是工资购买力的约束、企业成本的组成，也可能是资产估值的依据。

因此，价值理论的更新不应只停留于“价格由多种因素决定”这一宽泛表述。它需要解释这些因素如何被制度组织、哪些人能够影响它们，以及价格信号无法充分表达哪些需要。本章保留社会劳动与社会关系的问题意识，同时把经验价格解释、生产核算和规范评价分开。下一章将沿着其中最重要的一条关系，考察劳动能力如何进入组织过程，以及剩余为什么需要同时从形成和控制两个方面分析。

\chapter{劳动、技术与剩余}
\section{劳动者出售的究竟是什么}
在雇佣关系中，劳动者通常不是出售已经完成的一件产品，而是承诺在一定时间、任务和组织规则下提供劳动能力。企业购买的也不只是若干小时，它还希望获得质量、配合、经验和解决意外问题的能力。劳动过程因此不可能由事前合同完全写尽，监督、信任、激励和协商始终存在。

% ARGUMENT P056
\phantomsection\label{arg:P056}
劳动合同的不完全性，来自生产中的许多情况无法事前逐项写明。设备故障、客户临时修改、同事缺席和质量问题，都需要现场判断。企业因此需要一定管理权限，劳动者也需要在规则之外运用经验。若把劳动当作完全标准化投入，就会漏掉协作、信任和主动判断；若把管理权视为没有边界，又可能把一切临时要求都转成劳动者必须承担的义务。

不同劳动形式的区别，在这里表现得更具体。按件交付的独立经营者通常对方法与进度有较大决定权，但可能承担返工和销售风险；受雇者一般获得较连续报酬，却接受组织指挥。现实中存在大量混合形式，需要观察谁决定价格、任务、工具和工作节奏，而不能只根据合同名称判断。对平台和外包劳动的研究，尤其要把形式自主与实际控制分别记录。

马克思区分劳动与劳动力，揭示合同交换和实际劳动过程之间的差别。这一点在流水线、办公室和平台派单中仍有解释力。工人可能名义上拥有拒绝订单的权利，却因账户评分、收入不稳定或缺少其他工作而难以使用；管理者也可能无法直接观察努力，需要用产量指标替代。权利的正式存在与实际可行性，应分别考察。\cite{Marx1}

% ARGUMENT P057
\phantomsection\label{arg:P057}
劳动力的再生产还包含工作之外的条件。一个人能够持续提供劳动，需要食物、居住、休息、健康、教育以及在一定时期内由他人承担照护。企业只在工作场所购买劳动能力，却依赖这些条件长期存在。若工资和保障不足以维持它们，差额可能由家庭成员、公共服务或劳动者自身健康承担。这里的“差额”是需要具体测量的生活约束，并非可以从任何工资数字直接推算的统一剥削率。

劳动过程中的控制也有不同程度。标准化流水线可以直接规定速度，专业工作可能通过项目期限和绩效评价间接控制。两者都需要劳动者承担执行责任，但后者表面上可能拥有更多自由。是否具有真正自主性，应看劳动者能否改变任务安排、提出异议、获得必要资源，以及拒绝超出约定的要求后是否受到不成比例的损失。自由度不能只由工作地点或是否打卡判断。

但是，不能把所有劳动合同都理解为完全单向控制。稀缺技能、团队不可替代性、行业竞争、公共保障和集体协商会改变议价条件。分析的任务是识别这些差别，而不是从雇佣关系的名称直接推出每一份合同的具体损害。

% ARGUMENT P058
\phantomsection\label{arg:P058}
稀缺技能能够改善谈判位置，但这种优势往往依赖具体组织。某项经验在当前企业非常重要，离开后却难以被其他雇主识别，劳动者就可能拥有内部重要性而缺少外部选择。团队共同创造的能力也不一定归每个人所有，企业可能通过客户关系和知识产权把收益保留在组织中。技能、可迁移性和合同权利因此需要一起研究。

现代劳动经济学关于雇主工资设定能力的研究，为这种分析提供了竞争性视角。Card的综述指出，搜索摩擦和工作差异使企业可能拥有一定工资设定空间，即使一个地区并非只有单一雇主。本文采用这一机制，不将国外估计系数直接用于中国。对于中国劳动者，还需具体考察住房、迁移、保障和信息条件怎样影响换工作的成本。\cite{Card2022}

\section{剩余首先是一个需要分解的余额}
企业销售收入扣除中间投入、折旧、工资、税费和融资等支出后，可能形成会计利润。每个扣除项的定义、确认时间和主体不同，得到的余额也不同。把会计利润直接等同于马克思意义的剩余价值，会混淆生产形成的收入、融资分配、存货估价和资产交易。

% ARGUMENT P059
\phantomsection\label{arg:P059}
利润与现金流的区别，对阶层后果尤其重要。企业可能已经确认销售收入，却尚未收到货款；账面利润为正时，仍可能没有现金支付工资。也可能因为折旧较大而账面利润较低，却拥有充足当期现金。会计确认、实际收付款和持续经营所需资源，分别描述不同问题。

资产出售也会改变利润。企业出售一处过去低价购入的物业，可能获得一次性收益，但这不代表主营生产效率提高。若把这类收益与劳动过程形成的营业盈余混在一起，就会错误解释工资与利润之间的关系。研究剩余时，应先清理生产、融资、重估和资产处置，再讨论分配。对普通职工而言，企业账面盈利并不自动保障工资支付；对小经营者而言，现金链断裂也不必然说明其产品没有长期市场。

在一个简化的增加值账户里，可以写为
\[
VA=W+T_{\text{净}}+D+\Pi_{\text{净}},
\]
其中劳动报酬为$W$，生产税净额为$T_{\text{净}}$，固定资本消耗为$D$，剩余部分为净营业盈余。实际统计还涉及个体经营者的混合收入，不宜将其全部归为资本所得。这个式子是核算分解，不证明剩余的因果来源，更没有自动决定每一项收入应当归谁。

% ARGUMENT P060
\phantomsection\label{arg:P060}
增加值分解中的折旧，指生产中固定资本的消耗，目的在于区分维持现有能力与形成净收入。实际统计中的固定资本消耗与企业税务折旧可能不同，后者受会计和税法规则影响。若某年改变折旧政策，会计利润可能变化，而设备真实消耗未必同步改变。用企业财报讨论社会剩余时，需要处理这一口径差别。

混合收入则说明劳动与资本在小经营者身上难以直接分开。可以按类似岗位工资估算经营者劳动报酬，再将其余部分作为一个分析余额，但结果会依赖替代工资和工作时间假设。本书不采用一个未经验证的全国比例把混合收入强行拆分。对于家庭经营，应尽可能同时报告劳动投入、资本占用、负债和净现金流，让读者看到不同解释所依赖的条件。

经典剩余价值分析的重要问题是：劳动者创造和实现的产品收入，如何在工资之外形成可由他人支配的剩余？当代分析应进一步区分资本维护、承担风险、创新激励、组织管理、市场势力和议价结果。区分不等于为现有分配背书，而是让批评针对具体机制。若不作区分，就无法解释两个利润率相同的企业为何可能具有完全不同的劳动条件。

% ARGUMENT P061
\phantomsection\label{arg:P061}
剩余价值理论与现代利润分解回答的问题并不完全相同。前者把工资劳动与生产资料控制的关系放在中心，追问剩余劳动如何被占有；现代企业分析则可能从风险、组织和市场结构解释观察到的利润。不能因为两者都使用“剩余”，就假定它们的计量对象已经一致。对话需要先说明经典价值账户怎样对应现代增加值和财务账户。

也不能把“承担风险”当作一切利润的充分理由。若利润由私人取得，重大损失却长期由劳动者、供应商或公共部门承担，风险回报的解释就需要重新审查。反过来，确实承担失败损失并创造新的生产能力，也应在评价中得到承认。本文的重点是比较控制权、收益权与最终责任是否匹配，以及劳动者是否具有足够条件参与分配，而不是仅凭收入名称决定其正当性。

社会也不能把全部当期产出立即消费。设备需要替换，技术需要试验，灾害需要准备，儿童、老人和无法就业者需要支持。《哥达纲领批判》对这些扣除的讨论表明，剩余的存在与剩余如何被控制，是不同问题。本文因此把规范性焦点放在扣除是否必要、决策能否问责以及受益是否可以说明。\cite{Gotha}

% ARGUMENT P062
\phantomsection\label{arg:P062}
必要扣除的规模也需要治理。设备更新可以维持长期能力，但管理者可能以更新为名进行低效扩张；风险储备可以保护服务连续性，也可能成为不透明的资金池。判断扣除是否必要，需要公开用途、期限、替代方案和事后结果。即使剩余由公共组织控制，也不能跳过这一步。

不同阶层对扣除的时间偏好可能不同。依靠工资维持当期生活的人，更关心收入能否及时兑现；长期投资者更可能接受延后收益；退休者和需要医疗支持的家庭则可能重视当前保障。社会不能只按其中一类人的偏好安排全部资源。合理程序需要把维护、创新、基本保障和当期消费的竞争关系说明白，并允许受影响群体表达需求。剩余的社会用途由此成为制度问题，而非简单的会计余额处理。

\section{工资不是生产率的一面透明镜子}
企业劳动生产率提高，提供了增加工资、降低价格、提高利润或缩短工时的空间，但具体实现哪一种，取决于竞争、议价和制度。工资可以随生产率上升，也可以在企业市场势力增强时增长较慢。反过来，一些公共服务岗位的工资不宜通过短期可售产出衡量，否则会误判教育和照护的质量。

% ARGUMENT P063
\phantomsection\label{arg:P063}
生产率提高形成的是分配空间，工资增长则是空间如何被使用的结果。企业若面临激烈产品竞争，可能通过降价把收益交给消费者；若劳动者能够转换雇主，企业可能需要提高工资以保留人员；若市场进入受限而劳动者退出成本高，更多收益可能留作利润。三种路径可以同时发生，且比例随行业变化。

公共服务中的生产率尤其需要质量调整。一名护士照顾更多患者，单位时间处理数量提高，但若同时降低照护质量，就不能简单称为效率进步。教育、养老和医疗的产出往往包括安全、理解和持续关注，难以无限压缩时间。工资评价因此应结合责任、技能、风险和服务质量，而不是只比较可计数的任务数量。这也解释了为什么劳动价值与市场工资之间不能建立不加条件的一一对应。

工资差距还受到地区生活成本、培训年限、职业风险、劳动供给与迁移成本影响。不能从工资高就推断社会贡献更大，也不能把工资差距的全部部分都归于歧视或剥夺。辨识不合理差距，需要在相近岗位、技能和条件之间比较，并研究制度改变后的结果。

% ARGUMENT P064
\phantomsection\label{arg:P064}
比较工资还要看补偿性差异。夜班、危险环境和偏远地区可能需要额外报酬，但若劳动者缺少信息或选择，这些风险未必得到充分补偿。观察到高工资，不能自动推断风险已经公平定价；观察到低工资，也不能排除岗位具有稳定、便利或其他非现金条件。需要把报酬与工作条件作为组合比较。

制度性排斥与技能差异也可能共同存在。教育机会不平等可以先形成技能差距，招聘规则再进一步扩大差异。若研究只控制最终学历，可能把前一阶段的制度影响排除在外；若完全不考虑技能，又可能把所有工资差距误认为同工不同酬。因此，应分别研究机会形成、进入岗位和岗位内部报酬三个阶段。这样才能说明改革应针对教育、准入、合同执行还是工资谈判，而不是对所有差距使用同一个解释。

农民工监测调查提供了一个有限而具体的观察。2025年外出农民工月均收入为5774元，本地农民工为4376元。这个差额不是迁移的纯因果收益，因为外出者在年龄、技能、行业和家庭责任上可能与本地就业者不同；外出的住房、交通和分离成本也不同。月均收入按从业月份计算，不能自动乘以十二成为全年收入。\cite{Migrant2025}

% ARGUMENT P065
\phantomsection\label{arg:P065}
迁移比较还涉及选择偏差。能够外出的人可能更年轻、更健康或拥有亲属网络，也可能正因为本地机会较少而外出。留下的人可能承担照护责任，或者在本地拥有土地与经营资源。两组平均收入的差异，同时包含这些构成差别和迁移本身的影响，不能直接解释为“离开家乡就能增加某个比例的收入”。

家庭层面的净收益还需扣除新增居住、交通和分离成本，并考虑留守成员的劳动。外出者收入增加，可能伴随祖辈承担更多育儿任务；若家庭能够共同迁移，则又涉及流入地教育和住房。对农民工社会位置的分析，应沿着就业地与家庭再生产地的连接展开。只有同时观察现金流和时间安排，才能判断迁移扩大了多少选择空间，以及成本是否集中到家庭中缺少发言权的成员身上。

\Chart{migrantpay}{2025年两类农民工的月均收入}{数据图；元／从业月。类别差额不能解释为迁移的因果效应。来源：2025年农民工监测调查报告。\cite{Migrant2025}}

\section{从工时控制走向任务控制}
传统工厂常用在岗时间与产量管理劳动，数字系统则能够细化到配送时限、接单率、转化率和客户评分。精细管理可能减少等待、改善匹配，也可能把需求波动与顾客行为转成劳动者的收入风险。当劳动者需要为无法控制的堵车、商家延误或平台规则变化承担扣款时，风险分配问题比“是否使用算法”更重要。

% ARGUMENT P066
\phantomsection\label{arg:P066}
任务控制把部分管理决定嵌入系统规则。劳动者看到的是订单、评分和奖励门槛，但规则背后仍有目标设定、数据选择和管理责任。算法能够提高匹配效率，也可能把难以度量的服务质量让位于容易计数的速度。若报酬高度依赖完成率，劳动者可能为了维持收入承担额外安全风险。

分析应把技术工具与制度选择分开。平台可以设计异常申诉、合理缓冲和风险共担，也可以把所有延误都归于执行者。两种安排使用的技术未必不同，区别在于谁的成本被计入目标函数。商家、顾客、骑手和平台之间还可能存在信息不对称，不能只由其中一方的记录认定责任。可解释的扣款、纠错渠道和可核验的工作记录，能够使效率改进与责任分担同时接受检查。

评价劳动强度应同时观察有效工作时间、待命时间、工作可预测性和拒绝任务的代价。把上线时长当作全部劳动，会夸大部分人的投入；只计算有报酬的执行时间，又会漏掉必须保持可用的时间。研究应先界定劳动者在哪些时段无法自由安排生活，再说明计量方案。

% ARGUMENT P067
\phantomsection\label{arg:P067}
待命时间是否属于劳动，需要看劳动者能否自由使用这段时间。若必须在指定范围内保持随时响应，虽然没有持续执行任务，生活选择仍受到限制；若可以完全离线并承担没有订单的机会成本，则性质有所不同。现实中还可能在多个平台之间切换，简单把各平台记录时长相加会重复计算。

因此，工时调查应记录同时活动、等待、通勤和执行任务，并说明统计边界。收入按小时计算时，分母选择会明显改变结果。只用付费任务时间可能高估实际时薪，直接用全部上线时间又可能低估。对不同阶层而言，时间可预测性也有独立价值：稳定下班时间有助于安排照护和学习，即使月收入没有变化。劳动质量的评价需要同时看到报酬水平、时间长度与生活可安排程度。

劳资冲突也可能隐藏在供应链里。品牌方缩短交期，承包商把波动转给临时工；大企业延迟付款，小企业压缩工资或增加工时。若只观察最终雇主，就会把由交易链产生的压力误写为一家小企业的孤立选择。责任应依照控制能力和合同影响追踪，不能只看形式雇佣关系。

% ARGUMENT P068
\phantomsection\label{arg:P068}
供应链中的成本传递可以分成三个环节。品牌或大客户改变价格、交期与验收规则，供应商据此调整库存和现金需求，最终雇主再决定工资、工时与用工规模。每一环节都有一定选择空间，但空间可能不对称。拥有多个客户的供应商更容易拒绝不利条件，依赖单一客户者则可能承担更多波动。

这并不免除直接雇主的责任。经营压力可以解释其行为，却不能自动正当化拖欠工资或转嫁安全风险。研究的作用，是在保留直接责任的同时识别上游控制，以便制度对准真正的传导路径。只处罚最末端的小企业，可能无法改变大客户持续延长账期的行为；只强调大企业责任，又可能忽略小企业自身治理。合同透明、及时付款与劳动保障需要相互配合，才可能减少沿链条逐级向弱者传递的成本。

\section{技术收益为什么有多种去向}
自动化可以替代某些任务，也可以创造新的任务、提高产品质量、降低价格并扩大需求。Acemoglu和Restrepo的任务框架强调替代效应与新任务效应之间的竞争。技术提高总生产能力，并不保证某一批劳动者立即受益；新岗位是否出现、劳动者能否转岗、收益是否流向消费，都会改变结果。\cite{Acemoglu2019}

% ARGUMENT P069
\phantomsection\label{arg:P069}
任务框架的重要之处，是把岗位拆成具体活动。一个职业包含资料检索、沟通、判断、执行和责任承担，技术可能替代其中一部分，也可能提高另一部分的价值。岗位名称没有变化，并不说明工作内容没有改变；某项任务能够自动化，也不意味着整个职业立即消失。

生产率收益能否创造新需求，还受到价格和收入分配影响。若效率提高使产品更便宜，消费者购买增加，行业用工可能得到部分补偿；若市场接近饱和，降价也未必带来足够销量。即使全社会出现新岗位，原有劳动者也可能因技能、地点或年龄无法转入。技术的总量收益与特定群体的过渡损失，应同时报告。用长期平均收益否认当期失业成本，会把最需要解释的调整过程省略。

以一家完全假设的零部件企业为例，设备使同等产量所需劳动下降两成。企业可能裁员维持产量，也可能降价扩张，或在订单稳定时缩短工时而维持工资。前者改善企业单位成本但增加失业风险，后两者需要需求、融资和协商条件。仅从设备投入增加，无法判断最终就业效应。

% ARGUMENT P070
\phantomsection\label{arg:P070}
这个假设企业的三种选择，可以进一步比较条件。维持产量并减少人员，适合需求有限且劳动节约可以迅速实现的情境；降价扩张需要订单有弹性、供应链能够配合；缩短工时而保持工资，则需要企业有可分配收益和相应协商安排。技术只改变可行集合，不能单独决定企业选择哪一种。

劳动者的处境还取决于转岗通道。被替代者若能够使用已有技能进入新岗位，损失可能较短；若技能高度专用、当地缺少雇主，失业就可能持续。家庭负债与照护责任会进一步限制搬迁。政策因此不能只统计培训人数，而应观察培训是否对应真实岗位、完成后收入是否恢复，以及没有成功转岗的人得到什么支持。技术转型的公平性，主要体现在这些具体条件中。

技术补贴因此不宜只按购置金额评价。需要观察新增产出和质量、就业转移、培训兑现、工资与工时，以及闲置设备。若补贴主要推动企业争相采购相同设备，而产品需求没有增加，可能加剧价格竞争与债务。若公共培训和可携带保障降低转岗成本，同样技术变化就可能产生较广的社会收益。

% ARGUMENT P071
\phantomsection\label{arg:P071}
设备补贴可能改变企业采购时间，而未必增加原本不会发生的有效投资。企业若只是提前购置计划中的设备，短期投资数据会提高，长期净效果却较小。补贴还可能被设备价格上涨吸收，或者使缺乏运营能力的企业购买不适用的系统。评价需要比较没有补贴时会怎样，而不是把全部采购额都视为政策贡献。

对劳动者的支持也要避免形式化。短期培训证书不能保证可迁移技能，企业承诺转岗不能替代实际岗位与收入记录。公共资金可以将部分支持与设备利用、质量改进和劳动者转型结果联系，但指标过多也会增加行政成本。合理安排应选择少数能反映真实结果的指标，设置后续核查与退出条件。技术收益的分配由此成为可观察的政策效果，而不是写在项目申请中的抽象承诺。

\section{劳动议价权的条件性解释}
本文把劳动者议价能力理解为在不接受当前条件时仍能维持基本生活并寻找替代选择的能力。它包含技能，但不只包含技能。可领取的保障、可获得的住房、跨地区服务、合同执行、及时支付以及集体表达渠道，都会影响拒绝不利合同的实际成本。

% ARGUMENT P072
\phantomsection\label{arg:P072}
外部选择可以理解为劳动者拒绝当前合同以后能够达到的生活状态。这个状态由替代工作、储蓄、家庭支持和公共保障共同决定。两个人技能相近，若一个能够在失业期继续居住和就医，另一个会立即失去基本生活条件，他们的谈判位置就可能不同。工资差距因而不完全来自生产能力差距。

但提高外部选择并不意味着劳动者总会选择离职。更可信的退出可能使现有关系得到改善，企业为了保留人员提高报酬或减少不合理要求。也可能出现另一种结果：低效率企业退出，劳动者暂时面临调整。净效果取决于需求、竞争和保障设计。本文把外部选择作为机制变量，而不把它视为无成本的政策按钮。它需要真实财政与服务资源，也需要可执行的权利。

据此提出一个候选命题：在其他条件相近时，降低失业和迁移的基本生活风险，可能使生产率收益更容易转化为工资和工时改善。机制不是简单要求企业多付钱，而是提高劳动者的外部选择。竞争性解释是，保障改善也可能改变劳动供给和企业用工，净效果必须由数据判断。

% ARGUMENT P073
\phantomsection\label{arg:P073}
保障与就业关系的研究提供了可参考的证据，但需要精确理解。一项关于中国农村制度变化的工作论文，利用特定监督安排改善福利分配的情境，研究保障有效性与迁移选择的联系。这提示保障并不必然使人退出劳动，也可能降低迁移风险。不过，其局部研究设计、对象与制度背景，不能直接外推为所有社会保险政策都会产生相同效果。\cite{Migration2023}

竞争性证据同样必要。最低工资研究中，有中国企业样本发现就业调整，也有其他国家研究强调雇主市场力量使工资提高未必带来同幅就业减少。差异可能来自政策强度、执行、市场结构与样本时期。本文不从这些研究中挑选一个系数作为普遍答案，而将其用于说明：提高劳动收入的工具需要结合市场环境，工资、工时、正式就业和退出都应进入评价。\cite{MinWage2014,Card2022}

如果保障显著改善，而同类劳动者的工资、拖欠情况和工时均未发生预期变化，就要检查执行是否到位、企业是否有足够剩余以及市场势力是否位于其他环节。不能把任何不利观察都解释成政策“不够彻底”。命题的适用边界包括财政能力、覆盖准确性、企业盈利和就业需求。

% ARGUMENT P074
\phantomsection\label{arg:P074}
命题没有得到支持时，首先应区分测量失败与理论失败。资格扩大但实际领取困难，说明政策处理没有充分发生；服务真正改善而结果不变，则需要认真考虑机制较弱或被其他因素抵消。二者都应记录，不能一律用“执行不到位”保护理论。若反复出现服务改善却没有预期结果，命题的适用范围就应缩小。

同样，工资改善也不能单独证明保障提高了议价权。同期订单增长、劳动短缺或行业技术变化，都可能提高报酬。可以比较相近群体的变化，观察换工作成本、拒绝不利任务的能力和劳动纠纷结果是否同步改变。只有中间机制与最终结果相互支持，解释才更可信。本书提出这类检验路径，但没有把未执行的研究写成已经获得的实证结论。

本章保留劳动过程与剩余控制的核心问题，修订把全部非工资收入视为同质结果的做法。真正需要比较的是，生产率提高之后，有多少转成生活改善，有多少用于合理积累，有多少来自权利不对称，以及调整失败时谁承担成本。

% ARGUMENT P075
\phantomsection\label{arg:P075}
本章对剩余问题的处理，最终落在三个相互区分的判断上：生产是否形成了可持续的净产出，净产出由谁决定用途，分配结果是否使劳动和生活能够继续。一个企业增加利润，可能同时改善效率和劳动条件，也可能依靠延长工时或把成本移出企业。仅凭利润率无法区分这两种情形。

各阶层的位置因此应放在劳动过程与家庭条件的连接中观察。高技能受薪者可能拥有较强谈判能力，却仍依赖连续工资；小雇主可能控制员工任务，却受上游和债务约束；资本控制者能够影响分配，但也可能面对真实投资风险。将这些关系具体化，可以避免把复杂性当作无法分析的借口，也避免用简单身份替代实际机制。下一章将进一步讨论企业内部如何配置这些权利。


\chapter{企业、所有权与控制权}
\section{企业为什么不能只按出资人分类}
所有权决定重要权利，却不能单独解释企业行为。比较国企、民企和混合所有制企业时，需要同时观察行业、规模、市场位置、政策任务、融资条件与治理结构。把电网企业和街边商店的利润率放在一起，不能识别所有制的影响；把上市民企与全部国企比较，也会受到上市选择和行业分布影响。

% ARGUMENT P076
\phantomsection\label{arg:P076}
企业分类首先需要区分法律形式与经济功能。一个集团可以拥有多种所有制子公司，同一子公司也可能同时承担商业业务和政策任务。若把集团总资产全部归到一种简单标签下，再用平均利润解释其全部行为，就会把内部差异抹平。研究的单位可以是企业、业务或项目，但必须前后一致。

行业差异尤其重要。自然垄断网络需要较高固定投入，消费服务更多依赖人员和地点，研发企业则可能长时间没有销售收入。相同资产负债率在这些行业中的意义并不相同。对所有制的有效比较，应尽量寻找任务和市场条件接近的对象，再观察治理与融资规则改变了什么。这样的比较可能得到有条件的结论，而不是一张永远有效的国企、民企效率排行榜。

所有制研究至少有三种问题。第一是生产效率：投入相同资源，能否提供质量相近的产品与服务。第二是分配：剩余通过什么渠道归属股东、劳动者、消费者和公众。第三是风险：经营失败时，谁承担可预见和不可预见的损失。三种评价可能给出不同排序，不能把其中一项冒充全部。

% ARGUMENT P077
\phantomsection\label{arg:P077}
效率、分配与风险之间可能出现取舍。一家企业通过扩大规模降低单位成本，同时集中市场入口，消费者短期得到低价，供应商长期议价能力却可能下降。另一家企业利润较低，但维持必要服务和稳定就业；若成本过高，又会挤占其他公共用途。只报告一个指标，无法判断这些结果的综合意义。

因此，比较需要把时间尺度拉开。当期效率改善是否依赖未来维护不足，工资稳定是否由不断新增债务支持，价格降低是否只发生在排除竞争者之前，都应进入分析。对不同阶层而言，同一安排的时间分布也不同：劳动者需要当前收入，投资者关注未来回报，服务使用者则需要长期可靠供给。企业制度的评价应说明这些诉求如何协调，以及无法同时满足时采用什么公开标准。

例如，一家承担普遍服务任务的公共企业，财务利润较低不一定说明效率低；但如果没有独立核算公共任务，也不能把全部低利润解释为社会贡献。一家民营创新企业利润较高，可能来自效率，也可能来自暂时领先或排他收费。评价需要可以检验的反事实：同样任务由其他组织完成，成本、质量和风险会怎样。

% ARGUMENT P078
\phantomsection\label{arg:P078}
反事实并不总能直接观察。一项公共服务只有一家供应者，研究者无法同时看到另一种组织完成同一任务的结果。可以采用相似地区比较、成本构成分析或制度变化前后观察，但每种方法都有条件。地区差异可能来自人口密度和地形，前后变化可能同时受到技术进步影响，单纯成本比较又可能漏掉质量。

较稳妥的做法，是先把任务拆成可核验部分：最低服务范围、可靠性、维护、价格与风险责任，再比较实现这些要求所需资源。如果某种组织形式减少成本却显著降低服务质量，就不能称为同等任务下效率提高；若服务质量相近而成本下降，则为改进提供更强证据。把任务说明清楚，也有助于区分公共资源配置者的决策责任与基层工作人员的执行责任。

\section{剩余控制权与不完全合同}
现代产权理论把所有权与合同未能事先规定时的决策权限联系起来。Grossman和Hart讨论的核心，是资产整合如何改变各方的谈判地位和投资激励。将资产放在一家企业内部，可以减少某些交易摩擦，也可能削弱另一方积累专门技能的动力。所有权因此是一种有收益也有代价的安排。\cite{Grossman1986}

% ARGUMENT P079
\phantomsection\label{arg:P079}
不完全合同可以用专用投资说明。一个供应商为大客户购置只能生产特定零件的设备，投资完成后再寻找其他买家很困难。客户可能利用这种依赖压价，供应商预见到风险后就不愿充分投资。长期合同、共同出资或并入同一企业，都可能缓解问题，但也会改变双方控制权与收益。

整合的代价在于，原供应商失去独立决定和外部市场机会，其主动改进的激励可能减弱。因此，产权理论强调的是权利安排如何影响事前投资，而不只是事后谁拿到利润。把这一点与中国中小企业处境联系起来，可以看到大客户账期、验收权和专用设备要求怎样形成依赖。单靠名义上的自愿合同，未必能消除投资完成后出现的谈判不对称；但依赖本身也不证明交易没有共同收益。

这与劳动过程研究可以对话。出资人不能在合同中预先写尽所有技术问题，劳动者也不能规定每次任务变化的全部条件。管理权有协调作用，但它也使某些调整可以单方面发生。一个好的理论既应解释企业为何需要协调，也应解释控制权为何需要边界和问责。

% ARGUMENT P080
\phantomsection\label{arg:P080}
劳动者同样会进行专用投资。他们学习企业内部系统、建立客户关系、熟悉团队协作，离开后这些经验可能部分贬值。企业也投入培训和组织成本，双方因此形成持续关系。稳定合同可以鼓励这种共同投资，过度随意的解约则可能使双方只关注短期可见产出。

管理权的合理边界，应围绕任务所需协调与劳动者承担的生活风险确定。企业需要调整生产，不意味着可以无限改变工作地点、时间或安全责任；劳动者需要稳定，也不意味着任何岗位都应永久保留。制度可以要求信息、协商、合理预告和过渡补偿，让调整成本进入企业决策。这样既承认组织变化的必要性，也防止便利全部归管理者、成本全部落在缺少选择的劳动者身上。

在中国企业中，还需识别融资机构、控股股东、实际控制人、管理层以及具有正式职责的组织之间的关系。股权比例相同，担保安排、关联交易、章程和监管要求不同，实际控制就可能不同。不能从工商登记表直接读取全部经济权力。

% ARGUMENT P081
\phantomsection\label{arg:P081}
实际控制可以通过几类可观察事项识别：谁任命主要管理者，谁批准重大投资，谁能够改变关联交易，谁掌握融资与担保，谁在经营困难时决定清偿顺序。股权是其中重要依据，但协议、债务契约和组织职责也可能改变控制。不同问题需要不同证据，不能用一个“实际控制人”名称解释所有决定。

控制还可能在正常时期与危机时期发生变化。企业经营稳定时，管理层拥有较大日常权限；债务违约后，债权人可能取得否决或重组权；公共支持介入后，又可能附加新的经营条件。阶层位置因此带有动态性。高级经理的收入与权力，不完全等同于最终所有者；普通股东的剩余索取权，也不等同于决定企业劳动条件的能力。将权利分解，能够更准确地识别谁对具体结果负责。

\section{“像中国一样增长”解释了哪段历史}
Song、Storesletten和Zilibotti提出的模型，用不同企业的生产率与融资约束，解释中国转型过程中资源重新配置、储蓄和对外资产积累的并存。这项研究把所有制差异放进金融摩擦与结构变化中，优点在于提出了连贯机制。它的历史背景和参数化设定，也限制了直接外推到2026年的范围。\cite{Song2011}

% ARGUMENT P082
\phantomsection\label{arg:P082}
转型增长模型提供的一个重要启发，是高储蓄与融资约束可以同时存在。一个经济体总体拥有较多储蓄，并不意味着每家有生产能力的企业都能得到资金。如果金融渠道偏向拥有抵押或特定信用关系的主体，部分高回报项目仍可能受限；受限企业则可能依靠留存利润逐步扩大。总量资金充足与微观资金不足，并不构成逻辑矛盾。

但这种机制的现实强度需要随时期重新测量。企业生产率、银行风险判断和资本市场条件都可能变化。模型中的高生产率企业也未必在所有时期都由同一种所有制代表。本文采用历史模型来理解资源重新配置的可能路径，同时要求对当代企业使用实际融资与经营资料。否则，解释过去的一组分类会变成对当前主体的固定评价，失去历史分析应有的开放性。

如果市场、政策和技术发生变化，旧模型中的企业分类就需要更新。一些民营大型平台具有强融资能力，一些地方国企可能缺乏经营现金流，创新企业也并不总是资本回报更高的企业。应该保留对融资和资源配置的分析，而不是把模型中的两类企业永久对应成现实中的两个阵营。

% ARGUMENT P083
\phantomsection\label{arg:P083}
大型民营企业可能通过规模、品牌和市场入口获得较强融资能力，小型公共企业却可能受到地方财政与经营现金流约束。这说明融资优势不仅来自所有制，还来自组织位置、担保预期和市场力量。研究可以比较贷款条件、债务期限和融资来源，再分析这些条件与风险是否相称。

创新项目则增加了信息问题。研发成果尚未形成时，外部投资者难以判断质量，缺少抵押也不一定意味着项目没有价值。债务、股权、公共研发支持和采购合同，各自承担不同功能。政策若只按企业身份配置资源，可能错过真正受限的创新；若完全依据短期财务回报，也可能低估长期能力。合理选择需要说明支持所针对的市场或组织障碍，以及项目失败时怎样分担损失。

Hsieh和Klenow关于中国与印度制造业资源错配的研究，提供了另一条思路：观察相似行业企业之间的边际收益差异，估计配置改善的潜在空间。但这类估计依赖生产函数、需求和测量假设。企业面临的质量、风险和调整成本差异，也可能被模型解释为扭曲。估计出的潜在收益不是一张可以立即兑现的政策支票。\cite{Hsieh2009}

% ARGUMENT P084
\phantomsection\label{arg:P084}
资源错配研究常把企业边际收益的差异视为线索：若同样追加资源在某处产生更高收益，重新配置可能提高总产出。但观察到收益差异，还需要排除质量、风险、市场加价和调整成本。一个企业正在建设新生产线，短期回报较低，不一定代表资源永久配置错误；一个企业价格较高，也可能反映品牌或市场势力，而非物理生产率更高。

因此，模型中的效率损失应理解为在特定假设下的反事实估计。它有助于提出问题，却不能直接给出关停名单或行政配置方案。政策应进一步观察资金为什么没有流动：信息不足、制度排斥、真实风险还是资源无法迅速转移。若机制判断错误，所谓纠正错配可能变成新的错配。理论测算与实施设计之间，需要具体的组织和市场证据。

两类研究共同启发本文：真正需要解释的是资源为何没有流向更有社会回报的用途，以及重新配置的成本由谁承担。关闭低效企业可能提高平均生产率，但受影响劳动者未必立即进入更好岗位。统计上的效率改善必须与实际迁移、培训和保障联系。

% ARGUMENT P085
\phantomsection\label{arg:P085}
重新配置的成本通常集中在地点和人员上。设备可以出售，工人的技能却未必容易转移；企业可以搬迁，家庭的住房、孩子教育和老人照护不能同时移动。一个行业平均生产率提高，可能部分来自低生产率企业退出，而非所有在岗者获得更好工具。总量改善与原有劳动者收入下降因此可能同时出现。

评价结构调整应区分生产收益与过渡分配。若新产业确实创造更高社会回报，可以支持转移，但需要估算培训、求职、迁移与收入中断的成本，并说明由谁承担。只要求劳动者自行适应，会把调整的一部分真实成本移出企业和政策账户；无期限保留所有旧岗位，又可能阻止新机会形成。可携带保障、及时信息和有期限支持，正是连接效率改善与生活连续性的中间制度。

\section{公众可兑现收益的传递链}
第一章提出的公众可兑现收益，可以进一步写成一条传递链：公共资产形成经营或服务能力，治理结构决定其用途，财政与公共服务制度决定分配，具体居民获得现金、服务、价格优惠或风险保障。任一环节受阻，都可能使名义公共财富与生活改善之间出现距离。

% ARGUMENT P086
\phantomsection\label{arg:P086}
传递链可以用一项公共资本收益安排逐步观察。企业先形成经营现金流，其中一部分用于维持资产和必要投资；可分配部分通过规定程序进入财政或服务安排；财政再决定具体用途；最终使用者通过资格和供给获得利益。任何一环的名义增加，都不能替代后续结果。例如上缴增加后若财政减少其他同类投入，居民得到的净增量可能小于上缴额。

这提示研究需要观察替代效应。新增资金是增加了总服务，还是替代原有预算；低价服务是减轻了家庭负担，还是主要补贴高使用量群体；企业承担风险是否减少社会损失，还是隐藏了不良项目。公众收益的研究对象是整个传递过程，而不是挑选最容易显示增长的一环。对掌握预算和治理权的主体，也应按其实际控制的环节分别问责。

这个概念应优先用分项指标表达。利润上缴按现金记录，公共服务按可获得性与质量记录，价格补偿按覆盖对象和财政成本记录，风险承接按担保条件和实际损失记录。没有可靠估值时，不把四者强行换算为一个总额。尤其不能同时把企业低价供给的全部消费者收益和政府对该企业的补贴都作为新增公共收益。

% ARGUMENT P087
\phantomsection\label{arg:P087}
分项指标需要保持可加性边界。现金分红可以按金额相加，但服务质量和风险保障不宜未经估价直接折成现金。即使采用货币估价，也要避免把企业补贴、服务产出和消费者收益全部相加，因为它们可能描述同一安排的不同侧面。一个公开仪表板可以并列展示这些结果，分析则解释其关联。

同时，平均覆盖率可能掩盖阶层差异。具有信息和时间优势的家庭更容易申请服务，流动人口和照护负担重者可能难以使用。评价需要比较资格人口、实际使用者和未能使用者，并记录费用、等待与距离。公共资源的名义归属相同，实际兑现仍可能不均。把这种差异揭示出来，才能判断制度改善应针对资金不足、供给位置还是办理规则。

考虑一个假设供水企业。它可以把盈利全部留用于扩建，也可以提高维护质量、降低低收入家庭的基本用水负担，或向财政上缴利润。如果扩建服务于真实新增需求，保留利润可能提高未来公共收益；如果管网漏损严重却继续建设低利用率工程，规模扩大就未必改善公共财富。这种比较需要项目数据，不能只看总资产。

% ARGUMENT P088
\phantomsection\label{arg:P088}
供水算例中的三种用途，对应不同时间和分配。扩建增加未来容量，维护提高当前可靠性，基本用水优惠直接减轻特定家庭负担，上缴则把用途决定转入财政。它们不能仅按当期利润比较，也不能假定某一种永远优先。需要观察新增需求、漏损、服务中断和财政机会成本。

若人口增长较快，延迟扩建可能使更多家庭无法接入；若人口稳定而管网老化，维护可能比新增工程更有收益。若基本生活负担较重，有针对性的费用安排可能优于普遍降价。每种判断都有可观察条件。这个例子展示的，是公共企业治理如何通过具体选择改变不同阶层的生活，而不是对任何真实供水企业的评价。真实决策仍需成本、工程和使用资料。

候选命题是：在公共资产规模给定时，收益传递的透明度和预算约束，比名义规模本身更能解释居民获得的服务改善。竞争性解释包括地区原有收入、人口密度和中央转移支付。若控制这些因素后，传递机制没有额外解释力，概念的用途就应限于描述，而不能宣称它是决定规律。

% ARGUMENT P089
\phantomsection\label{arg:P089}
“透明度”本身也需要定义。发布更多文件不一定让公众更容易理解，关键是任务、成本、收益与责任能否被对应起来。若财务报表没有说明公共任务补偿，或者服务报告没有呈现债务与维护，信息数量增加仍可能无法支持判断。有效透明应降低外部主体核查和提出异议的成本。

这一机制可能受到组织能力限制。人口密集、财力较强地区更容易同时实现透明与优质服务，二者的相关性不必然表示前者造成后者。研究应寻找具体规则变化，观察信息改善是否引发预算调整、问责或服务变化。如果公开以后没有任何纠错渠道，结果可能很小。因而，候选命题应把透明、可回应性与执行能力分开，而不能把一个抽象的治理质量指数当作全部原因。

\Chart{rights}{从名义所有权到实际收益}{机制示意图；箭头表示需要分别验证的传递环节，不是已估计的因果系数。图由本文构建。}

\section{中小经营者处于什么位置}
小企业主往往既投入资本又亲自劳动，还可能用家庭住房担保。其经营收入包含本人劳动报酬、资本回报、风险补偿和租金，不能全部视为无需劳动获得的收益。若销售收入不足以支付店主按同类岗位计算的劳动报酬，账面仍有盈余也未必说明经营回报良好。

% ARGUMENT P090
\phantomsection\label{arg:P090}
中小经营者的收入结构，可以通过家庭与企业账户的连接理解。经营者投入存款、承担租约、亲自工作，亏损时可能减少家庭消费，盈利时又可能将收入全部留在店铺周转。账面企业与家庭法律上是否分离，影响失败后责任范围；实际中即使法律上有限责任，个人担保也可能把风险重新连回家庭。

因此，经营规模、雇佣人数与本人劳动依赖应分别记录。没有雇员的维修工，主要出售自己的技能；雇用少数员工但本人仍全日工作的店主，兼有劳动与管理位置；能够脱离日常劳动而持续依靠利润生活的所有者，则具有不同经济基础。用一个“小老板”概念覆盖这些差异，会使税收、保障和劳动政策难以找到真实对象。

与此同时，小企业主可以对雇员行使管理权，并将经营风险转为工资拖欠或临时用工。一个主体同时承受压力并拥有局部权力，并不矛盾。阶级或社会位置分析如果只能给每个人贴一个标签，就难以处理这类关系。

% ARGUMENT P091
\phantomsection\label{arg:P091}
局部控制权意味着一个主体可以在某段关系中较强，在另一段关系中较弱。小企业主面对大客户可能缺少议价能力，面对没有其他工作的临时工却可能决定报酬与任务。两种事实可以同时成立，且经常相互传递：上游压缩账期空间，经营者再通过延迟工资缓解现金压力。

这并不要求研究者在同情与批评之间二选一。可以支持公平交易和及时回款，同时要求履行劳动责任。将压力来源说明白，是为了改进制度，而不是取消责任。Wright关于矛盾性阶级位置的研究，提供了分析交叉控制关系的思想资源；本文借鉴其问题意识，但不把国外分类直接当作中国调查已经确认的群体边界。\cite{Wright1985}

政策也应据此分开。保障小企业及时收款和公平进入市场，主要针对它们面对大客户、平台和融资机构的弱势；劳动保护针对它们对雇员负有的责任。不能以企业困难为由长期取消基本劳动标准，也不能忽略过高固定合规成本对小规模经营的影响。简化程序、提高执行确定性和提供共同服务，可以同时改善两侧条件。

% ARGUMENT P092
\phantomsection\label{arg:P092}
小规模经营面对固定合规成本时，平均负担可能较高。复杂报表、重复证明和不确定检查，会占用经营者本已有限的时间。简化程序、提供公共会计和法律咨询，可以减少这些成本，而无需降低工资支付和安全底线。制度应区分必要保障与可以改进的办理方式。

对员工而言，小企业经营困难并不减少基本生活需要；对经营者而言，骤然增加难以预期的固定支出也可能导致退出。政策可以采用清晰过渡、分级服务和一致执法，帮助双方形成稳定预期。若优惠只按企业规模长期提供，还要防止人为拆分和规模停滞。支持的目标应是改善真实生产与公平交易条件，而不是永久保护任何既有经营形式。

\section{比较治理需要哪些证据}
一份有说服力的企业比较，应当说明样本进入和退出、行业价格、资产估值、政策任务以及融资补贴。税前利润、净利润、营业现金流和资本回报率回答不同问题。依靠应收账款确认的利润，不能自动支付工资；出售资产形成的利润，也不能代表主营业务竞争力。

% ARGUMENT P093
\phantomsection\label{arg:P093}
企业资料的选择还影响对成功的理解。只观察上市或存续企业，容易高估投资回报，因为失败企业已经退出样本；只观察获得支持的企业，又可能低估政策效果，因为它们原本处境更困难。比较需要了解样本是怎样进入的，并尽可能记录退出与并购。

财务指标之间也应相互校验。利润增长但经营现金流持续恶化，可能反映回款困难；资产回报提高但维护支出异常下降，可能透支未来能力；工资总额增加但人员与工时扩张更快，劳动者实际报酬未必改善。这些现象都不是单独定罪的依据，而是进一步核查的线索。企业治理研究应把经营、劳动和风险结果放在一起，才能辨认效率改善与成本转移。

本文没有据此建立新的企业面板回归，因此不能给出国企或民企总体效率的因果排名。可以提出的是研究设计：在行业内寻找治理或融资规则变化，对相近企业比较经营、服务与风险结果，并检查改革之前趋势是否相似。政策同时影响多个条件时，应承认识别困难，而不是通过复杂公式隐藏它。

% ARGUMENT P094
\phantomsection\label{arg:P094}
政策研究中的比较组，应尽量代表没有改革时的可能变化。若改革只选择管理较好、负债较轻的企业，改革后表现更好可能部分来自选择；若只选择陷入困难的企业，简单比较又可能低估效果。需要说明进入规则，检查改革前趋势，并对同期行业和融资变化作出合理处理。

结果指标也应事先确定。只在事后挑选利润、资产或就业中改善最大的一项，会夸大证据。对公共任务企业，应同时报告服务质量与财政成本；对商业企业，应包括现金流、投资、劳动条件和退出。本文没有可支撑这类新估计的微观样本，因此将这些内容作为研究设计。理论的完整性可以体现为明确需要什么证据，而不必伪装成已经完成所有实证工作。

本章对资本所有权的修订，是把法律归属、实际控制、收益传递和损失责任放在同一框架中。资本关系既表现为资产占有，也表现为决定别人何时得到钱、承担什么任务以及能否退出的能力。评价企业制度，最终需要沿这些关系追踪结果。

% ARGUMENT P095
\phantomsection\label{arg:P095}
资本控制在日常生活中经常表现为时间安排的权力。决定何时付款，可以改变供应商的融资需求；决定何时结算，可以改变劳动者的生活连续性；决定合同能否退出，可以影响家庭是否敢于迁移或照护亲属。所有权分析如果只看资产登记，就会漏掉这些持续发生的关系。

这也使企业制度与社会阶层连接起来。公共资源配置者控制预算与准入，企业控制者决定投资与用工，管理者执行并影响具体任务，小经营者在上下游之间承受与传递压力，普通劳动者则依赖报酬与保障。各位置并非永久不变，也不决定个人道德或政治立场。它们说明的是在特定制度关系中，谁拥有何种选择、谁承担何种后果，为下一章分析信用和未来承诺提供基础。

\chapter{货币、信用与资本循环}
\section{货币使生产跨越时间，也创造未来要求}
企业通常先支付工资和原料，再销售产品回收资金。农户、制造商、建筑商和软件企业的周期不同，但都可能面对收支时间不一致。货币与信用使生产可以在收入尚未实现时启动，同时把未来收入转化为今天需要履行的承诺。

% ARGUMENT P096
\phantomsection\label{arg:P096}
生产需要时间，所以收入与支出不在同一时点发生。制造商先采购、生产、运输，再等待客户付款；农业还受到季节影响；软件研发可能在形成产品以前长期支付人员费用。即使最终项目有收益，过程中缺少现金也可能使它无法完成。信用的重要作用，是让未来可实现的收入支持当前必要投入。

但把未来收入提前使用，也意味着必须承担预测错误。订单可能取消，价格可能下降，建设可能延误，家庭成员也可能失去工作。信用因此同时扩大行动能力和失败后果。对资产充足者，错误可能主要减少投资收益；对以住房担保的小经营者，错误可能影响居住；对依靠工资偿债的家庭，收入中断可能迅速压缩基本消费。资本循环需要把时间和社会位置一同纳入，而不能只画一个自动回到起点的箭头。

马克思用货币、生产资料、商品与增殖后的货币之间的循环，强调资本必须经过生产和销售才能完成周转。这一分析提醒我们，技术上生产出来的东西不一定能按预期价格售出，账面价值也不等于现金。更新这套框架，需要加入银行创造存款、证券融资、公共债务和复杂供应链付款。\cite{Marx2}

% ARGUMENT P097
\phantomsection\label{arg:P097}
资本循环中的每次转换都有具体条件。货币购买生产要素，需要可供使用的劳动、设备和材料；生产形成商品，需要组织与技术；商品变成销售收入，需要有效需求、交付和收款。任何一处受阻，都可能让前一阶段投入无法按预期回收。增加生产能力不能自动解决销售，增加信用也不能替代实际交付。

周转速度还影响同样资本能够支持多少活动。及时回款使企业较快再次采购，账期延长则增加资金占用。大型客户将付款时间推迟，可能在自身账上表现为现金改善，在供应商账上却形成额外融资需求。若供应商借款更贵，交易链的总成本可能上升。付款规则由此成为分配权力的一部分：它决定谁提供无形的周转融资，以及谁在收入尚未到账时先承担工资和生活开支。

信用不是对现成财富的简单搬运。银行发放贷款时，通常同时形成贷款资产与借款人的存款。英格兰银行关于现代货币创造的说明清楚区分了这一过程和银行仅把储户已有存款转借出去的图景。这是账户机制解释，并不意味着中国货币制度与英国完全相同。\cite{BOE2014}

% ARGUMENT P098
\phantomsection\label{arg:P098}
用一笔纯粹的记账示例说明，银行发放一百单位贷款时，资产端增加一百单位对借款人的债权，负债端增加一百单位存款；借款人同时取得存款资产和贷款负债。双方净财富并没有因这一步记账自动增加一百。它增加的是当前购买力与未来偿付关系，真实财富仍需通过生产、服务或有用资产形成来改变。

当借款人用存款支付另一家银行的客户时，两家银行之间需要清算；这使单家银行面临流动性约束。借款人后来用存款偿还本金，贷款资产与相关存款会相应收缩。利息支付则涉及收入分配与银行费用，不能与本金消灭简单混同。这个账户例子说明了信用创造的基本机制，但不取代对中国准备金、监管、政策工具和具体结算制度的研究。

银行也不能无限制地创造可持续信用。资本、流动性、支付清算、监管、盈利和合格借款需求都会约束扩张。单家银行创造的存款被转往其他银行后，还需要安排结算资金。贷款最后能否偿还，则取决于借款人的收入或资产处置，而不是放款时的记账本身。

% ARGUMENT P099
\phantomsection\label{arg:P099}
资本约束与流动性约束不是同一回事。银行可能拥有能够覆盖长期损失的资本，却在短期支付中缺少可用结算资金；也可能暂时有充足流动性，但贷款损失已经侵蚀资本。前者需要处理期限与支付，后者需要确认损失和补充承担能力。把二者都称为“没钱”，会导致政策对象混淆。

借款需求也是约束。利率下降以后，企业若看不到可盈利订单，未必愿意新增投资；家庭若已有较重债务，也可能选择还款。信用供给改善与实际借款增加之间，需要预期收益和资产负债表条件。由此不能从银行有放贷能力，直接推断所有有价值项目都获得资金，也不能从贷款增长缓慢，直接推断银行拒绝支持生产。需要分别观察需求、风险和渠道。

\section{本金、利息与利润不能混成一笔}
借款本金增加当期可动用资金，同时增加负债，通常不构成收入。偿还本金减少资金和负债，不等同于生产成本。利息是使用信用所支付的收入转移或服务相关支出，其核算与主体有关。银行利润还要扣除资金成本、运营费用和信用损失，不能把全部利息收入直接当作净利润。

% ARGUMENT P100
\phantomsection\label{arg:P100}
本金与收入的区别，可以用经营贷款解释。企业借入资金购买设备，贷款本身不增加营业收入；设备在使用中逐步消耗，其资本消耗进入相应生产核算；企业以后偿还本金，是清偿融资承诺。若把借款和偿还都当作生产收入与成本，融资规模变化就会被误记为产出变化。

金融服务的核算比简单利息转移更复杂，部分利差对应支付、筛选和管理等中介服务。本文在讨论家庭现金压力时记录实际付款，在讨论增加值时遵循相应统计口径，在讨论风险时再记录信用损失。三种观察不应随意相加。对读者而言，最关键的是先判断正在回答哪个问题：生产了多少、当期能用多少现金，还是未来可能损失多少。

对一个家庭而言，月供中的本金和利息都会占用现金，所以生活压力研究应计算两者；对全国收入核算而言，两者不能按同样方式计入。这个差别解释了为什么现金流紧张和净资产为正可以同时出现。一户拥有昂贵住房的家庭，仍可能因失业无法按期还款。

% ARGUMENT P101
\phantomsection\label{arg:P101}
家庭的资产负债表与现金预算可能给出不同信号。住房市值高于贷款余额，意味着按当前估值净资产为正；每月收入不足以覆盖还款与必要生活支出，则意味着流动性紧张。卖房理论上可以释放资金，但出售需要时间，成交价可能低于估值，搬迁还带来新的居住成本。净资产不能自动转化为当期支付能力。

不同阶层的差别，部分体现为谁能够跨过这种时间缺口。拥有流动金融资产、稳定养老金或亲属支持的家庭，可以等待收入恢复；缺少这些条件的家庭可能被迫出售资产、借入高成本资金或减少医疗支出。因而，研究债务脆弱性应同时观察资产流动性、还款期限、收入稳定性和基本开支。只按住房价值判断家庭富裕，可能低估真实风险。

资本循环研究因此应同时看存量、流量与到期结构。债务余额高说明未来承诺较大，但若期限长、利率低、收入稳定，短期压力可能较小；余额较低的企业也可能因几笔短债同时到期而陷入危机。只比较债务与GDP，会漏掉币种、期限、利率和现金来源。

% ARGUMENT P102
\phantomsection\label{arg:P102}
到期结构可以用两个假设企业比较。两家债务余额相同，一家在较长时期内分散偿还，另一家必须在短时间集中偿还；即使全年利润相近，后者对再融资的依赖也更强。若市场突然收紧，它可能在项目长期仍有价值时遭遇支付困难。这说明债务比率不能单独决定危机发生时间。

币种与利率也会改变承诺。收入主要以人民币取得、债务却以外币计价的主体，会受到汇率变化影响；浮动利率债务则把利率调整传到未来付款。公共部门和大型企业可能拥有更复杂的风险管理工具，小经营者与家庭未必能够使用。政策研究应区分这些暴露，不把所有借款主体视为具有相同金融能力。信用的分配不仅是额度分配，也是期限、价格与风险条款的分配。

\section{2026年信用结构中的观察与限制}
截至本书检索截止时已核实的8月报告显示，社会融资规模存量为464.8万亿元，其中政府债券为103.69万亿元；前八个月住户贷款减少1.03万亿元，企事业单位贷款增加11.26万亿元。这里前两项是期末存量，后两项是期间净增减，不能并成同一张总量账。报告由人民银行发布，本文使用新华网完整转载核验。\cite{Money2026}

% ARGUMENT P103
\phantomsection\label{arg:P103}
存量与增量的区别决定了这些数字怎样解读。存量记录截至某一时点仍未消失的融资，期间增量记录新增、偿还和调整共同形成的净变化。一个部门融资存量很大，但当期增量可以很小；另一个部门存量较小，却可能快速增加。把二者放在同一图中比较而不注明单位，会制造误导。

本书对资料截止采用保守处理：八月报告是本次能够核实正文与具体数值的版本，检索未取得可确认在截止日前发布的九月同口径全文，因此不自行推算下一期。这个限制意味着信用结构部分描述的是已核实观察期，而非声称覆盖截至十月十四日每一项最新发布。后续若取得早于截止日的可靠原表，应按版本记录更新，不能只替换数字而保留旧的解释。

\Chart{credit}{2026年前八个月贷款增减的部门差异}{数据图；万亿元，人民币贷款净增减。住户与企事业单位并非贷款全部分类，图中两柱之和不等于全部新增贷款。\cite{Money2026}}

这些数据与企业融资扩张、家庭去杠杆可以并存的解释相容，却不足以证明每个家庭都在主动还债。净减少还可能涉及新贷款减少、到期偿还、核销和统计调整；需要更细分的数据才能辨认。企业贷款增加也不自动等于新增生产投资，因为资金可用于周转、置换债务或购买资产。

% ARGUMENT P104
\phantomsection\label{arg:P104}
住户贷款净减少可以来自多种家庭行为。部分家庭可能提前还贷，部分家庭可能减少新购房借款，也可能只是到期偿还超过新增。不同原因对应不同需求后果：提前还贷可能降低未来利息负担，但压缩当期流动资金；新增借款减少可能反映住房需求变化，也可能反映融资条件。总量净额无法区分这些路径。

企业贷款增加同样需要看用途。用于支付被拖欠的供应链款项，可能改善多个企业现金流；用于置换高成本债务，主要改变未来偿付；用于购买旧资产，则未必增加当期生产能力。若要判断信用是否支持实体活动，应追踪资金用途、项目产出和回款，而不只比较贷款增长与投资增长。两者之间的差额不能自动称为资金空转，因为周转、期限和统计范围本来就不同。

社会融资规模包含非金融企业境内股票融资，不能直接命名为全社会债务。M2是特定口径的货币存量，也不是居民可以共同取出后购买商品的“闲钱总额”。借款形成的存款、交易周转和金融机构资金安排都会影响它。用M2除以人口得出每个人“应有的钱”，没有分配意义。

% ARGUMENT P105
\phantomsection\label{arg:P105}
货币存量具有对应的资产与负债，分布在不同部门和主体之间。一个人的存款可能是企业支付工资形成的，企业存款可能等待采购，金融机构资金还涉及特定统计范围。把总量平均分配到人口，只得到一个没有实际产权含义的算术值。它既不表示每个人应取得的收入，也不表示所有存款能够在不改变价格和资产负债关系的情况下同时转成消费。

股票融资进入社会融资规模，也说明融资并不全部是债务。股权没有与贷款相同的固定偿付结构，风险和控制安排不同。研究总体杠杆时，应使用适合的问题口径，再解释公共企业、非金融企业与家庭的边界。概念准确并非文字上的挑剔，它直接决定我们是否把风险算重、算漏，或者把一种融资变化误判为另一种经济行为。

国际比较可用BIS非金融部门信贷统计，但必须核对部门与工具覆盖。其企业分类和公共企业处理，与日常语言中的“私人企业”不完全一致。本文未把没有核实观察期的最新比率塞入图表，也不将BIS转录的中国基础数据视为独立调查。\cite{BISCredit}

% ARGUMENT P106
\phantomsection\label{arg:P106}
国际数据库的价值，在于提供较一致的分类与跨国比较入口，但一致不等于完全相同。各国基础统计、企业部门边界、债券估值与历史修订可能不同。某些公共企业在国民账户中属于非金融企业，而日常语言容易把它们全部归为政府。若不读方法，就可能把同一债务在企业与公共部门讨论中重复计算。

因此，国际比较先回答覆盖问题，再比较水平。若一个来源包括更广泛工具，较高比率未必表示真实负担更大；若采用市场价值，利率变化也可能改变债务估值。本文使用BIS方法说明这些差异，不以机构名称替代核查。国外机构整理中国官方资料时，提供的是可比框架和方法处理，独立性应按实际基础资料判断，而不是按发布地判断。

\section{金融收益何时服务生产，何时形成租金}
利息可以补偿时间、违约风险、流动性和中介服务；但这并不证明任何观察到的利差都合理。借款者缺少替代渠道、信息不透明或受到捆绑销售时，融资价格也可能包含市场势力收益。区分这两部分，需要比较违约损失、服务成本和可替代融资，而不是根据利率高低进行道德归类。

% ARGUMENT P107
\phantomsection\label{arg:P107}
利率中的风险补偿，只有结合预期损失才能理解。借款人违约概率较高、回收率较低，贷款者可能要求较高价格；但若价格本身过高，又可能使原本能够经营的项目无法持续。信息不对称与市场力量会使这个关系更复杂，不能简单认为高利率总是准确反映高风险。

低利率也不必然意味着融资公平。若只有少数主体能够获得低成本资金，而其他有合理项目的主体被排除，价格优势就可能强化资源集中。评价需要比较同等风险条件下的可获得性，以及担保或公共支持由谁承担。对家庭和小企业而言，手续费、抵押要求、提前还款条件和强制搭售，都属于真实融资成本。研究金融收益应把这些条款放在一起，避免只围绕名义利率讨论。

股权融资不要求固定还本，却把未来收益和部分控制权交给投资者。它适合一些高不确定性项目，但也可能形成短期估值压力。债务融资让原股东保留更多上行收益，同时使固定偿付要求压在经营之上。企业为什么选择某种融资，取决于信息、控制、税制和市场条件。

% ARGUMENT P108
\phantomsection\label{arg:P108}
股权与债务对控制权的影响不同。股权投资者承担较多经营波动，通常要求分享上行收益和一定治理权；债权人主要要求约定付款，但在违约或契约触发时也可能取得重要影响。企业融资不是只比较哪个利率低，还涉及创始人愿意让出多少控制，以及项目收益能否支持固定承诺。

高不确定性研发更适合具有损失吸收能力的资金，但股权市场也可能追逐短期估值，促使企业优先展示容易吸引投资的指标。稳定现金流项目使用债务可能有效，过高负债则增加经营对短期收入的依赖。没有一种融资工具天然代表先进或落后。应根据生产周期、信息和风险分担选择，并观察普通劳动者是否在融资安排失败时承担了原本未参与决定的损失。

金融机构的社会功能，是筛选、监督、支付、分担和跨期协调。评价金融发展不能只看资产规模，也不能把金融收入一概视为脱离生产的寄生。更具体的问题是：同样融资成本下，真实有用的项目是否更容易获得资金；失败能否有序处理；家庭是否被诱导承担难以理解的风险。

% ARGUMENT P109
\phantomsection\label{arg:P109}
金融服务是否有用，可以通过交易成本与风险改善观察。更可靠的支付减少现金保管和结算摩擦，更好的信息使有能力的借款者获得资金，有序重组避免有价值资产被仓促出售。这些作用可以支持真实生产，不能因为它们没有制造物质产品就被忽略。

但金融资产增长也可能主要反映更多中间层和相互索取，而没有相应改善最终使用者的融资。判断需要看扣除费用后的结果：家庭是否更理解风险，企业是否得到适合期限的资金，失败是否及时确认。专业能力较弱的家庭容易依赖销售渠道，其选择并不总能形成有效约束。披露、适当性和可执行的责任，因而是金融市场运行条件的一部分，而非市场之外附加的道德要求。

\section{风险的传递矩阵}
设一家企业因项目失败产生10单位经济损失。股东吸收4，银行确认3，供应商承担2，财政承担1，这四项构成同一损失的分担，不能再加上“企业损失10”得到20。如果银行随后获得财政注资，其中一部分风险又转移到公共部门，应更新最终承担者，而不是重复累计。

% ARGUMENT P110
\phantomsection\label{arg:P110}
这一损失算例采用的是事后经济损失，不是所有合同名义金额。项目原有资产、债务和预期收益都可能很大，但最终无法回收的部分才进入损失分配。若银行将贷款展期而没有确认减值，账面损失可能暂时较小，经济损失却未必消失。反过来，悲观估值后来恢复，也可能减少最终损失。

财政注资还需要区分资金支出与净损失。政府投入资金并取得有价值股权，不能把全部投入立即算成损失；应扣除可回收价值并考虑风险。对供应商和劳动者而言，拖欠造成的融资与生活成本也可能超出名义欠款的时间价值。完整矩阵应记录初始损失、后续转移和最终承担，分别处理真实资源损失、金融估值与流动性影响，避免把所有不利变化简单叠加。

风险转移可以减少总损失。展期帮助仍有经营能力的企业跨过短期冲击，可能避免低价抛售和停工；公共担保使有社会价值但缺抵押的项目得以开展，也可能提高真实收益。风险转移也可能仅推迟确认，令没有前景的项目继续消耗资源。关键区别是预期可回收价值是否改变，而不是账面到期日是否延后。

% ARGUMENT P111
\phantomsection\label{arg:P111}
展期是否有价值，取决于企业未来可产生的现金流和立即清算的损失。如果困难来自暂时订单中断，继续运营可能保存客户、团队和专用设备；若产品长期失去需求，展期只是增加新的投入。判断需要公开经营假设，并定期检查，而不是以“稳定”或“市场出清”作为无需论证的结论。

风险转移还会改变未来激励。一次有条件支持可以帮助跨过异常冲击，反复无条件支持则可能使管理者预期失败不由自己承担。政策因此需要将维持必要功能与调整原有权利结合：例如要求资本承担损失、限制新增高风险项目、披露现金用途。普通家庭也可能需要跨期支持，但应围绕基本生活和恢复能力设计，避免让救助主要补偿高风险资产持有者的预期收益。

\Chart{risk}{同一损失的初始分担与再次转移}{假设示意图；单位为任意记账单位。银行承担的3中若有2最终由财政吸收，最终分担为股东4、银行1、供应商2、财政3，总计仍为10。不代表中国实际损失结构。}

候选命题是，控制决策但可以向外转移损失的组织，更可能选择对自身有利而对整体过度冒险的项目。这与软预算约束和道德风险已有研究相连，本文不把它冒称新规律。新增的研究要求是绘制中国具体制度下的传递矩阵，识别最终落在税收、公共服务、工资、存款收益或资产重估上的渠道。

% ARGUMENT P112
\phantomsection\label{arg:P112}
风险矩阵中的位置，不应按名义身份预先填写。公共部门有时承担最终损失，也可能通过征税、费用或服务调整将部分负担传给居民；银行确认减值后，损失可能由资本、未来收益或公共支持吸收；企业股东承担股权损失，但若存在个人担保，其责任又可能延伸到家庭。每条路径都需要合同和财务资料。

本文的候选命题进一步要求观察决策之前的信息。若一个主体知道可以把损失转给别人，同时能够保留成功收益，就可能提高冒险倾向；若强问责、声誉或职业风险抵消这种激励，结果可能不同。因而，风险转移能力只是条件之一。研究应比较不同约束下的项目选择，不把所有得到支持的主体都视为过度冒险者，也不把最终损失的归属反推为最初已经确定的意图。

\section{危机预警应看承诺能否兑现}
Minsky区分能够用经营现金流支付本息、需要滚动本金，以及连利息都依赖再融资的状态。这套分类有助于观察融资脆弱性，但不能把最后一种状态当作对借款者的道德指控，也不能未经数据便把整个经济归入其中。初创研发项目可能暂时没有现金流，其合理性取决于股权支持和前景。\cite{Minsky1992}

% ARGUMENT P113
\phantomsection\label{arg:P113}
融资状态可以随经济条件改变。同一企业在订单稳定时能够支付本息，需求下降后可能只能支付利息，再融资中断后又可能需要出售资产。状态不是永久身份，而是现金流、承诺与市场条件的关系。利率变化、账期延长和抵押价值下降，都可能推动这种转换。

初创研发没有当期收入，也不能直接与依靠新增借款维持旧债的成熟企业混为一谈。前者若由能够承受失败的股权资金支持，固定偿付压力可能较小；后者若必须持续借新还旧，则更依赖市场信心。研究应观察融资工具是否与项目不确定性匹配。Minsky的分类提供了现金流视角，但判断具体中国主体仍需到期结构、经营资料和责任安排，不能仅凭某个负债总额下结论。

更可操作的预警包括经营现金流与到期偿付的差额、应收款期限、担保链、融资来源集中度和抵押物相关性。若多个部门用同一类资产作安全基础，资产价格下降就可能同时削弱它们的融资。宏观总资产看似充足，也不保证在同一时刻都能变现。

% ARGUMENT P114
\phantomsection\label{arg:P114}
共同抵押物会形成相关风险。若银行、企业和家庭都依赖同一地区房地产估值，价格下降可能同时减少抵押能力、削弱销售与增加谨慎储蓄。各方单独看似拥有保障，合起来却依赖相同收入来源。分散法律主体并不等于分散经济风险。

预警还应关注反应是否相互加强。企业为了偿债出售资产，可能压低其他企业抵押价值；家庭减少消费，可能进一步影响企业收入；地方缩减支出，又可能加剧回款困难。单个主体的合理自保汇总以后，会改变其他主体面对的条件。政策需要识别哪些连接具有系统性，并以公开、有限和可核查的方式维持必要支付。若只把压力从一个账户推向另一个账户，总体脆弱性可能并未下降。

本章因此把货币资本更新为一套组织未来资源要求的制度。创造信用能推动生产，不能替代生产；延长承诺能争取时间，不能自动消除损失。金融稳定的目标，应是让必要支付与有用生产持续，同时使失败的真实成本得到可解释的分担。

% ARGUMENT P115
\phantomsection\label{arg:P115}
货币制度的政治经济学含义，在于它分配跨越时间的能力。能够取得稳定融资的主体，可以等待收益形成、抵御暂时冲击并保留资产；无法取得融资者，即使拥有技能或有用项目，也可能被迫接受不利交易。信用因此既是生产协调工具，也是社会位置形成的一部分。

不过，扩大信用并不能代替收入与保障改革。家庭如果缺少持续现金来源，更多借款可能只是延迟生活压力；地方如果没有稳定收入和有用项目，低利率也不能永久解决债务。金融稳定需要真实生产、合理分配和可信责任共同支持。下一章将考察中国经济中最重要的连接之一：土地与住房如何把家庭生活、企业融资、银行资产和地方财政放在同一个循环中。


\chapter{土地、住房与地方财政}
\section{土地不是普通的可复制商品}
机器可以增产，土地的自然位置不能复制。城市扩张、交通建设和公共服务，可以改变某块土地的用途与收益，但不会消除位置的稀缺性。土地价格因此不仅反映开发劳动，也反映规划、交通、公共投资、邻近活动和对未来收益的预期。

% ARGUMENT P116
\phantomsection\label{arg:P116}
土地的稀缺性并不只来自总面积有限，更来自某个位置与就业、交通和公共服务的结合。远离岗位的空地与交通便利的城市地块，即使面积相同，也不能互相替代。开发新城可以增加建筑供给，但形成成熟生活环境还需要企业、学校、医疗和居民活动共同进入，因而存在时间与协调成本。

位置收益也会随条件变化。新交通改善通达性，可能提高需求；人口与产业迁移则可能使过去稀缺的位置失去优势。因此，地价不是一种由自然永久保证的收入。资本化把未来区位收益折为当前价格，预期变化便会引起资产重估。研究中国城市住房时，应同时观察真实居住需求和未来收益预期，而不能把土地不可复制直接推成房价只能上涨。

中国大陆的土地制度区分国家所有与农民集体所有，并通过用途管制、征收和使用权安排组织开发。土地管理法第九条规定城市市区土地属于国家所有，农村和城市郊区土地除法律规定属于国家所有的以外，属于农民集体所有。不能把城市住房交易写成土地所有权的完全私人买卖，也不能把集体所有简化为成员可以随时各自出售的一份土地。\cite{LandLaw}

% ARGUMENT P117
\phantomsection\label{arg:P117}
土地所有权、使用权、经营权与收益权需要分别讨论。法律上的公共或集体归属，并不意味着所有居民或集体成员都直接决定每次用途变化；具体权限还通过规划、组织程序和合同实现。个人取得住房或农业经营权益，也不等于取得不受用途和期限约束的完整土地所有权。

这些制度安排影响不同主体的退出能力。城市购房者可以在规则允许范围内转让房屋权益，农村家庭的承包、宅基地与集体成员资格则具有不同条件。将两者直接视为同一种可自由抵押和出售的资产，会夸大部分家庭的可用财富。本文仅据全国法律说明基本框架；涉及具体交易和补偿时，还需核对当地有效规则。阶层分析应关注权利能否实际使用、收益怎样分配，以及规则变化是否让某些主体承担集中损失。

法律边界影响市场进入、开发收益和补偿分配。研究土地资本化，应分别观察原有使用者、集体组织、地方政府、开发商、购房者和金融机构的权利变化。仅看出让成交价，无法知道增值收益中多少来自公共投资、多少来自用途变化，又有多少由原有使用者获得。

% ARGUMENT P118
\phantomsection\label{arg:P118}
土地增值的形成，可以按过程分解。原有用途产生一组收益，规划允许新的用途，公共设施改善可达性，企业与居民进入形成需求，融资条件再把未来收益资本化。成交价是这些因素共同作用的结果，不能全部归于某一次开发劳动，也不能全部认定为某一主体无成本取得的收益。

原有使用者可能获得补偿和安置，也可能失去生产、居住或社区关系。补偿金额能否弥补这些变化，需要看后续就业、居住成本和长期保障。一次性现金高低不能完全代表最终生活改善。对地方和开发商而言，项目收益也要扣除建设、维护与风险。把各方放在同一时间线上，才能比较谁先取得现金、谁承担长期义务，以及增值收益是否与实际投入和损失相称。

\section{地租理论需要加入公共投资和财政组织}
经典地租分析把土地的差异性和所有权对收益的影响纳入资本分配。其持久问题是：不依靠当期新增劳动的稀缺条件，如何形成持续索取收入的权利？在现代城市，这个问题需要加入公共基础设施、规划权、信用条件和居民对教育医疗的需求。\cite{Marx3}

% ARGUMENT P119
\phantomsection\label{arg:P119}
地租研究把稀缺条件与排他权利连接起来。在城市中，便利位置的收益来自许多主体共同活动，却可以通过特定产权和使用安排被资本化。公共投资、企业集聚和居民需求共同提高某地价值，收益却未必按贡献自动分配。这是地租问题持续具有解释力的原因。

中国土地财政的历史还需要财政制度背景。1994年起实施的分税制，重新安排中央与地方税收及转移支付关系；它改变了地方收入结构，但不能单独解释后来每一轮土地开发。城市化、住房改革、金融条件、产业竞争和地方组织能力共同形成路径。若把土地财政全部归因于一次财政改革，就无法解释不同地方和不同时期的差异。本文采用多条件形成过程，并将财政制度原文与后续机制推论分开。\cite{Tax1994}

一条新轨道交通可能提高沿线便利性和房产价格。价格增加不全是业主本人投资的结果，也不全是可以无成本收回的“公共创造价值”：居民可能早已通过税费与票价承担一部分成本，开发商也可能支付配套费用。若讨论回收增值收益，就必须说明成本分摊、初始购买价格和过渡安排。

% ARGUMENT P120
\phantomsection\label{arg:P120}
公共投资提高地价，并不意味着投资成本可以全部通过周边资产持有者回收。项目可能服务更大范围人口，部分收益体现为时间节省和就业机会，部分才进入房价。若只向某一小范围征收，可能忽略跨区受益；若全部由一般财政承担，又可能使没有资产的人补贴特定物业升值。

可讨论的安排包括使用者付费、受益范围内的合理成本分担和一般税收支持，但每种方式都需要估值与过渡。较早购入资产的人可能获得意外增值，较晚购买者却已经在价格中支付了这部分预期；同样税负对两者影响不同。现金收入较低的自住家庭还可能面对流动性问题。因此，回收公共创造价值需要明确时间、对象和支付能力，不能只凭一个总增值数字制定统一负担。

土地出让可以把未来较长时期的部分使用收益集中为当期收入。这样有助于快速建设，却也可能使本届可用资金与未来维护责任分离。若地方把一次性收入当作可持续收入安排经常支出，房地产市场转弱时就会出现约束。问题不在于土地出让收入天然无用，而在于收入期限与支出期限是否匹配。

% ARGUMENT P121
\phantomsection\label{arg:P121}
一次性收入与经常支出的错配，可以用城市建设周期解释。土地开发带来当期资金，支持道路和设施建设；设施建成以后，保养、人员与服务支出会持续发生。如果未来税源和使用收入没有同步形成，早期扩张就可能留下长期缺口。建设规模越大，并不必然意味着持续财政能力越强。

1998年的住房改革明确提出停止住房实物分配、逐步实行货币化，并发展住房金融与交易市场。这一制度转变扩大了家庭通过市场取得住房的作用，也使购房时间、单位历史安排和融资条件成为代际资产差异的重要来源。本文将其作为历史转折，不据文件目标推定所有地区实施效果相同。后来住房与土地收入、银行融资相互连接，需要继续通过市场与财政资料解释。\cite{Housing1998}

\section{房地产连接四张资产负债表}
房地产开发同时联系开发商预售回款、居民按揭、银行贷款和地方土地相关收入。在扩张阶段，销量和价格上升可以改善各方抵押与收入预期；在调整阶段，它们又可能同步承压。高度相关的安全来源，并不等于四份独立保障。

% ARGUMENT P122
\phantomsection\label{arg:P122}
四张资产负债表之所以联动，是因为它们依赖部分共同现金流。开发商需要销售回款，地方需要土地交易和相关税收，银行依赖开发与按揭偿付，家庭则依赖收入和房屋可使用价值。销售转弱时，开发商减少拿地，地方收入变化，项目建设与就业又可能受影响，从而反馈到家庭购房能力。

但这种循环的强度取决于具体条件。企业是否分散经营、地方税源是否多样、银行贷款是否集中、家庭杠杆是否较高，都会改变冲击传播。一个城市价格调整不能直接代表全国统一危机。研究应识别连接最紧密的地区与主体，并区分资产价格调整、项目现金缺口和家庭生活风险。四张表提供传导结构，实际损失仍需逐项核实。

对居民而言，住房兼有居住服务、储蓄载体和杠杆资产的性质。自住房价格下跌不立即减少其可居住面积，却可能影响换房、消费信心和抵押能力；租房家庭可能受益于租金或购房门槛下降，也可能因建筑和相关行业收缩失去工作。不能用“居民财富减少”概括所有家庭的福利变化。

% ARGUMENT P123
\phantomsection\label{arg:P123}
住房财富效应具有明显异质性。拥有多套出租住房的人，可能同时关注租金与资产价格；只有自住房的家庭，更关注居住稳定和换房成本；高负债新购房者的净资产对价格变化更敏感；租房者则可能希望租金与购房门槛下降。职业相同、工资相近的家庭，也会因这些条件形成不同利益。

还应区分价格变化与现金收入变化。房价下降而工作稳定，家庭可能继续居住并按期偿还；房价不变但失业，现金压力也可能很大。因此，政策如果只关注价格稳定，可能遗漏收入和交付问题。对社会福利更直接的尺度，是基本居住能否维持、流动人口能否接近工作、债务调整是否有序，以及资产损失是否被转给缺少决策权的家庭成员。

对开发商而言，预售收入使建设得以融资，同时形成交付义务。项目不能交付时，住户可能同时承担按揭与额外住房支出。解决问题需要区分项目建设资金缺口、企业集团债务和股东损失。以保护居住功能为目标，与保护原有股东全部权益，不是同一个目标。

% ARGUMENT P124
\phantomsection\label{arg:P124}
预售把部分融资风险放在购房者身上。购房者支付资金并承担按揭，是基于未来交付承诺；如果资金被用于其他项目或集团债务，特定项目的交付能力就可能与已收款项分离。分析需要追踪项目资金，而不能只看集团总资产是否足够。

处理未交付住房时，首先应估算完成建设需要的真实资源，再判断项目可回收价值和原有资金去向。公共支持若进入项目，应明确取得什么权利、谁先承担损失、未来如何回收。保护家庭居住和维护支付秩序可以成为政策目标，但不能自动扩大为对所有关联债务的担保。区分项目、集团与股东，使支持能够对准必要功能，同时减少将原有投资风险无差别转入公共账户的可能。

对银行而言，抵押物价格和借款人收入共同决定回收能力。只要借款人持续偿付，账面抵押比率变化不一定立即成为损失；但失业与房价同时下降，会放大风险。对地方而言，成交量、土地价格、相关税收和建设支出联动，地方融资平台还可能增加隐含联系。并非每个城市的强度都相同。

% ARGUMENT P125
\phantomsection\label{arg:P125}
银行风险同时取决于抵押与现金流。抵押物提供违约后的回收来源，正常偿付则主要依靠借款人收入。房价下降并不立即等于贷款损失，但会减少缓冲；若收入同时受损，回收压力就更大。地方经济越依赖建设和房地产，家庭收入与住房价格越可能一起变化。

地方平台和公共项目还可能通过担保与土地预期增加关联。Bai、Hsieh和Song关于财政扩张的历史研究，说明融资平台如何在特定时期延续表外投资渠道；本文用它理解组织路径，不把其历史估计改写为2026年规模。要判断当前风险，需要新的债务、项目和现金流资料。历史机制可以提示查找方向，却不能代替当期测量。\cite{Bai2016}

\section{财政数据说明了什么}
2026年1至8月，地方一般公共预算本级收入同比增长3.1\%，国有土地使用权出让收入同比下降28.6\%。两类收入分属不同预算，不能合成一个未经说明的“地方税收”指标。同期全国一般公共预算债务付息支出增长5.4\%，提示偿债支出与其他用途之间需要安排，但不能单凭增速断言公共服务已经全面削减。\cite{Fiscal2026}

% ARGUMENT P126
\phantomsection\label{arg:P126}
土地收入下降以后，财政调整并没有唯一方向。地方可以减少新增项目、动用结余、争取转移支付、调整资产经营、延缓付款或改变服务支出。每种选择会影响不同群体：减少低效建设主要影响相关投资链条，拖欠采购会压缩供应商现金，削减服务可能增加家庭自付与照护时间。

因此，从收入变化到阶层后果，需要至少观察支出结构、付款进度和服务结果。全国收入增速只能说明约束有所变化，不能直接证明某类群体已经承担损失。本文将“可能通过哪些渠道发生”与“已有资料确认发生”分开。若地方通过提高效率和稳定转移支付吸收冲击，家庭生活未必恶化；若主要依靠隐性拖欠和服务降质，财务压力就可能转入企业与家庭账户。

\Chart{fiscal}{2026年1至8月部分财政项目的增速}{数据图；同比，\%。不同预算及支出项目并列展示，不作加总。来源：财政部国库司。\cite{Fiscal2026}}

一般公共预算收入与支出的差额，不等于同额新增债务，更不能直接称作官方赤字。转移支付、债务融资、结转结余和调入资金影响平衡。政府性基金预算也有相应资金来源。一个诚实的财政分析，应先解释预算账户，再讨论可持续性。

% ARGUMENT P127
\phantomsection\label{arg:P127}
预算平衡还涉及不同层级。地方本级收入不足以覆盖支出，并不必然意味着同额债务增加，因为上级转移支付是其正常收入来源之一。全国汇总时又需要避免把中央支出与地方收到的同笔转移重复相加。理解财政账户，首先要知道观察的是本级、合并还是全国口径。

经常收入与一次性资金也应分开。处置资产可以改善当期现金，却可能减少未来收益；调入结余可以跨越短期缺口，但不能无限持续；新增债务则形成未来承诺。可持续性分析需要把这些来源放入跨期框架，并观察支出是否形成有用服务或收入。一个年度的收支差额只是入口，长期判断还要依赖人口、税基、维护和偿债安排。

全国数字还会掩盖地区差异。人口流入、产业与土地需求较强的城市，可能仍有融资和建设空间；需求较弱地区的问题可能是存量设施维护和债务偿还。不能因为全国需要住房调整，就给所有地方开同一剂扩张或收缩处方。政策的空间尺度应匹配问题。

% ARGUMENT P128
\phantomsection\label{arg:P128}
空间差异意味着相同政策可能产生不同结果。在就业增长、公共交通完善的地区，适度增加住房与服务可能降低迁移成本；在需求持续减弱、设施闲置较多的地区，继续新建可能加重维护和债务。判断应以实际人口、岗位和使用需求为依据，而非只追求建设数量。

同一城市内部也存在差异。中心区住房供给紧张与外围库存较高可以并存，因为交通、学校和就业距离使它们不是完全替代品。把全市库存除以人口，无法直接判断每类住房是否过剩。政策需要匹配具体居住需求，同时考虑改造和通勤成本。对低收入劳动者而言，距离岗位更近且稳定可负担的住房，可能比名义面积更大的远郊住房更有价值。

\section{怎样评价债务置换和住房收储}
债务置换可能降低利息、延长期限、减少短期再融资风险。若新债仍对应相同低回报项目，它改善的是偿付安排而不一定是最终净值。若争取的时间用于完成有用项目、处置低效资产和改革财政关系，就可能改善可回收价值。两种情况应分别追踪。

% ARGUMENT P129
\phantomsection\label{arg:P129}
置换争取的时间是否被有效使用，是政策成败的关键。若只是延后付款而没有处理项目收益和支出责任，未来缺口可能继续累积；若同时改善回款、完成接近投用的项目、停止无效扩张，就可能提高可回收价值。债务安排与经营改革应分别记录，才能知道效果来自哪一部分。

置换还会改变债权人之间的风险。新债可能具有不同担保、清偿顺序和持有人，原有银行或供应商获得更稳定资产，而公共部门承担更多责任。这个转移不必然不合理，但需要公开。评价应比较转移前后的总成本、支付可靠性和激励，而不是把法律形式更规范直接等同于经济风险已经消除。

一个假设地方有100单位债务，原年利率5\%，置换后为3\%，当年利息少2单位。在没有其他变化时，2是当期现金流改善，不是100单位债务被消灭，也不能将未来每年的节省不折现相加称为无条件财富。若为置换提供新的担保，担保风险也须记录。

% ARGUMENT P130
\phantomsection\label{arg:P130}
利息节省可以支持其他用途，但它并非没有承担者。债权人若自愿接受较低收益，可能因为违约风险降低、流动性改善或市场利率变化；若依靠公共支持降低融资成本，支持本身也有资源含义。需要说明收益下降由谁吸收，以及是否造成新的不透明补贴。

将未来节省折为现值，还要选择折现率和期限，并考虑本金、担保与后续新增债务。若置换以后继续扩张同类低回报项目，当期利息减少可能很快被新增承诺抵消。这个算例的结论因此很有限：期限与价格调整可以改善现金流，但最终负担取决于整个资产负债关系。对家庭和基层服务而言，最重要的是节省是否真正转化为稳定支付，而不只是报表上的财务指标改善。

收购存量住房用于保障居住，可能同时帮助交付、改善闲置利用和提供公共服务。但收购价格、区位、租赁需求与维护成本决定结果。在人口持续流出的地点，即使购买价格低，也可能形成长期低利用率公共资产。评价应围绕受益家庭和持续服务，而不是消化库存数字本身。

% ARGUMENT P131
\phantomsection\label{arg:P131}
住房收储的收益来自真实使用。闲置住房转成可负担租赁，可以减少新建资源并帮助需要居住的人；但位置、户型、质量和配套不合适时，低价取得也未必划算。改造、运营、维修和空置成本都应进入长期预算。

受益群体的确定也会改变效果。若主要帮助有稳定收入、原本就容易获得住房的人，政策对最脆弱家庭的改善可能有限；若资格过严，流动劳动者可能无法使用；若租金过低而没有持续补贴，运营质量又可能下降。需要在可负担性、覆盖和财政能力之间形成清楚规则。库存消化是资产市场结果，基本居住改善是社会服务结果，两者可以相互支持，但不能用前者的数量替代后者的效果。

政策应公开项目层面的价格依据、质量、入住率、运营收入与补贴需求。若以明显高于合理价值的价格收购，损失只是从原持有者移至公共部门；若以过低价格强行转移，也可能损害合同信用并阻碍正常融资。透明的评估和有序损失分担，比分别宣布“救”或“不救”更具体。

% ARGUMENT P132
\phantomsection\label{arg:P132}
项目评估应公开可比较的依据，包括近期交易、房屋质量、实际租赁需求、改造成本和长期维护。单一评估值无法消除不确定性，可以采用区间和敏感性分析，观察入住率或租金偏离预期时需要多少补贴。这样才能让公共决策看到风险，而不是只看到一个确定采购价格。

损失分担还涉及机会成本。用财政购买住房，就意味着相同资源不能同时用于其他服务；让债权人承担损失，可能影响其资本与后续融资；要求原股东追加投入，则取决于可追索资产和法律责任。政策比较应把这些后果放在一起。透明程序的作用，是让分担可解释、可核查，并减少拥有谈判优势的主体把损失转给不在谈判桌上的家庭与纳税人。

\section{城乡成本转移的启发与边界}
温铁军的历史研究把城乡关系放在工业化成本与危机吸收中，强调不能只观察城市部门的稳定而忽略农村承担的调整。这为本书的风险追踪提供了重要启发。使用这套解释时，必须把返乡、家庭支持、土地保障与公共财政分别记录，不能把农村描绘成容量无限、没有代价的缓冲空间。\cite{Wen2021}

% ARGUMENT P133
\phantomsection\label{arg:P133}
温铁军的成本转移视角提醒我们，城市部门的稳定可能依赖城乡之间的资源与家庭联系。外出劳动者在工作中断后返回农村，短期减少城市就业和居住压力，农村家庭却可能增加生活与照护负担。这样的调整不会全部表现为正式财政支出，却会改变社会再生产条件。

这一解释最有价值的部分，是要求沿着成本去向继续追问。它不应被扩大为每次危机都由农村承担的单一规律。不同历史时期的土地、人口、就业与财政条件不同，农村也可能获得公共投入和转移。本文将其与部门账户结合，分别观察收入、时间、资产和公共服务，既识别缓冲作用，也检查缓冲是否以家庭能力被透支为代价。

返乡可能提供居住与部分生活保障，也可能意味着货币收入下降、照护责任增加和技能闲置。留在农村的家庭成员承受的时间与资源成本，不能因为没有形成银行不良贷款就当作不存在。研究“软着陆”需要说明是哪个部门较平稳，以及另一个部门承担了什么。

% ARGUMENT P134
\phantomsection\label{arg:P134}
返乡的经济意义取决于家庭仍然拥有何种资源。可使用住房、土地经营、亲属支持和本地岗位，会影响回到乡村后的生活；若这些条件已经减弱，返乡未必能够提供足够缓冲。土地提供部分食物与居住基础，也不能自动替代现金收入、医疗和教育服务。

家庭内部的分担尤其重要。留守老人可能承担更多育儿，年轻成员可能减少技能积累，女性可能退出有酬劳动。城市企业减少的成本，并不一定消失，而可能变成这些时间与机会损失。研究需要追踪家庭成员，而不只记录外出者是否重新就业。所谓平稳调整，应说明哪些人的生活保持连续，哪些人支付了不可见代价，并比较有没有公共制度能够减少这种被动转移。

这也构成对历史解释的反证要求。如果某次城市调整主要由中央财政、金融资本或城市资产持有者承担，而农村收入和照护负担没有相应变化，就不能仍强行称作向乡村转移。随着人口和家庭结构改变，过去能发挥作用的安排也可能逐渐减弱。

% ARGUMENT P135
\phantomsection\label{arg:P135}
反证可以通过不同冲击比较进行。某次调整若主要由金融资本减值吸收，农村家庭并未增加负担，就说明城乡转移机制在该情境下较弱；若公共保障扩大以后，同样就业冲击对返乡和家庭照护的影响减小，也说明制度能够改变路径。这样的结果不会否定所有历史观察，却会限制理论的普遍范围。

人口老龄化、家庭小型化和长期城市定居，也可能降低过去依赖亲属缓冲的能力。一个家庭中可承担照护的人更少时，同样冲击会产生更集中负担。历史制度的韧性因此不能被当作无限资源。政策需要观察旧缓冲正在怎样变化，并建设可以跨地区、跨职业延续的保障，使风险处理不再过度依赖某些家庭成员的无酬承担。

本章把地租更新为稀缺位置、公共投资、正式产权和融资共同形成的收益关系，把住房调整理解为使用功能、资产重估和风险分担的联合问题。政策目标应是可持续的居住与公共财政，而不是让某一种资产价格永远沿同一方向变化。

% ARGUMENT P136
\phantomsection\label{arg:P136}
地租、住房与财政之所以必须联合分析，是因为同一块土地同时承载生产、居住、公共服务和资产收益。不同用途之间存在真实取舍，不能只通过提高资产价格解决。价格上涨可能支持融资，却提高新进入者成本；价格下降可能降低门槛，却损害既有高负债家庭。制度需要在这些冲突中明确目标。

本章采用的评价顺序，是先保障基本居住与必要公共功能，再识别有用投资和真实损失，最后安排资产与债务的有序调整。这个顺序是一项公开的价值选择，其实施仍受财政和法律条件约束。它不承诺消除全部损失，而要求损失可解释、责任与控制相联系，并尽量避免让缺少选择的人承担不可逆生活损害。家庭再生产由此成为理解住房与财政调整不可缺少的一环。

\chapter{家庭与社会再生产}
\section{生产并不从工厂门口开始}
劳动者每天能够工作，依赖睡眠、饮食、健康、交通和有人照顾需要照顾的家庭成员。下一代劳动能力的形成还依赖教育和成长环境。企业可以把这些条件视作已经存在，但全社会不能这样假定。若维持劳动能力的成本长期由家庭无酬吸收，生产账户上的低成本可能对应生活账户上的高压力。

% ARGUMENT P137
\phantomsection\label{arg:P137}
劳动能力不是每天自动恢复的存量。长时间工作可能减少休息，疾病可能使技能暂时无法使用，照护中断可能迫使家庭成员离开岗位。企业短期看不到这些变化的全部成本，但它们会通过缺勤、人员流动、健康和下一代能力反馈到生产。社会再生产分析把这些延迟结果纳入同一过程。

这里也要看到家庭的主动安排。亲属可以共享居住、协调时间和相互支持，这些合作具有真实价值；问题在于合作是否建立在可选择条件上，还是某些成员缺少退出途径。若一名家庭成员被迫长期承担照护，其收入、技能和未来保障可能下降，家庭总收入看似稳定却掩盖内部机会损失。以家庭为单位分析时，仍需要继续观察成员之间的权利与负担。

Fraser将照护与资本积累的关系解释为社会再生产矛盾，强调生产依赖却可能侵蚀其生活基础。本文借用这一研究问题，但不把其中关于资本主义整体的判断直接当作中国事实；中国的财政、公共产权和家庭制度需要单独检验。\cite{Fraser2016}

% ARGUMENT P138
\phantomsection\label{arg:P138}
社会再生产矛盾可以表述为一个条件性反馈：生产组织要求较低成本和较高劳动投入，而维持健康、照护和成长需要时间与资源；若前者持续挤压后者，短期成本优势可能转化为长期能力损耗。这个机制需要观察工时、服务、健康和家庭安排，不能仅凭某个宏观增长率判定。

中国情境中，公共服务、单位制度、市场购买和亲属支持共同承担再生产任务。它们的组合会随地区、职业和生命阶段变化，不能用一种纯粹市场化家庭模型概括。公共供给可能缓解家庭压力，也可能因资格与质量差异形成新的分层。本文因此把社会再生产理论作为提出问题的工具，再用具体制度与资料检验其成立范围，而不把某一理论传统的总体结论直接移植。

社会再生产分析把这部分条件放回政治经济学。它与经典资本再生产问题相连，又扩大了考察对象：不仅问机器、原料和商品能否进入下一轮生产，也问劳动者及其家庭能否持续生活。本文不因此把亲情全部商品化，而是研究制度怎样分配必要劳动和照护责任。

% ARGUMENT P139
\phantomsection\label{arg:P139}
经典资本再生产关注生产资料与消费资料怎样支持下一轮生产，家庭再生产则进一步追问劳动者怎样恢复、成长与延续。两者相互联系，却不能用同一个产值指标全部计量。机器维护有相对明确的工程标准，人的生活还包含尊严、关系和自主选择，不能只按未来劳动供给评价。

因此，改善照护的价值不只在于让更多成年人进入有酬工作。它也可能使失能者得到更好的生活，使照护者获得休息，使儿童获得稳定关系。即使这些改善没有立即提高GDP，也具有独立意义。本书把生产效率与基本生活并列为价值标准，正是为了避免把人仅仅视为生产的投入。政治经济学需要解释生产怎样服务生活，而不是让生活的全部价值取决于其对产出的贡献。

企业、家庭、市场服务和公共部门可以承担同一类任务的不同部分。托育服务增加，可能使家庭成员恢复有报酬工作；但如果服务时间与工作时间不匹配，名义覆盖就未必转化为实际帮助。社会再生产的评价对象，是生活条件的可用性，而不只是设施数量。

% ARGUMENT P140
\phantomsection\label{arg:P140}
任务在家庭、企业与公共部门之间转移时，成本结构会变化。企业提供食堂可能减少员工准备餐食的时间，公共托育可能使家庭共享专业设施，市场服务可能提高选择但增加现金支付。哪一种安排更有效，取决于规模、质量监督、地点与使用时间，不能仅按公共或私人标签判断。

服务可用性还包含可靠性。偶尔能够使用的照护，未必足以支持稳定就业；频繁停供会迫使家庭保留备用安排，实际减少的负担就较小。评估应记录服务中断、替代成本和家庭是否能够据此安排工作，而不是只统计开设机构与登记名额。对缺少亲属支援的流动家庭，可靠供给的价值可能尤其高；对愿意亲属照护的家庭，则需要保留选择与必要支持。

\section{人口变化是约束，不是单一命运}
2025年末，中国大陆60岁及以上人口占23.0\%，65岁及以上占15.9\%；全年出生人口792万人、死亡人口1131万人。年龄结构改变需要养老、医疗、照护和职业安排作出响应，但人口数量不能独自决定经济增长或个人福祉。生产率、健康寿命、迁移和服务组织同样重要。\cite{NBS2025}

% ARGUMENT P141
\phantomsection\label{arg:P141}
人口结构影响的是需求与供给的组合。老年人口增加会提高某些医疗和照护需求，劳动年龄人口变化会影响潜在劳动供给，但实际就业还取决于健康、技能、参与率和岗位组织。同样人数的劳动者，使用不同技术和制度，能够提供的产出与生活质量也不同。

人口变化具有较强惯性，却不等于结果已经确定。教育和健康改善需要时间，住房和公共设施也具有长期寿命，政策应考虑不同阶段的需求转换。例如某地儿童减少而老人照护需求增加，现有设施能否改造、人员能否转训，就成为具体问题。单纯扩大建设总量可能错过结构变化。人口分析的价值，在于帮助提前安排资源，而不是将复杂经济问题归结为一个出生率或抚养比。

把老人一概视为负担，会忽略其劳动、照护、知识和消费作用；把延长劳动参与视为没有代价，也会忽略职业健康、身体差异和家庭照护。退休安排涉及群体之间的异质性，办公室职业与高强度体力职业难以用同一条健康假设概括。

% ARGUMENT P142
\phantomsection\label{arg:P142}
老年人的经济位置差异很大。拥有稳定养老金和住房的人，与主要依靠亲属支持、健康状况较差的人，面对同样价格或医疗冲击时承受能力不同。年龄是生命阶段，不自动构成同一阶级。老人还可能承担育儿、家务与社区活动，这些投入会改变其他家庭成员的就业条件。

延长劳动参与需要区分自愿选择与被迫维持生活。部分专业人员愿意继续工作，部分体力劳动者却可能因健康难以维持原强度。若只调整年龄规则而不改变岗位、安全和照护安排，成本可能集中到特定职业。政策评价应比较健康、收入和可选择岗位，并允许不同劳动经历进入制度设计。人口老龄化带来的财政问题真实存在，但解决方式需要明确代际与职业之间怎样分担。

出生人口减少对教育设施、住房和消费的影响也存在时间差。当前儿童数量变化不能立即等同于当前劳动力减少。改善育儿条件可以提高家庭选择空间，但生育决定还受个人偏好、伴侣关系和长期生活预期影响。不能把提高出生数量当作评价每个家庭或每项福利的唯一标准。

% ARGUMENT P143
\phantomsection\label{arg:P143}
人口变化对不同市场的影响有滞后。出生减少先影响产科、托育和幼儿教育，之后逐步影响学校与劳动市场；住房需求则还受家庭分户、迁移、存量质量和居住面积影响。直接把出生变化推成当年住房需求同比变化，会忽略这些中间条件。

育儿政策的目标也需要清楚。降低养育压力可以改善儿童与父母生活，无论最终出生数量怎样变化；若只以新增出生人数评价，可能低估对已有家庭的收益。生育决定涉及身体、关系和长期承诺，家庭有权作出不同选择。公共政策更适合减少因服务不足和不公平负担造成的被迫选择，而不是把个人决定视为需要完成的生产指标。

\section{无酬劳动需要正确的分母}
2024年全国时间利用调查是本书采用的较新公开时间资料，但它不是2026年实时调查。调查公布的无酬劳动参与者平均时间，男性为每天112分钟，女性为209分钟。这里分母是参与这类活动的人；不能说全国所有男性或女性平均恰好投入这些时间，也不能把差额直接解释为歧视的因果估计。\cite{TimeUse2024}

% ARGUMENT P144
\phantomsection\label{arg:P144}
参与者均值与全体均值之间，还隔着参与率。若某类活动只有一部分人参与，全体平均时间等于参与率与参与者平均时间的组合；两个群体即使参与者时长相近，也可能因参与率不同而具有不同总负担。要理解性别或年龄差异，需要同时观察这两个维度。

时间日记还可能存在同时活动。例如一边做饭一边照看孩子，调查怎样记录主要活动和伴随活动，会影响总量。直接把不同任务时间相加可能超过实际可用时间，完全忽略伴随照护又可能低估约束。本书保留调查所用口径，不自行把已公布数字重构为全国总照护小时。图表用于展示公开观察，机制分析则说明需要进一步取得哪些信息。

\Chart{care}{2024年无酬劳动参与者的平均时间}{数据图；分钟／日，按性别。参与者均值，非全部人口均值；活动分类和调查日规则见原报告。\cite{TimeUse2024}}

这项差异足以提示研究任务：照护怎样影响就业时长、岗位选择、收入与养老金积累？但具体解释仍需家庭结构、就业状态、年龄和城乡资料。只比较男女均值，无法排除不同生活阶段的组成差异。即便如此，组成差异也可能受到制度影响，所以“可以被控制变量解释”并不自动等于“无需政策关注”。

% ARGUMENT P145
\phantomsection\label{arg:P145}
照护对职业的影响可能逐步累积。减少工作时长先影响当期收入，随后可能影响培训、晋升和社保缴费，长期又影响养老金与职业转换。一次短期退出如果发生在技能快速更新的行业，回归成本可能较高。家庭内部看似临时的安排，因此可能形成持续的社会位置差异。

这些结果不能全部归于个人偏好。偏好本身会在服务可得性、工资差距和家庭协商中形成；但也不能否认家庭成员对照护的不同意愿。研究应比较具有相似照护需要、但服务和制度条件不同的家庭，观察选择是否变化。若公共服务改善以后更多人能够实现原先希望的安排，说明约束确实发挥了作用；若选择不变，也需要尊重其可能具有的非经济理由。

无酬劳动的可见化有两种常见误区。一种是把所有家务按高薪职业机会成本估值，得到很大的总数后直接加到GDP；另一种是因为估值困难，就将其彻底排除。更好的办法是先建立时间与服务的卫星账户，再对替代成本和机会成本分别给出范围。

% ARGUMENT P146
\phantomsection\label{arg:P146}
卫星账户的作用，是在保持核心国民核算一致的同时，展示被其边界排除的重要活动。时间、任务、受益对象和替代服务价格可以分别记录，货币估价则作为不同假设下的结果。这样既让无酬劳动可见，也避免把估计值误当作已经发生的市场收入。

估价还应注意分配后果。如果用个人市场工资估算同样照护任务，高收入者的家务会被赋予更高货币值，这反映的是机会成本而非受照护者得到的服务质量。若用统一替代工资，则又可能忽略专业程度和地区差异。两种方案都可以用于特定问题，但必须解释用途。政策不应只根据一个全国巨大估值作判断，而应识别哪些家庭负担最重、哪些服务能够以合理成本改善生活。

\section{户籍、迁移与服务可携带性}
人口迁移会提高劳动与岗位的匹配，也会产生家庭分离、重复住房、跨地就医和教育安排成本。是否能够稳定定居，不仅取决于就业合同，还取决于公共服务资格与可负担住房。城乡关系分析因此应从“人是否移动”推进到“权利和服务是否随人移动”。

% ARGUMENT P147
\phantomsection\label{arg:P147}
迁移的收益通常来自岗位与能力更好匹配，但其成本不只是一张车票。寻找住房、重新建立照护网络、办理服务资格和适应工作，都需要资源。家庭成员不能同时迁移时，还会产生两地生活与情感支持的安排。工资差额只有扣除这些条件以后，才能接近家庭实际获得。

服务可携带性也有不同层次。权益记录能够转移，是账户层面的连续；异地可以实际就医或入学，是供给层面的连续；支付与结算及时，则是现金层面的连续。某一层改善不代表其他层已经解决。对流动劳动者而言，稳定预期尤其重要：如果服务资格会随短期就业变化中断，家庭就可能不敢长期定居。这种不确定性会反过来影响劳动流动与消费安排。

中国各地服务政策与实施条件存在差别，不能用一条全国口号代替具体资格规则。本书不列举未经逐地核查的入学和医保待遇，而提出可比较维度：申请门槛、等待时间、个人支付、跨地结算、可延续性以及家庭成员覆盖。名义参保与实际可获得服务，可能因这些条件产生距离。

% ARGUMENT P148
\phantomsection\label{arg:P148}
比较地方制度时，可以从一个具体家庭的完整过程出发：需要什么证明、在哪里申请、等待多久、是否需要先垫资、换工作以后权益是否继续。每个环节都可能形成实际门槛。信息充分、时间宽裕的家庭更容易完成程序，工时长或频繁迁移的家庭则可能承担更高办理成本。

这使行政程序本身具有分配效果。简化手续和跨地互认，可能在不改变名义待遇的情况下提高实际获得；但若供给能力不足，仅取消证明也可能增加排队。政策需要把资格、办理与资源扩容配合起来。本文不以一条全国性原则推断所有家庭已经平等获得服务，而将可用性作为需要持续核查的结果。

人口流入地需要增加服务供给，流出地却可能仍有老人照护和设施维护支出。如果财政资金不能适当随实际服务人口调整，地方就可能把新增常住人口主要视为成本。改善跨地区财政关系，既是公平问题，也是提高劳动配置效率的问题。

% ARGUMENT P149
\phantomsection\label{arg:P149}
流入地与流出地之间的成本，不能只按户籍人口或当期就业人数单独分配。儿童和老人可能留在原籍，劳动者在城市使用交通、医疗和居住设施，两地都承担部分公共服务。若资金只跟随其中一个指标，另一部分需求就可能缺少支持。

更合理的财政安排需要考虑实际服务人口、年龄结构、地区成本和既有设施。流入地增加服务可以促进就业匹配，流出地维持基本保障则减少家庭分离带来的损失。跨地区协调的收益可能超出单个地方，所以需要更大范围的分担。与此同时，转移支付也应避免奖励低效供给，资金使用和服务结果需要公开。公平与效率在这里通过人口流动的具体条件连接起来。

\section{家庭风险如何变成需求问题}
家庭储蓄既可以为购房和教育准备，也可以应对失业、疾病和养老的不确定性。储蓄行为本身不能证明保障不足，因为收入增长、资产选择和人口结构也会影响储蓄。要识别预防性动机，需要观察相似家庭在保障改善或风险冲击前后的行为。

% ARGUMENT P150
\phantomsection\label{arg:P150}
家庭储蓄的动机可以同时存在。准备首付款、养老、教育、创业和应对疾病，可能使用同一账户，单看存款余额无法分解。收入上升时，储蓄增加也可能只是基本消费满足后的正常选择；资产价格变化则会改变家庭希望持有的流动资金。

Chamon、Liu和Prasad利用较早时期中国家庭资料研究收入不确定性与储蓄，为预防性机制提供了历史证据。其样本与模型对应特定年代，不能据此认定2026年的全部储蓄变化都由保障不足造成。本文采用它来说明风险与行为之间存在可研究的联系，同时把住房、人口和收入预期作为竞争解释保留。真正的当期判断需要更新家庭资料，而非把历史系数直接移用。\cite{Chamon2010}

不过，机制可以清楚表述：如果未来必要支出的尾部风险很大、保险不充分，家庭可能增加安全资产并减少当期非必要消费。企业面对较弱终端需求，可能减少招聘，进一步强化家庭谨慎。并不是每个家庭都会进入这个循环，也不意味着所有储蓄都应被刺激花掉；关键是降低无法有效自保的基本风险。

% ARGUMENT P151
\phantomsection\label{arg:P151}
尾部风险指发生概率可能较低、但一旦发生就严重影响生活的支出或收入损失。家庭若无法通过保险或公共制度有效分担，就可能保留较多流动资金，即使平均收入预期并不差。对于资源较少的家庭，困难还可能表现为根本无法储蓄，只能在冲击后压缩必要消费或借款。

因此，预防性储蓄与脆弱性不能简单按存款高低排序。储蓄较多者可能在主动自保，储蓄很少者可能更缺少自保能力。政策若只希望把存款变成消费，会忽略家庭为何需要缓冲。更有效的问题是哪些基本风险可以被可靠分担，服务是否能在需要时兑现，以及减少不确定性以后家庭是否获得更大的生活选择。消费变化只是其中一个结果。

候选命题是，生产扩张只有在工资、服务与保障使家庭再生产能力同步改善时，才更容易形成稳定内需。这里的同步并非要求每年同幅增长，而是避免长期把维护劳动力的成本推到收入和时间都不足的家庭。竞争性解释包括收入预期、住房调整与消费偏好，实证必须允许它们同时存在。

% ARGUMENT P152
\phantomsection\label{arg:P152}
生产扩张与家庭再生产的联系，可以沿着收入和时间两条路径理解。工资与公共服务提高，可能支持更多稳定需求；工作时间过长或照护缺口扩大，则可能减少家庭消费时间与劳动参与。两条路径会相互影响，例如照护负担使一名家庭成员退出工作，既减少收入，也增加其他成员的压力。

这里的命题需要限定。“同步改善”不是说每一项投资都必须当年增加消费，也不是说家庭保障单独决定增长。追赶性投资可以先形成未来能力，外需也能在一定时期支持生产。需要警惕的是长期把家庭时间、健康和风险承受当作没有约束的投入。若这些条件持续恶化，生产与需求的连接可能变弱。实际强度仍需通过家庭和企业资料检验。

\section{公共服务与现金支持怎样配合}
现金支持让家庭按自身需要选择，但不能凭空创造附近的托育位、医护人员或无障碍设施。公共服务直接供给可以扩大可用资源，却可能存在选址不当、质量差和排队问题。两者应根据约束配合，而不是被当作互相排斥的意识形态选择。

% ARGUMENT P153
\phantomsection\label{arg:P153}
现金与服务的效果，取决于约束所在。供给充足而价格过高时，现金支持能够扩大选择；专业人员和设施不足时，单纯增加购买力可能更多推高价格。直接供给可以扩大能力，但如果忽略家庭需要，也可能形成低利用率。两者不是固定的优劣排序，而是针对不同问题的工具。

资金来源同样影响净效果。补贴由提高同一群体税费筹集，实际净收入改善可能有限；服务扩张若挤占其他必要项目，也存在机会成本。评价应比较家庭最终获得、供给质量和财政持续性，而不是把发放金额或建设面积直接当作福利增加。对不同阶层，服务的可达性与现金的灵活性具有不同价值，政策需要保留一定选择空间。

假设某地母亲退出就业的主要原因是缺少晚间托育，统一发放一小笔补贴未必改变其约束；如果当地已有合格服务但费用过高，补贴就可能发挥作用。政策需要先识别时间、价格还是质量瓶颈，再决定财政资源投向。受益者也不只是家长，儿童和照护劳动者的安全与待遇都应纳入评价。

% ARGUMENT P154
\phantomsection\label{arg:P154}
托育时间错配说明，需求不仅是数量问题。早晚班劳动者、轮班服务人员和固定办公时间的白领，所需时段可能不同。统一营业时间容易服务其中一部分家庭，而使其他家庭继续依靠亲属。若政策只统计名额，便无法看到这种差异。

设计服务时还需考虑工作人员。延长营业时间需要合理排班、报酬与安全，不能把家长的时间压力简单转给低薪照护劳动者。儿童的休息和稳定关系也应进入质量标准。一个方案可能同时改善家长就业和儿童照护，也可能只提高劳动供给而降低服务质量。完整评价需要把家长、儿童、从业者与公共预算放在同一机制中，而不是只选择其中一方的便利。

改善照护供给还可能创造就业，但不能把就业数量直接当作净新增财富而忽略原有家务替代。家庭无酬活动转为市场服务，会提高被GDP记录的活动量，其中一部分只是核算边界改变。真实改善应看服务质量、照护者选择空间和劳动生产率。

% ARGUMENT P155
\phantomsection\label{arg:P155}
家庭活动转为有酬服务后，统计边界与真实效率可能同时变化。若专业机构通过共享设施和分工提高质量，确实可能增加可用服务；若只是把同样任务从家庭转到市场，GDP记录增加的一部分来自交易形式改变。二者需要分别说明。

释放出来的家庭时间也有多种用途。有人增加有酬工作，有人学习、休息或照顾其他成员，这些选择不应全部按新增劳动供给评价。对长期承担照护的人，能够获得连续休息本身就是生活改善。政策研究可以记录时间去向和自主选择程度，不必强行把所有收益折成下一年度产值。社会再生产分析由此扩大了评价维度，同时保持核算上的清楚边界。

\section{从隐形补贴走向公开分担}
家庭以无酬劳动和风险自担支持生产，可以被称为对生产体系的隐形支持；但这个说法必须落实到任务和机会成本，不能暗示每个家庭都向某一家企业支付了可以精确追索的款项。企业税费、劳动报酬、公共服务和家庭活动共同构成一个复杂体系。

% ARGUMENT P156
\phantomsection\label{arg:P156}
“隐形支持”是一个分析比喻，用来指出生产依赖未充分显现在企业账上的条件。它不能直接推出某个家庭对某家公司拥有确定债权，因为劳动报酬、税费、公共服务和家庭活动之间存在复杂对应。若要提出具体补偿，需要明确责任、受益与可执行机制。

更可操作的制度路径，是将普遍需要的保障通过公共或保险安排分担，将雇主能够控制的工作风险纳入其责任，将家庭自主选择保留下来。这样可以减少某些成本长期落在缺少议价能力的成员身上，同时避免把所有社会责任都压给单一企业。责任分配应依据控制能力、受益范围和风险共担效率，而不是仅凭某项活动发生在家庭还是工作场所。

本文主张优先公开照护时间、基本保障缺口和跨地服务障碍，再讨论怎样分担。家庭仍可以选择亲属照护，公共政策的目标是减少被迫承担和不可逆退出，而不是把一种生活方式规定为唯一正确答案。

% ARGUMENT P157
\phantomsection\label{arg:P157}
被迫承担与自主选择之间，往往取决于有没有可行替代。家庭成员愿意照护亲人，并不表示他愿意完全失去收入、休息和未来保障。提供替代服务、照护支持与可恢复的职业安排，可以使选择更接近真实意愿。政策目标是扩大可行集合，而非替家庭决定唯一生活方式。

公开负担还可以改善家庭内部协商。若照护时间和职业损失长期不可见，收入较高者可能被误认为贡献更多，承担无酬劳动者则缺少表达依据。记录并不自动解决权力差异，但能够使讨论有具体对象。研究应尊重家庭隐私，并避免把个人道德评价替代制度分析；重点在于哪些公共条件可以减少长期不对称，使成员在风险发生时仍保有基本自主。

本章对再生产理论的更新，是把社会持续生产的条件从设备和劳动力数量，扩展为时间、健康、保障与代际能力。若这些条件被透支，某些企业的成本优势可能不可持续；若得到改善，收益也可能表现为更健康的生活和更自由的选择，而不只是下一年的GDP增速。

% ARGUMENT P158
\phantomsection\label{arg:P158}
再生产能力具有累积性。健康、教育和技能在较长时期形成，持续透支以后恢复也需要时间。短期用工成本较低，可能伴随后续医疗、照护和技能更新的更高支出；这些成本若由公共部门和家庭承担，企业账上的优势就不能直接等同于社会效率。

与此同时，改善生活条件也可能产生生产反馈：更可靠的照护支持就业，健康改善减少中断，服务可携带性促进匹配。但这些反馈不是评价福利的唯一理由。基本保障首先回应人的生活需要。将生产与生活双向连接，可以解释制度怎样形成长期能力，也可以避免把所有家庭政策都只当作刺激消费或增加劳动力的手段。第八章将在这一基础上，系统分析不同社会位置的资源与约束。


\chapter{收入、财富与社会位置}
\section{一个平均数装不下分配}
收入是某一时期的流量，财富是某一时点的资产扣除负债后的存量。高收入者不一定已有大量资产，拥有住房的退休者也可能现金收入较低。比较生活条件还需家庭规模、公共服务与必要支出。把任何一个数字直接称为阶级地位，都会丢失重要信息。

% ARGUMENT P159
\phantomsection\label{arg:P159}
把收入与财富分开，是为了看清生活怎样延续。收入支付当期食物、交通、住房和照护；资产可能提供居住服务、利息或经营条件，也可能在收入中断时被出售。负债则把未来收入的一部分预先约定给债权人。同样拥有一套住房，无贷款的家庭和必须按月还本付息的家庭，在失业时面对的时间压力不同。前者可能缩减其他支出后等待，后者还须维持合同支付。这一差异来自资产用途和支付期限，不是来自对谁更勤劳的判断。

分析单位也影响结论。把每个人视为完全独立，会漏掉夫妻共用住房、父母照护和家庭共同债务；把家庭看作没有内部差异的整体，又会掩盖谁拥有产权、谁可以动用存款以及谁承担无酬劳动。一户家庭有较高净资产，并不保证其中每位成员都能够离开不利工作或家庭关系。因此，本章在家庭预算之外保留个人的支配权限、社会保险权益和时间约束。

这里使用“社会位置”，是要描述一种能够持续一段时间的资源关系。收入下降一次不一定意味着阶层位置改变，但长期丧失专业资格、稳定客户、住房保障或经营控制权，可能改变未来取得收入的方式。相反，一次奖金未必带来位置跃升，而获得稳定服务、减少高息债务或积累可以转换行业的技能，可能扩大长期选择。位置既不能由某一天的资产价格确定，也不能由个人自我认同单独确定。

后文的比较以生产、控制、收益和风险为主轴，再记录保障与家庭条件。它不是用六个变量机械计算一个总分：有公共配置权而没有私人股权，与有私人股权而没有公共权限，是两种不同关系；把两者加总成高低排序，会抹掉责任来源。解释社会分层，需要保留这种结构差异。

2026年上半年居民收入平均数与中位数之比约为1.207，即平均数比中位数高约20.7\%。这是由公开值计算的描述量，不是基尼系数，也不能据此重建全部收入分布。两个社会可以有相同平均数和中位数，却在最顶部或最底部存在很大差别。\cite{Income2026}

% ARGUMENT P160
\phantomsection\label{arg:P160}
平均数与中位数的差距，首先说明把全国收入总额平均分配得到的数，与处在排序中间位置的居民收入不同。它不能告诉我们差距集中在哪一端。假设上端少数收入增长，中位数可以不变而平均数上升；假设底端居民收入改善，平均数也可以上升，而中位数仍然暂时不变。因此，单看这个比值无法判断是哪一种变化。

全国居民人均数也不是一份典型工资单。家庭中的儿童、退休者和暂未就业者都进入人口范围，劳动收入、经营收入和转移收入按统计规则共同进入家庭可支配收入。将22981元除以六个月，可以得到相同口径的人均月度平均量，却不能把结果称作全国在职者平均月薪。分母转换需要就业人数、家庭成员结构和收入来源等额外资料。\cite{Income2026}

比较城乡居民时还要处理生活条件。城市现金收入较高，不表示扣除住房、交通和照护以后所有城市家庭都更宽裕；农村现金支出较低，也不表示医疗可及性、劳动机会和照护质量自动相同。可以同时展示收入、必要支出、实物自给和公共服务，但不能为了得到一个方便的排序，任意给未观察的服务赋值。对于工资相近的流动家庭与本地家庭，更直接的问题是同一项基本需要分别要付出多少现金、时间和迁移成本。

\Chart{income}{2026年上半年居民收入的不同统计量}{数据图；元／人，半年。全国平均数、中位数与城乡平均数分别描述不同分组，不代表四类互不重叠的人。\cite{Income2026}}

工资、经营、财产与转移收入是收入来源分类，不是四类固定身份。领取工资的人可以同时有利息和租金，个体经营者也可以领取养老金。2026年上半年全国人均可支配收入中，工资性收入占57.9\%，经营净收入占15.8\%，财产净收入占8.0\%，转移净收入占18.3\%。这不能直接转化为劳动与资本在国民收入中的份额，因为两个统计体系的覆盖与定义不同。\cite{Income2026}

% ARGUMENT P161
\phantomsection\label{arg:P161}
收入来源能够帮助识别依赖关系，但要从统计分类再向合同和活动深入一步。经营净收入中可能包含家庭成员没有单独领取的劳动报酬，也可能包含设备与店面资本的回报；财产净收入可能来自小额利息，也可能来自大规模出租资产。两者同属一种来源，不意味着其规模、稳定性与控制条件相同。给家庭贴上“经营者”标签以前，应先问经营是否依赖本人每天劳动，是否雇佣他人，以及扣除工具折旧与融资成本后还剩多少。

工资性收入也有内部差异。固定薪酬、按件工资、绩效奖金和高比例佣金，使劳动者对订单变化承担不同程度的风险。一名有稳定底薪的服务人员与收入完全依赖接单的劳动者，即使全年所得相同，月度波动也可能不同；若必要支出刚性较高，波动本身就会带来借款或违约成本。只比较年均收入，无法看到这种代价。

本书因此区分两种“依赖”：一个来源占家庭收入的比例，以及失去这个来源后能否找到替代。两家店都从同一客户取得一半收入，其中一家能够转接其他订单，另一家的设备只适合该客户，实际依赖程度并不一样。家庭与企业在这里具有相似结构。可持续的选择空间，来自来源多样性、资产可转换性和有效保障的共同作用，而不是收入来源名称的数量。

\section{财富为什么比收入更难观察}
收入调查通常较难完整覆盖顶端家庭和隐性收入，财富调查还面对企业股权估值、共同持有、负债遗漏和非交易资产定价。没有近期市场成交的企业股权，采用账面价值还是未来收益折现，会得到不同结果。自住房的估值也会随市场和地区变化。

% ARGUMENT P162
\phantomsection\label{arg:P162}
估值方法为什么影响分层判断，可以从一间非上市企业理解。若按账面净资产计量，早年买入的土地和机器未必反映当前可出售价格；若按未来利润折现计量，则必须预测需求、竞争和风险；若企业依赖创办人个人劳动，其全部利润还不能视为可以脱离劳动持续取得的资本收益。把一种估值称作唯一真实财富，会把这些假设隐藏起来。

负债同样需要观察法律关系。公司欠债不一定等额成为股东个人债务，但个人担保、家庭共同借款和抵押安排可能将二者连接。若统计只记录股权价值而遗漏担保风险，会高估家庭缓冲；若将全部公司债务再从股东已经扣除公司负债的股权价值中减去，又会重复扣减。家庭财富表必须先明确公司与家庭的边界，再把跨边界的担保单列为或有责任。

住房估值还涉及可交易性。一套缺乏买方、存在权利限制或需要大修的房屋，即使附近挂牌价格较高，也未必能迅速按该价格出售。家庭为筹措医疗费用而急售，取得的价格可能低于正常等待条件下的估值。因此，净资产与流动性资产需要同时报告。前者帮助理解长期资源，后者帮助理解能否及时支付；高净资产、低流动性的家庭与两者都缺少的家庭，困难相似之处在现金，政策含义却不完全相同。

Piketty、Yang和Zucman关于1978至2015年中国财富和收入分配的研究，尝试结合国民账户、调查和税收等资料，对长期变化作综合估计。其意义在于让宏观积累与私人财富分布进入同一框架；其局限则来自基础资料覆盖、调整假设和历史终点。本文把它用于理解形成过程，不当作2026年的抽样调查。\cite{PYZ2019}

% ARGUMENT P163
\phantomsection\label{arg:P163}
综合估计的必要性来自资料各有所长。国民账户有助于建立总体收入和财富边界，住户调查有助于观察分布，税收和其他行政资料可能补充调查较难覆盖的部分。将它们拼接时，需要解决统计单位、收入定义和缺失部分的分配。研究者作出的调整是测量工作的一部分，应当随结果一同被检验，不能因为最终呈现为一条平滑曲线，就把它当成逐户完整清点。\cite{PYZ2019}

这类长期研究可以支持对产权变化和私人积累的讨论，却不能单独说明某一名家庭成员为什么上升或下降。住房制度改革、企业发展、教育机会和迁移都可能改变分布，宏观共同变化也可能包含不同群体的相反经历。若要区分这些机制，需要把政策时间、地区差异与家庭经历连接，而不只是观察总体财富收入比。

对本书而言，历史估计提供一个可比较的形成过程：资产积累与收入分配必须共同研究。它没有授权我们把截至2015年的分布直接延长到2026年。尤其在住房市场、人口结构和企业估值改变以后，延用旧比例可能产生显著误差。可保留的判断是测量框架和历史发现；关于当前大小与变化速度的判断，则应等待同口径的新资料。

《世界不平等报告2026》的中国国别表将相关数值标为2024年研究估计，并分别按收入和财富排序。由于下文记录的表内一致性问题，本书不使用该页数值绘制精确分配比较图或估计各阶层规模，只在来源审计中保留其原列数值。其收入边界与国家统计局居民可支配收入不同，不能直接拼接。\cite{WID2026}

% ARGUMENT P164
\phantomsection\label{arg:P164}
理解国别估计，须先阅读它所使用的收入边界。国别表图注说明，收入是在养老金和失业保险体系运作之后、个人所得税之前衡量；它与居民可支配收入、企业工资总额均非同一个指标。人口单位、收入单位和排序单位若没有对齐，即使两个数字都称为“前百分之十”，也不能直接比较。

本书把这张表保留在来源审计中，是为了公开核查过程，而不是借它给各阶层安排一个精确人数或财富比例。经济控制关系与分布分位也不是一一对应：高薪专业人员可能进入较高收入分位，却没有决定企业投资的权限；收入暂时偏低的控股股东仍可能控制大量生产资源。分位数回答分布问题，阶层分析回答关系问题，两者需要相互补充。

在存在未解释的表内冲突时，即使某几项数值自身能够满足恒等关系，也不能据此保证整个估计过程没有其他问题。因而扩充版撤下以该表为依据的精确份额比较图。这里的处理不是断言该研究整体无效，而是将证据使用范围收缩到本书能够独立核验的程度；若发布方补充一致的版本与口径说明，可以重新评估采用范围。

\Chart{classflow}{社会位置的形成：资源、制度与实际处境}{机制示意图；箭头表示需要检验的传递关系，不表示经验系数或人口规模。各类位置可以交叉，住房、年龄、迁移和照护责任贯穿其中。}

核查发现，该页底部人口组的均值、份额与组间倍数不能在同一标注口径下同时成立。原PDF视觉与文本核查结果一致。未取得解释或修订以前，本书暂停将该表用于精确分配结论；以下算式仅审计原表，不构成替代估计。\cite{WID2026}

% ARGUMENT P165
\phantomsection\label{arg:P165}
核查可以写成一个简单恒等式：同一统计总体中，某组收入份额应等于该组人口份额乘以该组平均收入，再除以全体平均收入。按该页列出的底部一半人口、组均值1988和全体均值14396计算，份额约为6.9\%，与表列13.7\%不同。按收入份额43.3\%和13.7\%推算前十分之一与底部一半的组均收入比，约为15.8，而页面另列31.4。这里的复算只证明所展示数字无法按同一口径同时成立，不证明哪一个数字应该被改成另一个。

财富列也有类似问题：底部一半的组均财富5447与全体均值86462推算约为3.15\%，而表列份额为6.3\%。人口范围、单位和组别标注若另有特殊定义，发布方需要给出解释；在解释缺位时，本书不自行替数据发布方修表，更不把任何一种“修正值”用作中国财富分布的新估计。详细算式与原页定位见附录。

这种审计比简单标注“数据可能有误”更有信息量。它指出冲突发生在哪个约束，也指出不确定性如何影响使用。若一项政策结论只依赖收入与资产集中可能造成选择差异这一机制，就不需要这些精确份额；若结论依赖最顶端到底占有多少，则应暂缓量化判断。降低结论强度，应落实到具体句子和图表中。

这种做法体现资料规则：国外来源并不因机构声誉而免于核算检查，国内资料也不因发布主体而自动失效。同一基础数据被不同机构整理，增加了方法视角，却不增加独立样本数量。确定性必须与资料质量相称。

% ARGUMENT P166
\phantomsection\label{arg:P166}
来源可靠性不是一个永久授予机构的总分。同一机构的统计表、预测模型、评论文章和政策倡议，证据角色不同。对统计表要核查口径和修订，对预测要核查时间与假设，对政策倡议要核查机制和取舍。把它们统一当作权威结论，会使读者无法判断争议究竟来自事实、模型还是价值标准。

独立验证还须沿数据来源向上追溯。若国内报告、国际机构报告和新闻报道都使用同一份住户调查，这三个文本不能构成三份独立样本；但不同研究对缺失收入采用不同调整，确实可以提供方法敏感性比较。证据的增加可能来自新观察，也可能来自更清楚的误差范围，两者都重要，却不能互相冒充。

当几种资料不一致时，研究者应先尝试解释统计边界，随后再判断是否存在真实冲突。例如税前与税后收入差异可能由税收转移造成，个人与家庭收入差异可能由家庭规模造成。只有完成这些转换以后仍无法相容的部分，才适合进入争议清单。阶层分析尤其需要这种耐心，因为一个看似很小的分母变化，就可能使“多数人”的表述失去依据。

\section{住房与代际积累}
在住房占家庭资产较大部分的情境中，购房时点、所在城市和负债程度会造成代际差异。但这里首先是机制判断，本书没有取得最新家庭微观资产表，不估计2026年全国住房财富占比。早期调查比例不能未经更新就改写成今天的事实。

% ARGUMENT P167
\phantomsection\label{arg:P167}
住房进入阶层形成，至少经过三条不同路径。居住路径改变家庭每月是否需要支付租金，以及能否稳定接近学校和岗位；资产路径改变家庭可出售、可抵押或可传给下一代的资源；负债路径则约定未来现金流。三条路径可能同时存在，也可能方向相反。房价上涨提高已有住房的账面净值，却也可能提高换房、子女购房和迁入同一城市的成本。

若一个地区的就业机会、公共服务和住房供给长期集中，住房还可能成为进入这些机会的门槛。此时问题不是房屋本身神秘地产生社会权力，而是位置、服务资格和抵押规则被绑定在同一资产上。把服务可及性与购买高价资产的能力部分分离，可能降低这一门槛，但效果取决于租住稳定性、服务质量和迁移政策的具体实施。

形成过程不能被简化为一代人整体获益、另一代人整体受损。同龄家庭在地区、单位经历、继承和借款上差异很大；较早进入住房市场者也可能买在就业衰退地区。比较代际差异应同时控制城市与购置条件，并记录没有买房的人。只追踪今天仍然持有高价值住宅的家庭，会遗漏失败投资与被迫出售，造成幸存者偏差。

假设两个家庭工资和储蓄率相同，一个在住房价格较低时获得住宅，另一个在较高价格时借款购买。未来同样工资增长，后者可能有更多现金用于还贷。若两户都有自住需求，前者的账面升值也不必然能够无成本兑现，因为换房成本可能同步上涨。资产效应因居住安排而异。

% ARGUMENT P168
\phantomsection\label{arg:P168}
把前面的假设具体化：设两户每月可用于安排的劳动收入相同，必要生活支出也相同，一户没有住房债务，另一户每月另需支付固定还款。收入下降时，后者的可压缩余量较小，更可能动用储蓄、增加工作时长或减少培训。其压力来自还款与收入发生在不同时间，而不是从房价跌幅直接推算。即使房价暂时不变，失业也可能造成现金流困难；即使房价下跌，收入稳定且债务较少者也未必被迫出售。

银行和家庭的视角还不相同。银行关注借款人履约与抵押回收，家庭关注居住连续性和未来收入。让银行避免确认损失的安排，不一定同时减少家庭负担；保护家庭居住的安排，也需要说明债权怎样处理。评价住房调整时应分别观察还款期限、利息、住房质量、工作机会与最终损失，不能只选择房价一条曲线。

对租户而言，价格下降可能经由新增出租供给降低租金，也可能因房东减少维护而降低质量；租金是否下降取决于租赁供需和合同调整。这说明买卖价格、租金与家庭福利之间存在传导环节。把所有租户写成房价下跌的直接受益者，同样超出了机制能够保证的范围。

代际转移不仅是遗产金额。父母提供首付款、照顾孙辈、支持失业期生活，都能改变年轻人的选择边界；父母需要长期照护，也可能形成相反约束。仅比较个人工资，可能把家庭支持形成的差异误读为个人能力差异。

% ARGUMENT P169
\phantomsection\label{arg:P169}
家庭支持的一项作用，是改变等待的能力。求职者如果能够在几个月内维持基本生活，就可能比较岗位、完成培训或拒绝拖欠风险较高的工作；必须立即支付照护或债务的人，则更可能接受眼前可得的岗位。这并不意味着等待总能带来更好结果，而是说明同样的才能和努力，在不同预算约束下会产生不同选择。

支持也可以表现为承担失败。青年创业失败以后，如果仍有稳定住处、基本医疗和可恢复的就业渠道，一次损失未必改变整个家庭的位置；如果经营债务与家庭住房相连，失败可能压缩所有成员的机会。所谓创业精神，因此不宜只解释为个人性格。制度提供怎样的退出方式、家庭能够承受多大损失，也影响谁愿意进入以及进入什么行业。

不过，家庭帮助不能全部折算成货币，更不应被一律视为需要取消的不平等。照护、陪伴和知识传递是社会生活的重要部分。公共政策较可行的目标，是使缺少这些支持的人仍能获得基本教育、健康、居住与重新开始的机会。要削弱出身造成的强制性约束，不必预设每个家庭必须拥有相同的生活安排。

因此，机会平等的研究需要把教育、住房、照护与可获得信用联系起来。更好的公共服务能够减少出身对生活机会的部分影响，但不能保证消除所有家庭差异。政策应说明试图缩小的是基本机会障碍还是全部结果差距，两者要求的资源和权利安排不同。

% ARGUMENT P170
\phantomsection\label{arg:P170}
机会研究需要区分进入门槛、过程条件和退出后果。免除一次报名费用可以降低进入门槛，但若学习期间没有生活来源，实际可参与者仍然有限；提供课程可以改善过程条件，但若毕业后缺少公平招聘和岗位，收益不能兑现；允许创业进入可以增加试错，但若退出后长期失去基本保障，失败成本仍然高度不均。三者共同决定机会是否真实。

因而，制度评价不宜只统计某项权利在文本上向所有人开放。要看低收入家庭、流动人口、照护者和残障者能否实际使用，使用需要付出多少额外时间与成本，以及取得服务后是否遭遇新的排斥。相同的报名资格并不自动意味着相同的可行选择；完全相同的结果也不能反过来证明过程已经公平。

将这些条件纳入阶层分析，可以解释社会位置怎样被再生产，也可以识别发生变化的入口。教育扩展、保障改善和企业进入有可能降低原有壁垒，同时产生新的竞争与资格分化。历史分析应追踪这种连续变化，而不是把任何一套当下分类当作永久结构。下一节将把这些维度落实到具体生产关系，说明相近收入背后为什么可能存在不同利益和责任。

\section{阶级分析怎样保留关系性}
经典阶级分析强调人们在生产资料控制和劳动关系中的位置，优点是把不平等与制度关系连接。其问题不在于关注所有权，而在于现实中控制、收入、劳动和风险日益交叉时，单一二分是否仍足以说明差异。

% ARGUMENT P171
\phantomsection\label{arg:P171}
关系性分析与收入排序的区别，可以通过同一个工厂看清。股东、经理、技术人员、生产工人和外包清洁人员共同参与生产，但他们决定投资、安排劳动、分享收益和承担停产后果的权限不同。即使某位技术人员的工资超过一位小股东的当期分红，这种组织关系也不会因此自动倒转。收入告诉我们当期取得多少，关系分析则解释为何以这种方式取得、何时可能失去。

从经典阶级分析发展到中间位置的讨论，并非本书首创。Wright以“矛盾的阶级位置”讨论同时受到支配又具有部分控制权的情形；中国社会分层研究也长期比较以职业和收入构成的等级描述，与以权威、劳动和生产资料关系构成的分类。刘欣、田丰的研究回顾提供了这些路径之间的区分。本文据此采用多维记录，但不把不同传统的所有概念视为天然兼容。\cite{Wright1985,LiuTian2018}

这种区分还限制了政治推论。共同经济位置可以使人们对工资、服务或合同形成相近关切，却不能保证他们拥有同一种政治态度。家庭经历、地区、年龄和具体组织会影响认同与行动。本文讨论的是利益发生的条件和可能冲突，不根据职业替个人判定立场，也不把任何一组人描绘成具有固定道德品质的整体。

阶层之间的边界也有历史性。企业发展可能使自雇者开始雇工，技术人员可能成为合伙人，经营失败又可能使雇主回到受雇位置。分类的用途是追踪这些转换经过了什么条件，而不是宣告一旦进入某类就不能改变。真正需要解释的是哪些转换开放、哪些转换成本高，以及失败后是否仍有基本生活条件。

本文采用四条轴线：是否必须依赖持续出卖劳动维持生活，是否能够支配他人的劳动和重大资源配置，是否拥有可持续资产收入，以及是否能够把失败成本转移出去。它们构成关系描述，不是一套已经统计验证的新阶级分类。收入高低另行记录，避免把收入与权力混为一谈。

% ARGUMENT P172
\phantomsection\label{arg:P172}
第一条轴线是劳动依赖。观察者可以询问：如果停止本人有酬劳动，在不出售基本居住资产的条件下，家庭能维持多久必要生活？这一问题把工资、存款、资产收益和保障联系起来。不过，它需要根据家庭需要确定支出边界，不能凭一个统一月数就宣布社会身份。短期储蓄较多的受薪者仍可能长期依赖劳动；资产收益看似较高但高度波动者，也可能无法稳定退出工作。

第二条轴线是控制。控制不是办公室大小，而是能否决定投资、录用解雇、定价、采购、预算或规则，以及他人能否有效挑战这些决定。要区分制定规则、解释规则与执行规则。窗口工作人员可以影响个别服务流程，却通常无权重设财政分配；企业主管可以安排班次，却未必能够改变总部的产量目标。只按单位性质划类，会把不同层级的责任混在一起。

第三条轴线是持续收益权。自住房提供居住服务，与多项出租资产提供可支配现金流不同；未归属的激励份额，与可以自由出售的股权也不同。记录资产时应写明收益条件、负债、流动性与权利期限。第四条轴线是损失承担和转移：有限责任、保险、担保、合同付款顺序与公共救助预期，都会影响同样经营风险最后由谁承受。实际损失转移必须从合同和结果观察，不能只凭主体“大”或“小”推定。

公共权力、社会保障与家庭积累进入上述轴线的方式不同。公共权力主要影响规则和配置权限，保障主要影响劳动中断后的生活，家庭积累主要影响可动用资源与时间。因此，本书以四条关系轴线组织解释，同时以六类信息进行记录。两者是概念层级不同，不是互相竞争的分类标准。后续研究可以分别检验这些变量是否具有额外解释力。

高级受雇经理可能收入很高、拥有管理权，却仍受董事会和资本市场约束；拥有一套自住房的普通劳动者有资产，却未必能停止工作；负债经营的小老板控制几名雇员，同时受制于平台和大客户。公共部门工作人员的保障可能较稳定，但多数人并不决定公共资本的总体配置。

% ARGUMENT P173
\phantomsection\label{arg:P173}
\subsection{公共资源配置者与公共服务执行者}
公共资源配置者的位置，来自经制度授予的预算、土地、采购、人事或公共资本决策权限。其权力能够影响资源用途和其他主体的机会，却不等于个人对公共资产拥有可出售的私人所有权。判断这个位置，应查看权限范围、集体决策程序、审计约束和责任追究，而不能把经其管理的资产总额写成个人财富。公共受托权的经济影响与私人收益需要分别说明。

配置者的行为还受目标和组织环境约束。若评价强调短期可见产出，维护、服务质量和远期债务可能较难得到同等重视；若责任可以追溯到长期结果，选择可能改变。这是激励机制的条件性推论，不是对所有任职者动机的断言。周黎安和Xiong关于地方激励的研究，为分析组织目标如何进入经济选择提供了历史与模型视角，但不能据此对每个地区作同样判断。\cite{Zhou2007,Xiong2018}

公共服务执行者包括承担具体教学、医疗、行政和运维任务的人员，内部也有管理者、专业人员与辅助岗位。相对稳定的职业保障可能扩大其长期安排能力，但多数执行者并不能决定公共资本整体用途。把他们全部纳入公共资源控制上层，会把组织内等级和劳动关系抹掉。对他们的分析应同时看岗位责任、工作负荷、薪酬支付与实际参与决策的渠道。

这两类位置之间可能有共同利益，例如稳定预算和服务质量，也可能存在考核、用人和责任分配上的张力。受服务居民则需要可用、可靠并可申诉的公共服务。讨论公共部门改革时，必须把管理权限、从业劳动和居民需要放在同一关系内；削减一项预算可以改善账面支出，却可能把工作强度转给执行者、把等待成本转给居民。

\subsection{大型资本所有者与企业控制者}
资本所有者的重要条件，是拥有能够持续参与收益和决策的权益，而不仅是一年收入较高。控股、董事会权限、融资安排和关键知识产权，可以使控制超出直接持股比例；高度分散的小股东则可能分享收益，却很少影响经营。企业控制者与所有者可能重合，也可能分离。分析应沿最终决策权追踪，避免只按工商登记中出现的名称排序。

这一位置的优势通常需要若干条件同时成立：可调动较大规模资源，能够分散部分风险，拥有融资和组织能力，并能参与设定与其他主体的交易条款。不过，规模本身不是无限支配力。技术替代、竞争、监管、债权约束和经营失败都可能削弱这些能力。只有观察到进入壁垒、合同优势或损失分担结果，才可以判断具体权力是否被用于形成租金或转移成本。

资本积累与社会生产之间也不是固定冲突。企业扩大有效产能、培训工人、改进产品，可能创造共同收益；在分配新增收益、确定工时、安排供应商付款和处置亏损时，则可能出现利益分歧。解释这种关系需要把创造的总收益与分配规则分别考察。不能从存在私人利润推出没有社会贡献，也不能从企业提供岗位推出现有合同必然公平。

\subsection{中小雇主与个体经营者}
中小雇主处在两种关系的交汇处：对雇员有一定用工和管理权限，对银行、房东、大客户或平台则可能缺少议价能力。向下能够安排劳动，并不意味着向上能够决定价格和结算。其利润可能包含本人长时间劳动的报酬、投入资本的回报和组织收益。若不分解这些成分，就容易把所有小经营者一律等同于大型资本控制者。

假设一家维修店雇用两名工人，店主同时维修、接单并承担租金。增加订单可以使三人都受益，但收入如何在工资与经营剩余之间分配仍需协商；大客户拖欠货款，又可能使店主尝试延后支付工资。后一个压力可以解释行为发生的条件，却不能免除已经形成的支付责任。把上下游现金流接起来，既能避免道德化归因，也能防止用经营困难掩盖风险转移。

个体经营与自雇者的关键区别，是其收入是否主要依赖本人劳动以及是否雇佣他人。独立医生、自由设计师、街边经营者与平台接单人员都可能被统计为自雇，但客户自主性、专业壁垒和保障相差很大。营业执照只说明一种登记形式，不能替代劳动控制分析。研究还应记录家庭成员无偿参与，以免把家人劳动形成的收入全部计作资本利润。

小经营者的政策利益也不是只有降低税费。及时结算、租约稳定、公平进入、低成本争议处理和退出后的基本保障，可能直接减少经营不确定性。部分措施同时有利于雇员，但工资、社保和安全投入依然存在分配问题。把共同利益与具体冲突分别说明，才能避免把“帮助小企业”当作一张无需核对最终受益者的空白支票。

\subsection{稳定受薪者、专业人员与市场化白领}
稳定受薪者的优势主要体现在收入连续性、保障和职业预期，未必体现为对他人劳动的控制。一份相对稳定的合同可以支持住房、育儿和长期学习，也可以减少短期收入波动造成的高成本借贷。不过，稳定程度仍取决于组织财力、岗位类型和制度兑现。不能把单位名称直接转换成终身安全。

专业技术与市场化白领内部，一端可能凭稀缺技能、客户关系或管理权限获得较大选择空间，另一端可能在高度标准化流程中接受严格绩效与替代压力。文凭、职业声望与实际自主性不是同一个变量。专业人员的高收入可能来自培训成本、稀缺能力，也可能部分来自资格壁垒和组织租金，需要具体研究。将“知识”笼统当作一种与机器完全同类的资本，会忽略其必须依附于健康、时间与持续学习的性质。

人工智能和组织重组可能改变这种内部差异。某些常规知识任务更容易标准化，责任判断、客户信任和组织协调则可能更加重要。技术冲击不一定先影响收入最低者，也可能改变原本被认为稳定的职业阶梯。对这一群体，应同时观察工作内容、技能可迁移性、劳动时长、绩效波动和家庭债务。高薪但高负债、技能高度专用的人，面对失业时可能缺少看上去那么大的缓冲。

\subsection{产业与服务业劳动者、平台劳动者}
产业与服务业劳动者主要依赖持续劳动取得收入，但其内部控制条件不同。熟练技术工人的工艺知识可能形成议价基础，临时替代性较强的岗位则可能更受排班和订单影响；面对面服务劳动还依赖地点、顾客评价和情绪劳动。工资多少、劳动多久、何时得到支付、能否拒绝不安全任务，应当合并观察。农民工统计提供了重要就业事实，但不能把其所有成员固定为同一种职业位置。\cite{Migrant2025}

劳动者的议价能力来自可行替代。附近是否有其他雇主、迁移是否需要重新安排子女教育、是否能够承受空档期、技能能否被其他企业认可，都会影响拒绝不利合同的能力。一项保障政策即使没有立即增加名义工资，也可能通过改善退出条件改变协商。相反，名义工资提高若伴随工作强度明显上升，生活改善的程度就要重新判断。

平台劳动者还面对交易入口与任务规则的控制。拥有车辆或工具，只说明承担了部分资产投入；价格、客户接触、派单和账号处罚如果由平台决定，其经营自主性就受到限制。另一方面，部分从业者确实可以跨平台接单或独立建立客户，不能一律按完全受控的典型处理。需要依据实际安排分类，而非依据合同标题预先得出劳动关系的法律结论。

这一群体的风险往往与工具折旧、等待时间和任务波动相连。若只看一单到账额，而不扣除设备、能源和无单等待，就会高估有效报酬；若把所有上线时间都视为同样强度劳动，也可能忽略其间真实可自由支配的部分。第九章将进一步分析记录和协商规则。共同的研究要求，是让从业者承担的时间与成本能够进入评价。

\subsection{农村家庭中的不同经济位置}
农村家庭不能按居住地点合并成单一阶层。以自家劳动从事农业、租入土地扩大经营、雇佣农业工人、主要依靠外出工资或主要依靠转移收入，各自面对不同控制与风险。土地承包、经营安排和实际劳动需要分开记录。经营面积较大不等于拥有土地所有权，拥有承包权益也不等于主要现金收入来自农业。

兼业家庭同时连接城市劳动市场与农村生活。城市岗位减少可能使成员返乡，但返乡需要居住、食物、照护和新的收入来源；保留土地可以提供部分缓冲，却不能自动承担所有教育、医疗与养老费用。温铁军关于城乡成本转移的讨论有助于提出这种跨部门问题，具体承担程度仍需家庭资料核验。不能把家庭吸收冲击直接称作社会已经无成本化解风险。\cite{Wen2021}

规模经营者更多暴露于投入价格、天气、产品销售与周转风险，农业雇工则关注工作季节、工资和安全，小农户可能更受市场进入与服务获取约束。乡村老年家庭的主要困难又可能是照护和医疗，而非扩大经营。相同政策对这些位置作用不同：补贴设备可能有利于扩大经营，却未必解决老年人的基本服务；改善交通与就医则可能同时支持多类家庭。

城乡关系也不是静止边界。家庭成员可以跨地区配置劳动和照护，年轻人的教育和迁移改变家庭未来收入结构。分析农村阶层，要跟随这些关系移动，而不是只在村庄常住人口中寻找全部经济活动。对外出成员与留守成员分别记录，才能看清城市工资怎样支持乡村生活，乡村照护又怎样支持城市就业。

这种多轴分析不应变成“每个人都很复杂，所以权力差异无法讨论”。相反，它使研究可以观察哪些人同时拥有控制权、稳定收益和损失转移能力，哪些人同时缺少这些条件。较稳定的组合可能形成社会位置，但需要调查验证其规模、流动与代际延续。

% ARGUMENT P174
\phantomsection\label{arg:P174}
前述类别是一幅关系图，不能当作已经完成的人口普查。有人同时是专业受薪者、小额出租者与照护者，有人从农业兼业转入平台接单；相互交叉并不使分类失效，而是要求研究先确定问题。例如研究雇佣控制时以本人岗位为单位，研究失业缓冲时以家庭预算为单位，研究公共权限时以职位实际授权为单位。不同单位的比例不能直接相加。

更值得检验的是有利条件是否相互强化。稳定收入能否帮助取得较低成本信用，信用能否支持住房与教育，家庭积累能否延长求职等待，控制权能否使风险转给其他主体；相反，收入波动、保障不足和照护压力是否会共同压缩选择。这些反馈可以产生持久分层，但每个连接仍需证据，不能因为机制听来合理就宣称全国必然如此。

\begin{longtable}{p{.18\textwidth}p{.24\textwidth}p{.24\textwidth}p{.22\textwidth}}
\caption{社会位置的关系比较：分析维度，不是人口规模估计}\label{tab:positions}\\
\toprule 位置 & 主要控制或依赖 & 需要观察的风险 & 不能据以直接判断的标签\\\midrule\endfirsthead
\toprule 位置 & 主要控制或依赖 & 需要观察的风险 & 容易混淆的标签\\\midrule\endhead
公共资源配置者 & 经授权决定部分预算与资源用途 & 决策责任、监督及公共损失归属 & 管理资产不等于私人财富\\
企业所有者与控制者 & 收益权、投资权及组织权限 & 经营损失、担保、竞争与融资 & 高收入不自动等于控股\\
中小雇主 & 局部用工控制与上游依赖并存 & 回款、租金、个人担保 & 小规模不消除雇主义务\\
自雇与个体经营者 & 本人劳动、客户关系及工具 & 停工、接单与保障中断 & 登记身份不等于实际自主\\
稳定受薪者 & 依赖劳动，收入连续性较强 & 组织支付、岗位和家庭照护 & 单位类型不等于配置权\\
专业人员与白领 & 技能、资格及不同程度管理权 & 技能贬值、绩效与债务压力 & 文凭不等于长期安全\\
产业与服务业劳动者 & 岗位、工时及劳动合同 & 工资中断、伤害、迁移成本 & 职业名称不等于同质处境\\
平台劳动者 & 任务入口、规则与工具投入 & 账号、等待、事故与收入波动 & 自备工具不等于独立定价\\
不同农村家庭 & 土地、经营、务工和转移组合 & 农业风险、城市就业与照护 & 农村居住不等于纯农业收入\\
\bottomrule
\end{longtable}

表中没有把青年、老人或女性列为与职业并列的互斥阶级，因为这些身份贯穿多种位置。它们会改变照护、健康、经验与制度待遇，必须作为交叉维度研究。这样既能看见共同约束，也避免把一位高资产退休者与一位缺少基本养老来源的老人误认为经济利益完全相同。

\section{分配、公平与效率的关系}
收入差距可以激励技能积累和风险承担，也可能来自排他权利、信息优势和历史机会。仅凭差距大小，无法确定每一部分的效率作用。有效政策应尽量减少缺乏社会贡献的租金，并保障基本机会，而不是简单把所有差异视为同一种现象。

% ARGUMENT P175
\phantomsection\label{arg:P175}
评价差异的来源，需要构造比较问题。若取消某项排他资格后服务质量保持、价格下降且新进入增加，原收益中可能存在由限制进入维持的部分；若降低合理培训回报后人才供给明显减少，则激励约束需要认真考虑。这些判断不能由“收入高”三个字直接推出。租金、创新回报与劳动报酬在一份收入中可能并存，分析应尽可能拆开其形成条件。

效率也需要明确时间和主体。某家企业通过延迟支付供应商保持现金，可以改善自己的短期周转，却增加供应商融资与停工成本；一种教育制度提高少数人的考试成绩，也可能增加其他家庭的时间与现金压力。局部产出提高不能直接作为社会净收益。本书要求同时核对成本所在账户，正是为了使效率评价与分配评价相互连接。

公平标准则包含独立的价值判断。有人更重视基本生活，有人更重视贡献与回报，有人更重视起点机会，这些标准不总产生相同排序。经验研究可以说明一个方案会怎样改变资源和行为，却不能仅凭统计结果替社会决定全部权重。清楚列出标准，比把自己的偏好包装成不容争论的经济规律更有利于讨论。

再分配也不只有税后转移。教育、医疗、托育、公共交通、住房与合同执行，可以改变市场收入形成之前的条件。劳动者更容易获得技能和迁移，企业更容易公平进入市场，这些安排可能同时改善效率与分配。其效果仍取决于质量和实际覆盖。

% ARGUMENT P176
\phantomsection\label{arg:P176}
市场收入形成以前的条件，常决定事后需要多少补救。如果公共交通扩大了可达岗位，劳动者可以比较更多雇主；如果托育与轮班时间相配，照护者可能保留就业；如果小企业能够及时追索货款，其融资需求可能减少。上述措施都可能改变初次收入，但具体效果依赖实际使用与供给质量，不能仅根据投入金额认定成功。

不同阶层之间还可能存在合作空间。雇员与小雇主都可能希望上游按时付款，居民与公共服务人员都可能需要稳定运营预算，消费者与小商家都可能从透明平台规则受益。但一项共同改善不意味着其他分配冲突消失。付款恢复以后，工资与利润如何分享仍需安排；平台费用降低以后，商家是否降价也取决于竞争。

将共同约束与分配争议分开，有助于确定改革顺序。先减少无效等待、任意收费和信息不透明，可能扩大可分配收益；随后如何分配仍需要制度与协商。研究不能因为存在合作收益就忽略权力差异，也不能因为存在权力差异就否认任何共同改善的可能。

另一方面，任何税收和支出都可能改变行为。提高某项资产税，需要考虑估值、流动性、税负转嫁和资本配置；扩大补贴，需要考虑退出门槛与行政成本。把这些问题说清楚，不等于反对再分配，而是避免政策目标被执行方式抵消。

% ARGUMENT P177
\phantomsection\label{arg:P177}
税负转嫁不是一句“最后都由消费者承担”可以概括的规律。转嫁取决于供求弹性、替代渠道、竞争结构、合同和时间。若某项稀缺资源供给固定，对其收益征税的结果可能与对可迁移投资征税不同；若租赁市场供给紧张，房东可能尝试涨租，但能否实现仍受租户支付能力和替代选择约束。没有这些信息，就不应承诺税负必然留在法定缴纳者身上。

政策还可能改变收入的记录方式。若不同身份适用的负担差距较大，主体可能转变登记、合同或报酬形式，表面上的类别规模变化不一定代表真实生产关系发生同等改变。因此，改革评估需要观察活动、收入与责任，而非只观察登记数量。对平台、自雇和小企业尤其如此。

避免副作用不意味着等待完美信息。可以先用明确边界、有限预算和可退出安排实施，再根据实际结果修订；但试点选择不能只挑最容易成功的群体。最需要帮助的人常具有复杂负债、迁移和照护条件，若把他们排除，得到的平均效果就不能代表政策目标人群。

\section{家庭调查应怎样使用}
CHFS申请说明显示，调查完成、内部开放和对外学术申请的范围并不一致。本书没有取得受限原始数据，不能声称使用了2025年全国微观样本。公开的旧波次结果可以用于历史机制研究，但必须标明真实调查期和样本限制。\cite{CHFS}

% ARGUMENT P178
\phantomsection\label{arg:P178}
公开调查报告可以告诉读者研究采用何种样本、调查问题与年份，但它不能替代未取得的原始数据。没有微观样本，就无法重新处理权重、缺失、异常值和家庭结构，也无法检查某个细分群体是否样本过少。引用已发表结论时，本书只采用公开材料实际支持的层级，不把二手汇总改写成自己计算的发现。

家庭调查中的“没有回答”也不能随意按零处理。拒绝披露资产、忘记债务、无法估值与确实没有该项资产，是不同状态。把它们合并会改变财富分布。相同问题还出现在控制权测量：受访者不知道组织决策过程，不表示组织中不存在控制关系。需要通过问卷设计、行政资料与组织案例相互补充。

资料开放条件本身是研究边界之一。公开申请范围可能随时间变化，本书记录核查日期和可读说明，而不承诺某一波次一定可以获得。后续如果取得新样本，应重新核对地域范围、抽样框与调查期，不能仅把报告标题中的年份更新就称作完成了2026年实证。\cite{CHFS}

若要检验本文的社会位置框架，后续调查应同时收集收入来源、资产与负债、管理权限、保障、照护责任和冲击应对。纵向追踪比一次横截面更有助于判断位置是否稳定。还需特别处理迁移人口、顶端家庭和退出样本的偏差。

% ARGUMENT P179
\phantomsection\label{arg:P179}
一个可操作的研究设计，应先固定需要解释的结果，例如失业后三个月是否削减基本支出、能否获得新岗位，或者企业停业以后家庭是否失去住房；再检验劳动依赖、控制、资产收益、保障与照护能否解释这些差异。若在已有收入和教育变量之外，新维度没有增加解释力，就应承认框架对于该问题的价值有限。

位置流动还需要区分绝对流动与相对流动。经济整体扩张可以使很多人的实际生活改善，但不同家庭在排序中的位置不一定变化；新增专业岗位可以扩大进入机会，但出身差异仍可能影响谁进入。只报告某类岗位人数增加，无法证明机会已经均等；只报告排序稳定，也不能否认绝对生活条件改善。

纵向研究会遇到退出样本的问题。迁移、破产、家庭解体或健康恶化者更难追踪，而这些恰是研究风险最关心的情况。若失去他们，剩余样本可能显得比实际更稳定。可以比较持续样本与补充样本、报告失访结构并进行敏感性分析；但任何处理都不能保证完全修复未观察差异。本书提出的是检验路径，没有假称已经执行这些步骤。

本章保留阶级分析对关系和控制的关切，修订单凭收入、职业或资产名义归属划定全部位置的做法。财富分布提供重要线索，但判断生活机会和风险，还需要知道资产能否兑现、保障是否可用以及谁有能力改变合同。

% ARGUMENT P180
\phantomsection\label{arg:P180}
将本章放回全书，阶层分析的功能是给各章的机制找到具体承担者。价格下降经过合同传给谁，技术进步由谁决定采用，债务展期保护了谁的现金流，照护缺口由哪个家庭成员补上，都需要明确社会位置。没有这一步，宏观增长和风险变化容易停留在没有人的部门名称之间。

与此同时，位置不能成为替代证据的万能答案。某人属于企业控制者，不足以证明其每次决定都为自己转移成本；某人依赖工资，也不足以证明每项工资差异都源于单一支配关系。研究仍须说明行为经过什么权限和合同实现，有哪些竞争解释，以及不同选择是否可行。关系分析提供问题的结构，具体结论则必须由观察约束。

本章因此把新版阶层分析从群体描述推进到可追踪的关系体系。重点不是再造一份永久名单，而是识别哪些制度使控制与责任分离，哪些家庭能够承担等待和失败，哪些人持续用时间与健康填补缺口。后续数字经济、全球生产和危机章节，将沿这些关系分析技术和冲击的传递。

\chapter{平台、数据与人工智能}
\section{数字化改变了协调成本}
平台把原本分散的买卖、支付、物流和信息连接起来，能够降低搜索和交易成本，扩大商家覆盖。它也可能积累对入口、排序、规则和结算的控制。当越来越多交易依赖同一入口时，平台就不仅是一个出售服务的企业，还成为其他企业和劳动者开展活动的条件。

% ARGUMENT P181
\phantomsection\label{arg:P181}
理解平台，可以先比较两种组织交易的方式。分散交易中，买家逐个寻找卖家、核验质量并安排支付；平台将搜索、评价、履约和结算集中，节约重复投入。商家因此可能接触原来无法触及的顾客，消费者也可能更容易比较。这部分收益来自协调改进，不能因为平台后来形成较强控制，就否认其最初和持续提供的服务。

控制产生于依赖逐渐累积。商家在某平台形成评价和回头客，劳动者积累接单记录，消费者习惯其支付与售后；如果这些资源离开平台就难以保留，退出成本会提高。平台修改费率或规则时，参与者不只比较当期价格，还须考虑放弃既有关系的损失。历史形成的投入由此进入当前议价。

这一机制在不同群体之间作用不同。大型商家可能维持自有渠道，小店可能更依赖单一入口；具有独立客户的专业人员可能把平台作为补充，主要依靠派单的劳动者则较难转移。分析平台权力，需要把这些替代渠道画出来。用户数量可以说明规模，但不能单独证明每类参与者都具有同等依赖。

双边市场理论说明，对用户一侧低收费、对商家或广告主另一侧收费，可以是协调网络效应的方式。不能从一侧免费直接推出该侧没有价值，也不能从另一侧收费直接推出垄断滥用。Rochet和Tirole的框架强调不同用户群之间的相互影响，适合用来识别收费结构，但具体福利效果仍取决于竞争和替代。\cite{Rochet2003}

% ARGUMENT P182
\phantomsection\label{arg:P182}
双边市场的特点，是一侧参与会改变另一侧获得的价值。平台若先降低消费者进入成本，可以吸引更多消费者，再使商家愿意付费；也可能先吸引商家供给，再改善消费者选择。因此，评价收费不能只看某一侧的价格，而应看总收费、参与规模和两侧相互作用。Rochet与Tirole的研究说明了这种价格结构问题，不为任何实际收费自动提供正当性。\cite{Rochet2003}

一个假设例子可以帮助比较。平台向消费者免收使用费，但向商家收取服务费。若服务费对应更低的获客和履约成本，商家净收益仍可能增加；若商家已经难以离开，而费率上升没有对应服务改进，其净收益可能下降。相同的“免费加抽成”形式，可以容纳不同竞争状态，判断须看成本、替代渠道和规则变化。

还要观察跨侧转嫁。商家可能把费用计入售价，平台也可能补贴某些商品吸引流量，最终消费者支付的价格未必直接显示平台收入。在有补贴和返还时，账面抽成率尤其不能替代净费率。研究应把佣金、广告、支付、配送和促销承担分别记录，避免从一个醒目的费率得出全部分配结论。

中国平台分析应把消费者便利、商家成本、劳动条件与创新进入放在一起。低价可能来自技术效率、交叉补贴或对供应商的压力；提高抽成可能用于服务改进，也可能反映缺少替代渠道。要区分这些机制，需要成本、质量、进入和合同的证据。

% ARGUMENT P183
\phantomsection\label{arg:P183}
低价形成的路径不同，长期含义也不同。技术改善降低每单真实资源耗用，可以支持较持久的低价；资本补贴可以暂时压低价格，但持续性取决于融资与未来收益；供应商让利可能来自效率改善，也可能来自议价不足。判断应追踪商品质量、交付时间、劳动投入和供应商现金流，而不仅比较消费者支付金额。

对家庭而言，便宜商品可能提高购买力，同时家庭成员也可能作为商家或劳动者承担费用和收入压力。消费者、经营者、雇员不是互不相交的三群人。因此，需要比较家庭净效果，但又不能用总体平均抵消集中损失。某些家庭可能主要消费受益，另一些家庭却集中暴露于就业和经营调整。

一个可检验的区分是：费用提高以后，平台履约质量、纠纷处理和有效成交是否同步改善？如果改善不足而退出困难持续，控制解释更值得考察；若服务成本确实上升、竞争者费率相近且商家净收益改善，单纯的租金解释就较弱。两种情况也可能混合存在，数据应允许这种结果。

\section{网络效应并不总导致永久垄断}
更多用户可以吸引更多商家，更多商家又吸引用户，这种正反馈有集中趋势。但用户是否同时使用多个平台、数据能否携带、切换成本有多高、服务是否存在差异，都影响集中程度。大规模与持久市场势力相关，却不是同一个概念。

% ARGUMENT P184
\phantomsection\label{arg:P184}
网络效应与规模经济需要区分。服务器和研发成本分摊到更多交易，是生产成本方面的规模经济；更多买家吸引更多卖家，是参与者之间的网络效应。两者都可能支持集中，但约束不同。技术进步可以降低建立新平台的固定成本，却未必解决新平台缺少交易双方的启动困难。

多平台使用是否有效，取决于重复参与的成本。消费者同时安装两个应用可能很容易，商家维护两套库存、广告和售后却更复杂，劳动者同时接不同平台任务也可能发生时间冲突。仅以“存在其他应用”证明竞争充分，会忽略参与另一渠道是否真正可行。

网络还可能出现拥挤、质量下降和偏好分化。更多商家使选择增加，也可能提高搜索难度；更多劳动者接单可以缩短等待，也可能降低人均订单。专业平台和地方服务因而可能保有空间。要判断平台能否长期集中，需要观察这些正负效应，而不能把早期增长的正反馈无限外推。

当平台从交易中学习需求并改善匹配时，数据可能提高效率；当它利用信息优势优先扶持自营业务、限制商家接触顾客或提高退出成本时，就可能改变竞争。规则透明并不要求公开每一行程序，而是让收费、处罚、排序的重要条件可理解、可申诉和可监督。

% ARGUMENT P185
\phantomsection\label{arg:P185}
平台同时充当规则制定者和交易参与者时，可能出现利益冲突。假设它能够观察商家的销售和退货信息，又经营同类商品，就可能更早识别需求；这种信息优势可能用于改进供给，也可能用于限制原商家接触顾客。哪一种情况发生，需要看数据用途、排序变化、合同安排和竞争结果，不能由存在自营业务直接推出违法或效率下降。

透明规则也有边界。完全公开反欺诈细节可能帮助规避，公开个人交易信息可能损害隐私；但费用如何形成、处罚依据是什么、重大规则何时改变、如何提出异议，可以向受影响者说明。可解释性不等于公开所有技术细节，而是使责任能够被理解和检验。

对小商家与劳动者，申诉渠道的成本特别重要。理论上可以申诉，但需要长时间等待或复杂材料，可能使权利无法使用。评价应记录恢复经营所需时间、争议期间收入损失及错误决定是否被纠正。一个渠道存在与一个渠道有效，是两个不同事实。

监管需要避免两种简单化。只按企业规模限制增长，可能损失网络协调收益；只以消费者当期价格低为标准，又可能忽略未来进入障碍和商家、劳动者的成本。评价时间跨度和涉及群体，应在政策开始时说明。

% ARGUMENT P186
\phantomsection\label{arg:P186}
监管比较的对象应当是具体安排。限制排他交易、要求费率通知、提供可迁移记录和强化争议处理，各自改变不同约束，成本也不同。规模上限虽然容易表达，却可能迫使重复建设；完全不干预又可能使入口控制积累为长期壁垒。需要先确认问题，再选择能改变该问题的工具。

时间跨度会影响评价。短期补贴可能让消费者受益，长期若进入减少、服务恶化，净效果可能不同；相反，一项提高保障的规则可能在初期增加价格，却减少事故和家庭损失。政策应预先说明考察多久、跟踪哪些群体，避免在不同时间窗口之间挑选有利结果。

还需要观察规则的组织后果。平台可能把任务转交新的承包商，也可能改变合同形式或减少服务区域。如果责任链随之变得更难追踪，文本上的保障提升未必兑现。有效评价应沿实际控制和支付关系追踪调整，不能只比较改革前后平台自身财务报表。

\section{数据的非竞争性与权利边界}
一份数据被一家公司使用，通常不妨碍另一家使用，但隐私、商业秘密和安全要求限制共享。数据也不自动具有高价值：质量、时效、代表性、可连接性和处理能力共同决定用途。把所有数据统一叫作“新石油”，容易忽略复制成本和权利结构的差异。

% ARGUMENT P187
\phantomsection\label{arg:P187}
非竞争性不意味着无成本。一份记录可以重复使用，但获取、清洗、标注、校验与安全管理都需要劳动和设备；将不同来源的数据连接，也可能产生质量与隐私风险。数据的复制成本较低，与形成可靠用途的总成本较低，是两件事。

数据的互补性尤其重要。设备运行数据如果缺少维修记录和工况说明，可能难以解释故障；大量用户行为如果不能区分异常、欺诈和真实偏好，也未必提高预测。数据资产的价值因此不能只用字节数或记录条数表示。拥有处理能力和相关业务知识的主体，可能更容易从相同数据中取得收益。

权利安排需要对应用途。允许统计分析，不一定意味着允许识别个人；允许某项公共服务使用，也不等于可以任意转售。技术上能够复制，不自动产生道德或法律上的使用资格。本书讨论的是使用收益、控制与责任的经济关系，不把数据这一总称当作具有唯一所有权答案的物品。

数据价值往往由多人活动、公共规则、设备和组织共同形成，很难为每次点击计算唯一边际贡献。要求精确按点击分红，可能产生高于收益的计量成本。可行选择包括特定数据用途的授权、集体协商、公共数据开放、可携带接口和对实际损害的责任安排，各自适用不同条件。

% ARGUMENT P188
\phantomsection\label{arg:P188}
共同形成的数据很难用单个行为的边际贡献准确分账，因为贡献取决于其他数据和模型。一个罕见病例记录可能对特定任务非常重要，在另一任务中几乎无用；一条普通记录单独价值很低，却可能与大量同类记录共同提高可靠性。这种互补性使统一按次数付费难以对应实际收益。

这不意味着参与者没有权利或无法取得收益。可以对明确用途约定报酬，对可识别损害建立责任，对公共投入形成的资源规定服务义务，也可以在具备代表性的范围内进行集体协商。选择哪一种方式，需要比较计量成本、谈判能力和执行条件，不能从无法精确估值跳到可以无条件占用。

对劳动者而言，还要区分完成明确生产任务与日常使用服务时留下记录。受雇标注、测试和审核具有具体劳动过程，其报酬与工作条件可以直接研究；一般用户产生记录，则涉及消费、授权和隐私等关系。把两者全部称为同一种数字劳动，可能遮蔽更可操作的责任边界。

公共数据开放可以促进创新，也可能让具备算力和渠道的大企业获得更大收益。若要扩大受益，需要考虑接口成本、质量说明、隐私保护和中小主体的使用能力。宣布数据“属于公共资源”，并不自动保证公众获得可兑现收益，这与第四章的产权传递问题相通。

% ARGUMENT P189
\phantomsection\label{arg:P189}
公共数据开放的实际门槛，往往在接口和持续维护。格式不一致、说明不完整、更新中断，会使有资源的大机构更容易清理和接入，小企业则难以承担前期成本。开放数量增加而可用性不变，未必扩大创新机会。可以把准确性、更新时效、接口稳定性和使用成本作为评价对象。

受益还取决于下游竞争。如果只有少数主体拥有互补算力、渠道或专有数据，开放公共数据可能主要提高其收益；若工具和基础能力更容易获得，更多经营者才可能进入。由此产生的政策问题，是怎样降低合理使用门槛并约束损害，而非简单按企业大小决定是否允许使用。

公众收益可以有多种形式：更好的服务、较低的价格、研究成果或按约定回流的收入。应事先说明目标和责任，避免把下载次数当作最终福利。公共数据的维护费用、隐私风险和服务改进应一起入账，这也是公共产权向实际收益传递问题在数字领域的延伸。

\section{人工智能暴露度不等于失业率}
生成式人工智能能够处理文字、代码、图像和部分知识任务，但职业由多种任务组成。某些任务可以自动完成，另一些仍需要现场判断、信任、责任和人与人之间的协调。ILO在2025年的职业暴露指数研究中，采用任务层面评估；这种方法提示可能受影响的范围，而不是预测同等比例岗位必然消失。\cite{ILO2025}

% ARGUMENT P190
\phantomsection\label{arg:P190}
职业暴露指标回答的是某些任务与技术能力之间有多少重合。它没有自动包含采用成本、市场需求、责任安排和企业组织，也没有直接观察技术使用以后的就业变化。一项职业包含起草、沟通、现场处置和责任签署，只有部分任务受到影响时，职业可能被重组而不是整体消失。\cite{ILO2025}

暴露与脆弱也不同。技术能够辅助某类专业人员，可能提高其服务能力；但是否改善收入，还取决于客户需求、职业准入和雇佣安排。另一类岗位即使直接暴露较低，也可能因上游企业重组或下游需求改变而受到间接影响。仅按技术匹配程度排列职业风险，会遗漏产业联系。

本书因此将技术可能性、实际采用、任务调整和最终就业分成四个观察阶段。前一阶段变化是后一阶段的条件之一，却不是同义词。对中国情境，后续研究应从企业实际工作流出发，分别记录哪些任务被取消、哪些被增加、审核由谁承担以及人员如何转岗；不能把全球指数当作全国失业预测。

从技术演示到稳定生产还要跨过数据质量、集成、错误责任、隐私、维护和培训等成本。若一项工具减少起草时间，却增加审核时间，它的净效率提升可能小于演示中的节省。对医疗、工程和法律等高责任任务，少量严重错误也可能限制完全替代。

% ARGUMENT P191
\phantomsection\label{arg:P191}
评估净效率应从完整任务链开始。假设原任务需要十小时，其中起草六小时、核验两小时、沟通两小时；工具把起草压到两小时，却使核验增加到三小时，那么总时长从十小时降到七小时，而不是降到起草所需的两小时。这是说明核算方法的假设例子，没有代表任何行业平均效果。还应加上培训、接口、维护与错误处理的成本。

错误成本也不能只按平均准确率排序。在低后果、可快速修正的任务中，少量错误可能容易处理；在后果严重且难以撤回的任务中，审核和责任要求会更高。企业采用因此取决于错误分布、可发现性与责任承担，而不只是演示表现。把被节省的时间全部计作利润，会漏掉为稳定使用所增加的劳动。

这会改变专业人员的位置。具有核验能力和承担责任资格的人可能成为互补环节，初级人员则可能减少练习基本任务的机会。短期效率提高是否削弱长期技能形成，需要追踪培训和职业晋升。技术治理不仅涉及当前岗位数量，也涉及下一代专业人员如何获得经验。

本文不把全球暴露比例直接乘以中国就业人数。中国的产业结构、企业规模、工资、软件投入和管理条件都会影响采用。小企业可能因缺少集成能力而滞后，也可能借助通用工具跨过部分固定成本。哪种机制更强，需要行业和企业数据。

% ARGUMENT P192
\phantomsection\label{arg:P192}
采用差异可以从固定成本和可变成本理解。大型企业可能有数据、管理和工程能力，却也可能面对复杂系统与较高合规成本；小企业可能缺少集成人员，却能够直接使用通用工具处理简单任务。因此，规模对采用的影响不是单向的，需要按任务和组织复杂性区分。

Brynjolfsson、Li和Raymond对特定客服工作场景的研究，观察到辅助工具带来的生产率改善在经验与技能不同的人员之间并不相同。这为知识扩散和人机互补提供了一个经验案例，但其企业、任务和质量控制条件限制了外推。本文不把该研究的效果量套用于中国所有白领，也不把短期处理效率直接等同于长期工资增长。\cite{AIWorkNBER}

研究中国采用过程时，需要关注采购者、实际使用者和受影响者是否同一群体。管理层可能按减少成本选择系统，员工关注工作负荷与技能，客户关注质量与责任。若指标只奖励处理速度，工具可能增加数量却降低复杂问题的解决质量。需要同时记录速度、错误、返工、客户结果和劳动条件，才可能判断真实净收益。

\section{技术收益如何在产业链分配}
人工智能的收益可能流向模型提供者、算力与设备供应商、拥有独特数据的企业、使用工具的员工或终端消费者。市场进入与互补资产决定分配。工具成本下降，既可能使应用企业受益，也可能因普遍竞争而主要转化为更低价格。

% ARGUMENT P193
\phantomsection\label{arg:P193}
收益分配取决于互补要素的稀缺性。若模型和基础工具竞争充分，使用价格下降，拥有客户关系、专有工作流程或可靠数据的企业可能保留更多收益；若关键算力或某项知识产权限制进入，上游供应者可能保留较高回报。不同阶段的瓶颈会移动，不能把某一时点的利润格局当作技术的永恒属性。

员工能否分享收益，还取决于其替代条件和组织规则。工具使稀缺专业能力更容易复制，可能削弱原有工资溢价；工具提高某些人员协调更多任务的能力，也可能增强其贡献与谈判位置。相同技术对新手、资深人员、管理者和外包劳动者的影响可能不同。应观察任务变化与合同，而不是以“知识劳动者”作为同质整体。

消费者通过价格和质量分享收益，公共部门则可能通过税收与公共服务获益，但这些通道都需要制度条件。竞争不足时，成本下降未必充分转成降价；跨主体收益转移也可能使税基变化。因此，技术生产率与广泛生活改善之间仍有传递环节。它们不是技术之外的次要分配问题，而会反过来影响采用和需求。

假设一项辅助系统把某任务时间从十小时降到六小时。若需求固定，企业可以减少用工；若低价扩大需求，任务总量可能增加；若员工能够参与协商，可以争取缩短工时或承担更有价值的新任务。仅凭四小时节省，无法推导工资必涨、岗位必减或劳动从此不再重要。

% ARGUMENT P194
\phantomsection\label{arg:P194}
把时间算例写成条件关系更清楚。原来每单位任务需要十小时，采用后需要六小时，单位时间下降40\%；若产量不变，总工时相应下降。若任务量扩大到原来的五分之三的倒数，即约1.67倍，总工时才回到原水平。扩张不足这个幅度，总工时可能减少；超过这个幅度，则可能增加。这只是一项固定任务结构下的算术关系，不包含新任务和质量变化。

需求是否扩大，取决于价格下降幅度、客户需求与其他供给约束。企业即使降低成本，也可能不降价；即使降价，某类服务的需求也未必显著增加。更进一步，总工时增加不保证岗位数量同比例增加，因为工时安排可能改变；岗位增加也不保证每名劳动者收入提高，因为技能与合同可能变化。

这个算例说明，技术对就业的推论至少需要需求、价格、任务结构和工时制度四组信息。它也给出了可以协商的空间：新增收益可以进入利润、工资、价格、培训或休息，具体组合受竞争和权利安排影响。把某一种结果说成技术的唯一必然后果，会提前关闭本来需要研究和协商的问题。

所谓数字剩余，应当分别记录提高效率形成的可分配收益、控制入口形成的租金、知识产权回报以及风险外部化。把这些全部命名为一种全新的剩余价值，并不会自动增加解释力。理论创新应能说明在什么竞争和权利安排下，哪一部分收益会改变。

% ARGUMENT P195
\phantomsection\label{arg:P195}
剩余的分解可以从可复算的经营账户出发。收入减去投入、折旧、工资、必要维护和融资成本后形成会计余额，但其中包含多少创新回报、市场租金或未被计入的社会成本，还需要进一步分析。平台规模扩大以后利润上升，可能同时来自真实协调效率、固定成本分摊和议价变化，不能靠给利润换一个名称解决识别。

若要提出新的理论命题，至少应说明什么观察与现有解释不同。例如，控制入口的企业是否能够在服务质量不变时持续提高净费率，且这种能力是否随数据迁移和替代渠道改善而减弱？这一命题比“数据产生新剩余”更具体，因为它给出了条件、结果和反证。

马克思主义对劳动控制与剩余分配的关切在这里仍有作用，但数字技术并不自动使所有非工资收入成为同一种对象。研究者需要同时容纳创新成本、组织服务和排他收益。区分这些部分，是为了更准确地判断哪些制度促进能力形成，哪些制度使收益与贡献、控制与责任脱节。

吉林省社会保险事业管理局2026年6月30日的实施说明显示，该省自7月1日起启动新就业形态从业人员职业伤害保障试点。此文件直接支持省级实施事实与所列保障范围；本次复核的人社部全国发布页返回验证页面，不能据此补充全国实际覆盖结论。地区启动、平台纳入、个人登记和待遇兑现需要分别核查，职业伤害保障也不等同于完整劳动关系认定与养老、医疗、失业保障。\cite{InjuryJilin2026}

% ARGUMENT P196
\phantomsection\label{arg:P196}
将制度扩围与实际保障分开，可以沿一项事故的处理链观察：从业者是否进入适用范围，任务是否有记录，事故能否确认，损害如何评定，待遇何时支付。任何环节出现缺口，都可能使名义覆盖无法完全转化为家庭缓冲。省级实施文件能够证明制度在该地区启动，却不能替代全国平台和个人覆盖的统计。

吉林公开文件将出行、即时配送和同城货运列入其试点行业，并说明医疗、伤残和死亡等保障范围及经办流程。这提供了核对制度机制的具体材料；本书没有取得全国统一理赔数据库，因而不报告全国实际受益人数，也不把省级启动推成每名平台劳动者已经全面参保。\cite{InjuryJilin2026}

第八章的社会位置在此体现为不同的信息和申诉能力。专职接单者、跨平台工作者与偶尔接单者，对记录连续性和平台配合的依赖可能不同。评价制度应追踪事故后的最终家庭损失，特别是停工期间收入、必要照护和未获补偿支出，而不只记录是否完成登记。

职业伤害保障适合针对与工作任务相关的风险，但其他收入波动和长期养老问题仍需不同安排。把全部保障责任按传统固定雇主方式分配，可能遇到多平台工作与不连续任务的困难；完全让个人承担，则可能造成保障缺口。可以比较按任务缴费、跨平台累计和公共基础保障的组合。

% ARGUMENT P197
\phantomsection\label{arg:P197}
工作伤害与一般收入波动需要不同的保险逻辑。前者可以围绕任务风险、事故确认和责任形成制度；后者还包括需求减少、工作选择和宏观变化，不能完全沿用同一认定方式。养老则需要较长时间积累与持续筹资。将这些问题统称为“给灵活就业上保险”，容易掩盖具体设计差异。

多平台工作需要可衔接的权益记录，但可携带不意味着无限收集个人信息。记录应围绕缴费、任务与待遇所需的最小信息，明确更正和查询渠道；平台之间结算也需要防止重复缴费、遗漏和责任互相推诿。制度越复杂，越需要考虑小平台和劳动者的实际操作能力。

基础保障与工作相关责任可以互补。公共制度保障基本生活，平台或相关经营主体承担能够控制的任务风险，个人保留合理选择空间。具体比例涉及筹资与价值取舍，本书不提出未经测算的全国统一费率。可检验的目标是保障中断减少、理赔更可靠，而不是只把更多费用形式上列到某一个账户。

政策成本可能由平台利润、商家收费、消费者价格或劳动者报酬共同承担。法律上规定谁缴费，并不能完全决定经济上谁承担。改革评估需要追踪这些调整，同时检查保障是否真正减少事故后的家庭损失。

% ARGUMENT P198
\phantomsection\label{arg:P198}
法定缴费者与最终承担者不同，可以通过价格调整理解。若平台具有较高利润且竞争约束使提价困难，新增费用可能更多减少利润；若商家依赖平台而消费者价格敏感，平台可能尝试向商家转嫁；若劳动供给缺少替代，报酬也可能受到影响。这些是不同条件下的路径，不能预先断言其中一种必然发生。

要判断净效果，应同时比较事故损失和正常时期收入。劳动者每单净收入略有下降，但严重伤害后的损失显著减少，是否属于改善，需要结合风险分布、偏好和基本保障标准；若收费增加而理赔困难，制度收益则可能没有兑现。平均收入不能代替风险保障的价值，参保人数也不能代替实际支付。

平台还可能通过安全设计减少总风险，例如合理时限、道路条件提示和纠纷处理。如果制度只在事故后补偿，却不改变能够预防的危险安排，就会错过减少真实成本的机会。第十二章将把费用归属与控制能力联系起来，使能够影响风险的人承担相应改进责任。

\section{对数字资本的可检验命题}
本章提出的条件性命题是：当平台控制不可替代的入口、退出成本高且争议处理不透明时，其效率收益更可能伴随向商家和劳动者转移风险。若互操作和申诉制度改善，控制造成的不利结果应有所缓解。竞争性解释包括服务质量提高、欺诈风险变化和市场整体需求。

% ARGUMENT P199
\phantomsection\label{arg:P199}
这一候选命题的关键变量需要可观察定义。入口不可替代，可以用转向其他渠道后的客户保留与收入变化研究；退出成本，可以记录迁移评价、数据、设备和客户关系的时间与费用；争议处理透明度，则可观察理由说明、处理期限和纠错结果。若变量始终停留在抽象词语，命题就无法被检验。

风险转移也要具体化。可以追踪结算延迟、规则变动造成的沉没投入、事故损失和收入波动。平台费用上升不必然等于风险转移，劳动者收入下降也可能由总体需求减少造成。需要结合规则调整时点、受影响程度和相近未受影响活动，尽可能排除共同冲击。

即使找到相关关系，也应检查改革是否针对已经恶化的市场实施。若困难平台更早被监管，简单比较监管前后可能把原有问题误算为政策后果。明确这种反向选择，能够防止理论在看到任何结果时都声称自己正确。

这可以通过具体规则调整前后的抽成、商家退出、工资波动和纠纷结果来检验，但需要避免把平台自主改革和市场变化混为一谈。如果在高控制条件下仍存在充分进入和有效协商，收益也广泛传递，就应缩小命题的适用范围。

% ARGUMENT P200
\phantomsection\label{arg:P200}
一个合适的比较设计，不应只追踪仍留在平台的商家。规则变化后退出的人可能承担了最大成本，也可能转到更好渠道；遗漏他们，会使剩余样本偏向有能力适应者。应结合进入、退出和持续经营三组观察，尽可能追踪离开后的收入与经营状态。

反证也应预先写清。如果退出成本高却仍有有效协商，平台规则可被纠正，效率收益通过稳定报酬和较低价格广泛传递，那么控制并未按本书所述路径转成不利后果。此时应进一步识别起保护作用的组织和制度，而不是把良好结果称为暂时例外而排除。

相反，如果互操作改善以后费用没有下降，也不能立即否定其全部作用。可能存在质量提升、迁移尚未完成或新进入需要时间。研究应设置合理观察期与中间指标，同时允许到期以后得出效果有限的结论。可反驳性要求对解释的边界真正作出承诺。

人工智能并未取消生产、控制、分配和风险四个问题，而是改变它们的表现。未来研究不应把技术神化为自动实现公平的力量，也不应把所有创新预先解释为控制工具。要追踪的是技术如何嵌入具体组织，以及谁能影响收益的去向。

% ARGUMENT P201
\phantomsection\label{arg:P201}
数字化把部分资源控制从厂房和机器延伸到接口、评价、数据和模型，但这些权利仍然需要合同、组织和公共制度支持。技术没有脱离社会关系独立运行。它可以降低许多人的进入成本，也可以使少数入口更难替代，具体方向取决于采用方式和竞争条件。

因此，本章的理论修订是把数字收益拆解为生产改进、组织服务、排他收益与风险承担，再观察它们如何分配。这样既能够识别新技术创造的真实能力，也能够追踪被转移到商家、劳动者和家庭的时间、设备与保障成本。分析的单位从单个平台利润扩展到完整交易关系。

制造业并不处于这套关系之外。工厂通过平台采购和销售，供应链使用数据协调，工程人员使用智能工具，产品交付又连接全球需求。下一章将把入口与组织问题放回实体生产和跨境联系，说明技术能力、贸易总额与国内各阶层实际所得之间为何仍有距离。


\chapter{制造业与全球联系}
\section{生产能力是一套协作系统}
制造业竞争力不仅是单厂工资低或设备先进，还包括供应商、技能、物流、能源、金融、工程经验与需求反馈。一个零部件设计的改进，需要上游材料、模具、检测和下游使用者共同配合。产业集群可以缩短试错过程，也可能把风险集中在相近地区和相似产品上。

% ARGUMENT P202
\phantomsection\label{arg:P202}
协作系统的意义，在于一家企业无法单独拥有完成产品所需的全部知识。材料批次变化需要工艺调整，设计改变需要模具与检测配合，客户反馈需要沿供应链传回。距离、沟通和共同经验影响这些环节的速度。集群可以使试错更快，但真正发挥作用的是关系与能力，而不是企业在地图上靠得近。

这种共同能力具有部分外溢。一家企业培训的工人可能到另一家工作，某项检测设施可以服务多个生产者，供应商改进工艺可以惠及不同客户。单个企业未必能够收回全部投入收益，因此公共基础设施和协作组织可能有作用；但公共支持是否必要，还要比较已有市场供给和具体缺口，不能从存在外溢推导所有投资都应补贴。

集聚也会集中风险。同一地区依赖相似客户、能源或运输节点时，一次冲击可能同时影响很多企业；工人虽有多个雇主可选，却可能发现它们同步减少订单。地域上的企业数量不等于真正独立的就业选择。对产业工人与小供应商而言，客户和产品的多样性有时比表面规模更能提供缓冲。

因此，不能把中国制造业概括为纯粹低工资生产，也不能从少数先进产品推断所有行业都已完成升级。技术能力、产品质量、服务网络和单位成本应分别观察。高技术分类是一种统计分类，不等于每个被归入其中的企业都具有高利润或自主关键技术。

% ARGUMENT P203
\phantomsection\label{arg:P203}
“升级”至少包含产品质量、生产效率、任务控制和收益分配几个不同问题。企业可以买入先进设备提高良品率，却仍缺少产品设计与客户渠道；也可以拥有某项独特工艺，在统计上属于传统行业。把行业标签当作技术能力，会忽略企业内部的真实差异。

利润也受竞争影响。技术进步普遍扩散以后，产品价格可能下降，消费者受益而企业利润率未必提高；某些低技术环节因为位置或渠道稀缺，反而可能取得较高回报。因此，利润高不能单独证明技术先进，利润低也不能直接证明没有生产能力。

对劳动者，升级是否改善生活需要观察技能、工资、工时和职业路径。自动化可以减少危险和重复任务，也可能取消初级岗位并提高剩余岗位强度；新设备需要维护人员，但被替代者未必能立即取得这些岗位。企业能力升级与同一批劳动者处境改善之间，需要培训、转岗和保障作连接。没有这些条件，平均生产率上升仍可能伴随局部长期损失。

2026年1至8月，规模以上工业增加值实际增长5.3\%，社会消费品零售总额名义增长1.1\%；房地产开发投资下降19.9\%。这些指标范围不同，不能用增速简单相减得出“过剩量”，但它们提示生产、消费和地产投资需要联合分析。消费品零售总额也不覆盖全部居民服务消费。\cite{Current2026}

% ARGUMENT P204
\phantomsection\label{arg:P204}
不同增长率不能直接相减，是因为它们统计对象与价格处理不同。工业增加值衡量特定工业范围的生产，零售总额记录特定商品交易，房地产开发投资反映一种投资活动。把5.3\%减去1.1\%所得差值称为供需缺口，相当于将不同分母和不同价格口径强行放在同一条尺上。\cite{Current2026}

研究供需关系需要更完整的核算。工业产品可以进入居民消费、企业投资、政府采购、库存和出口；零售商品也可以来自进口，居民还购买大量服务。只有把产出用途与价格变化对应起来，才能判断哪些行业出现库存积压或有效需求不足。总体增速的分化提供线索，却不是对过剩规模的直接测量。

阶层后果也需沿具体行业展开。房地产投资减少可能影响建筑工人和材料供应商，出口扩张可能改善部分制造业岗位，服务消费变化可能影响本地小经营者。同一家庭如果成员分处不同产业，收入变化还可能相互抵消或共同恶化。宏观资料适合描述环境，家庭处境则需要进一步观察其就业组合与负债。

\section{出口总额与国内所得之间}
出口包含国内增加值，也可能包含进口零部件和原材料。产品在多国之间反复加工，毛贸易额会多次记录中间品价值。评价出口对大陆劳动者和财政的贡献，需要国内增加值、工资、税收与利润归属；评价技术依赖则需要看关键投入来源和替代成本。

% ARGUMENT P205
\phantomsection\label{arg:P205}
用一个纯粹假设的两国加工例子说明：境外零部件价值40，进口后在国内加工增加价值60，成品以100出口。出口毛额为100，国内增加值在这个简化例子中是60；不能把100全部解释为国内新创造的收入。如果零部件已经包含此前从国内出口的投入，再简单按进口价值扣减也可能漏掉返回的国内增加值，需要完整追踪生产链。

Koopman、Wang与Wei的框架正是把毛出口分解为不同来源的增加值和重复计算部分，连接贸易统计与国民核算。本文借用这种核算区分，不声称已经利用最新投入产出表计算出2026年中国出口国内增加值率。缺少相应资料时，不能给所有产品统一乘上一个来自旧年份的比例。\cite{ValueTradeNBER}

国内增加值也不等于国内劳动者工资。它还包含折旧、经营盈余、混合收入和相关税费；企业所有权与跨境收益汇出又会影响国民收入归属。要从出口追踪到居民生活，至少还须经过要素收入、税收转移和企业分配。增加值核算避免重复计算，但并未独自回答分配问题。

海关数据经新华社10月14日发布，2026年前三季度货物出口23.20万亿元、进口16.74万亿元，合计39.94万亿元；相应同比增速为16.4\%和22.3\%。这是人民币名义货物贸易，本文据此计算货物差额6.46万亿元。它不是经常账户余额，不包括服务和初次、二次收入，也不是对GDP增长的净出口贡献。\cite{Trade2026}

% ARGUMENT P206
\phantomsection\label{arg:P206}
货物差额是一定时期货物出口与进口名义金额的差。它可以帮助观察对外商品交易，却不能直接说明这一时期总储蓄、全部国际收支或福利得失。经常账户还包括服务和收入等项目，GDP中的净出口则要按照国民核算范围及价格处理计算。不同表之间可以建立联系，不能只因都含“出口减进口”便当作同一个数字。\cite{Trade2026}

人民币计价还受到汇率换算与合同价格影响。相同外币收入换成人民币可能变化，企业进口投入和外币债务也可能同时变化。只看到出口金额增长而忽略成本，会高估利润改善；只看到进口增加而视为国内需求流失，又可能遗漏进口设备和材料支持的生产。

这组公开数据的用途是描述前三季度货物贸易规模。本书保留新华社发布与海关数据之间的来源关系，不将新闻发布当作另一个独立统计样本。计算6.46万亿元是对已列公开数的减法复算，不是新的调查结果，也不用于替代经常账户或增长贡献数据。

\Chart{trade}{2026年前三季度大陆货物贸易}{数据图；万亿元人民币，名义值。进出口为海关数据的新华社发布值；差额为本文计算。三柱不可相加，差额不等于经常账户余额。\cite{Trade2026}}

出口增长可能反映数量增加、价格、汇率、产品升级或提前交货。若贸易政策变化使客户提前采购，短期强劲未必代表持续需求；若出口品质量提高，相同数量也可能取得更多收入。这些都是竞争性解释，本书没有取得足够细分数据来估计各自贡献。

% ARGUMENT P207
\phantomsection\label{arg:P207}
数量、价格与结构需要逐层分解。某类产品出口金额上升，可能是销量增加，也可能是同类产品价格提高，还可能是高价型号在组合中占比增加。若用总额除以总件数所得平均单价衡量质量，就可能把产品组合变化当成同款升级。细分产品、目的地与企业资料可以减少这种混淆。

提前交货则涉及时间重新安排。客户为了应对可能的运输或政策变化先建立库存，会提高当期出口，但随后采购可能减少。确认这一机制需要订单、交货周期与海外库存等资料，不能仅凭出口突然增长就断言抢出口。它只是需要与竞争力提升共同检验的解释之一。

企业内部分配还可能滞后。订单增长首先增加加班和流动资金需求，稳定工资或新增招聘可能以后才调整；若企业认为需求短暂，就可能采用临时用工。不同劳动合同因此承担不同风险。把出口增长当月直接推为所有制造业工人收入同幅增长，会跳过经营预期与用工选择这两个环节。

\section{国际分工既能创造收益，也能形成依赖}
进入全球价值链可以扩大市场、学习标准、获得技术和摊薄固定成本。企业也可能受制于品牌、关键设备、软件许可或终端渠道。谁能决定设计、标准、交期和采购替代，影响价值链中的议价能力。国家间贸易差额不能直接揭示这些企业间权力关系。

% ARGUMENT P208
\phantomsection\label{arg:P208}
价值链中的控制权常体现在谁能改变产品规格、谁可以替换供应商、谁掌握客户和售后。供应商为某一买方购买专用设备以后，转接其他订单可能需要较大成本；买方若能轻易替换供应商，议价关系便可能不对称。相反，掌握难以替代工艺的供应者也可能对大型买方有较强议价能力。企业规模与链条位置需要共同分析。

进入国际生产网络可以带来学习，因为客户要求质量、交期和可追溯性，供应商在反复交付中积累能力。但学习不是自动赠予：培训、认证和试错需要资金，也可能由供应商承担失败风险。能否把经验转成独立设计与多客户经营，决定依赖是否随时间降低。

国家贸易余额无法替代这种企业分析。一国可以整体取得货物顺差，同时其中一些工厂承受很薄利润；另一国的消费者、品牌商与当地劳动者也可能面对不同结果。国际分配需要区分国家账户与企业、家庭账户，避免将一个整体国家拟人化为只有一种收入和一种意志的行动者。

经典关于世界市场与资本积累的研究，提醒人们关注跨境分配和不平等关系；当代分析还需要区分产品、任务和企业。不能把出口国整体视为被动被剥夺者，也不能把贸易双方签约当作所有相关人都受益的证明。消费者、股东、劳动者和受到环境影响的居民可能有不同结果。

% ARGUMENT P209
\phantomsection\label{arg:P209}
跨境交易创造总收益，与所有参与者都改善，是两个不同命题。更便宜的投入可以降低下游成本，但受到竞争的原供应者和工人可能需要转型；更大的出口市场可以提高国内企业规模，却也可能增加特定地区资源与环境压力。总量收益不能直接抵消集中且持久的损失。

Autor、Dorn与Hanson关于美国贸易调整的研究提醒我们，地区劳动市场的适应可能比简单模型假定的更慢，消费者收益与局部就业成本可以同时存在。本文采用这一比较机制，不将美国估计的幅度和持续年限移植为中国预测，更不从该研究推断当前每项贸易冲击的方向。\cite{TradeAdjust2016}

对于中国，关键问题是受冲击者能否迁移、转岗或取得生活支持。住房债务、技能专用性、照护责任和地区服务会影响调整速度，这又把贸易分析带回第八章。一个行业的新岗位增长，未必发生在原失业者所在地区，也未必要求相同技能。补偿与培训的设计应当对准这些实际距离。

一条价值链中，生产环节的低利润率可能来自竞争激烈，也可能与品牌和渠道集中有关。要证明租金转移，需要比较进入壁垒、知识产权、合同和风险承担。不能单凭最终售价与出厂价相差很大就推断全部差额是剥削，因为运输、库存、零售服务、营销和失败产品也有成本。

% ARGUMENT P210
\phantomsection\label{arg:P210}
最终售价与出厂价之间至少经过物流、库存、营销、零售、售后和税费。某些费用是真实服务，某些可能包含垄断租金，还有一些用于承担滞销与产品失败。仅将两端价格相减，再把差额全部归为对制造环节的剥夺，会把不同经济活动合并。

可以尝试比较同类产品在不同渠道中的费用和利润，并观察进入障碍改变后价格差怎样变化。若渠道开放后服务质量保持而价差明显下降，租金解释更有支持；若价差主要随库存风险、退货率和服务内容变化，成本解释可能更强。两者混合存在时，应分别估计，不必强迫结论只有一种来源。

对制造企业与劳动者，增强工艺和客户能力可能改善议价，但并不保证总价差都转成工资。控制新的品牌和渠道也可能使企业内部权力重新集中。产业升级政策若以居民生活为目标，还需跟踪工资、培训、利润分配与公共收益；品牌归属变化本身不是全部验收标准。

\section{技术自主需要定义具体能力}
“自主”可以指能够设计、制造、维护、替代或在外部供应中断时维持基本功能。这些目标的成本不同。某种设备可以在国内组装，却仍依赖关键软件；某种材料能够生产，却可能在质量稳定性和规模上存在差距。政策应写明要获得的能力，而不是只提高本地采购比例。

% ARGUMENT P211
\phantomsection\label{arg:P211}
能力分解能够使政策目标可检验。会设计图纸不等于能以稳定良率批量生产，能生产不等于能维护整个使用周期，能维护不等于关键材料中断后仍能运行。每个环节需要不同人员、设备和知识。把它们统一写为“已经自主”，会使实际薄弱点从指标中消失。

成本比较也应包括时间。正常时期采购国际产品可能成本较低，建立替代能力则可能需要持续研发和小规模试产；这种投入提供的是应对中断的选择权，不能只按当期销量评价。但若没有明确风险场景、服务对象和性能要求，备用能力也可能成为维持无效项目的理由。

对专业人员与产业工人，技术能力需要职业积累。只购买设备而不建立工艺训练、维护和知识传承，可能形成对外部服务的新依赖。企业为了短期节约培训成本，也可能削弱长期能力。因此，“自主”的社会条件还包括稳定技能体系与合理职业回报，而不是只有设备产地。

供应多元化、库存、替代技术和国内研发是不同风险管理方式。为一切投入追求完全本地化，可能提高成本并失去国际协作；对关键不可替代环节完全不准备，则可能在冲击下损失更大。合理组合取决于中断概率、社会后果和替代时间，而不是一条适用于所有产品的比例。

% ARGUMENT P212
\phantomsection\label{arg:P212}
库存、替代供应与自主研发应按冲击持续时间比较。短暂中断可以由库存缓冲，较长中断需要可替换供应或工艺，根本性技术封锁才更可能要求长期研发与产品重新设计。三种工具对应不同时间尺度，也有不同机会成本。库存占用资金并可能过时，多供应商增加认证与协调成本，研发可能失败。

一个可用的决策框架，是比较准备成本与不同冲击下的损失，而不是假定能够准确知道所有概率。对于后果严重、恢复时间长的关键功能，可以进行情景和压力测试；对于容易替代的普通投入，则不必使用相同强度的本地化要求。韧性是保持必要功能的能力，不是任何时候都不使用外部资源。

还要检查风险是否被分散。表面上有三家供应商，如果它们依赖同一种上游材料或同一个运输节点，实际风险仍可能高度相关。企业和公共部门都需要沿供应链向上核查。小企业通常较难独立完成这种工作，公共信息和共性检测可能比普遍补贴新增产能更有针对性。

公共支持适合具有学习外溢、公共基础或重大风险保障价值的领域，但需要避免补贴被规模指标俘获。企业承诺产能不等于形成有效能力，专利数量也不自动说明技术质量。可采用阶段性质量测试、真实客户采用、维护能力和退出规则，降低永久补贴的风险。

% ARGUMENT P213
\phantomsection\label{arg:P213}
公共支持的经济理由，应当与具体市场缺口相对应。基础研究成果容易外溢，企业可能投资不足；共性检测设施具有共享价值，分散重复建设可能低效；关键功能的备用能力可能服务更广泛的社会需要。这些理由不同，支持工具也应不同，不能用同一种按产能补助覆盖全部问题。

阶段性考核可以把资金与学习结果联系起来。研发初期关注可验证的技术问题，中期关注性能和稳定性，接近应用时关注真实客户采用和维护。失败并不自动意味着浪费，因为探索本来具有不确定性；但长期没有学习进展而持续追加支持，则需要重新评估。接受合理失败与拒绝无条件续期可以同时成立。

成本承担者也应清楚。财政资金来自其他可选用途，政策性融资仍有偿付或损失，供应商参与试制也会投入资源。支持项目取得收益以后，公共部门可以通过预先约定的股权、回收或公共服务义务分享部分结果，但安排应避免抑制必要创新。最重要的是让投入、结果和责任可以相互核对。

\section{竞争、产能与地方重复投资}
产能超过当期需求并不总是错误。设备建设有时间差，需要为未来需求和备用安排留空间。某个行业同时扩张，也可能是企业对增长机会作出的分散判断。问题在于需求预期是否一致偏乐观，融资和地方补贴是否使低回报项目持续进入，以及退出机制是否有效。

% ARGUMENT P214
\phantomsection\label{arg:P214}
正常备用与低效过剩的区别，需要结合需求和技术寿命。医院设备、能源系统与一般消费品工厂的合理利用率不会相同；某些设施需要为峰值与故障留余量，某些生产线则可以通过灵活排班调整。统一以“没有满负荷”判定浪费，会忽视服务可靠性的价值。

判断问题应观察预期如何形成、亏损如何处理。如果很多企业根据同一乐观需求扩张，又相信损失最终可以被外部承担，投资可能超过社会需要。若亏损项目能够及时退出、设备转用和人员转岗，调整成本可能较小；若债务、地方目标和专用资产使退出困难，低回报可能持续。

过剩还可能是相对价格问题。设备存在不等于产品没有使用价值，而可能是按当前成本无法卖到覆盖投资的价格。降价能否扩大需求、工艺能否降低成本、资产能否转用，决定损失大小。政策应比较这些路径，而不是把继续生产或立即关闭当作唯一答案。

地方政府获得就业与税源的回报，可能积极支持产业集群；若邻近地区使用相似招商模板，就可能形成重复投入。全国层面看似先进的产业方向，在具体区域也可能缺少人才、供应链和客户。技术类别不能替代项目可行性。

% ARGUMENT P215
\phantomsection\label{arg:P215}
地方支持集群具有可以理解的动机：企业落地可能带来岗位、采购和未来收入，公共基础设施也可以形成区域能力。但同一地方未必承担全部跨区域重复建设后果，企业也可能在地区之间比较优惠。若信息不足与竞争性补贴相互加强，单个地方看似合理的选择可能形成总体低效。

项目评价因而需要比较本地真正具有的条件。工程人才、供应商、客户距离和维护能力不能由行业名称代替。一座新厂房可以很快建成，熟练协作网络则往往需要较长积累；若先形成大量固定资产再等待生态出现，预算和债务就承担了时间风险。

对地区居民，判断招商成效需要看岗位持续性、工资与公共服务，而不只看承诺投资。外来企业获得优惠以后是否形成稳定采购，地方债务是否挤压维护和教育，也影响实际结果。对小供应商与劳动者的支付责任应进入项目考核，使地方收益与承担的社会成本相联系。

价格下降在这里有双重含义。它可以表示效率提高并让消费者受益，也可以表示竞争把融资和环境成本转向其他主体。如果工资拖欠、供应商账期和公共补贴增加，产品低价的社会成本就需要重新计算。不能把所有降价都叫作恶性竞争，也不能把低价本身当作效率证明。

% ARGUMENT P216
\phantomsection\label{arg:P216}
低价竞争可以推动效率，也可能把未支付成本隐藏起来。若企业减少材料浪费、提高良品率和缩短物流，降价有真实资源基础；若通过拖欠工资、延长供应商账期或减少必要环保投入降低现金支出，低价的一部分由其他主体承担。两者可能同时发生，需要分项核对。

价格战中的现金流尤其重要。企业暂时以低价扩大市场，可能利用规模分摊固定成本，也可能依赖持续融资。是否可持续，需要比较每单位贡献、固定成本、客户保留与资金来源。销量增加但应收账款更快增加，未必改善支付能力。

劳动者与消费者在这里同样交叉。家庭购买低价商品可能受益，同时作为雇员或小经营者承受付款和收入压力。政策若直接限制降价，可能保护低效生产者并损害消费者；若只追求最低价格，又可能纵容拖欠与外部成本。更有针对性的方向是保障支付、安全、环境与公平进入，让价格竞争在较完整的成本约束下进行。

\section{外需不能替代所有内部调整}
出口可以吸收国内生产，带来收入和学习，但国外需求有其规模、周期和政策约束。一个经济体的外部盈余对应其他经济体的相应余额，不可能所有经济体同时依靠越来越大的净出口解决内需不足。这里的约束来自国际核算与相互依赖，不是反对贸易。

% ARGUMENT P217
\phantomsection\label{arg:P217}
国际核算约束的含义，是全球各经济体对外净余额需要相互对应，不是每一对贸易伙伴之间都必须平衡。一个地区可以从另一个地区进口投入，再向第三地销售最终产品。双边差额不能完整反映生产链和资金流，也不能单独说明谁在整体交易中取得更多福利。

外需可以扩大市场，但无法保证无限吸收新增供给。海外家庭收入、企业投资、贸易成本和政策都会变化。即使单个企业可以凭质量扩大份额，整个经济体也需要面对外部市场规模与其他供给者的反应。将企业层面的成功经验直接推成全国持续依赖净出口的方案，会忽略加总后的约束。

因此，改善国内家庭收入与保障有双重意义：首先改善生活，其次使生产拥有更稳定的需求来源。它不要求关闭外部市场，也不保证所有行业都能转向内销。出口产品的规格、渠道和价格可能与国内需求不同，转向需要实际适配。宏观再平衡应承认这些企业层面的成本。

对中国大陆而言，扩大出口与改善家庭消费能力可以互相支持。企业从全球市场获得规模和知识，劳动者与公共部门获得收入后形成稳定内需；如果收益主要留在投资和资产积累中，而家庭保障与收入传递较弱，这个循环就可能偏向持续扩张供给。是否如此，需要收入分配和企业再投资数据验证。

% ARGUMENT P218
\phantomsection\label{arg:P218}
出口收益进入家庭，可以经过工资、经营收入、公共税收与资产分配等通道。企业保留利润用于有效投资，也可能支持未来岗位；但若家庭长期收入增长较弱、保障不足而投资持续扩大，终端需求可能跟不上新增能力。这里需要比较时间路径，不能要求每一笔投资在当年全部转成消费。

分配还影响需求结构。较高收入家庭可能增加不同类型的服务和资产配置，低收入家庭可能更优先满足基本消费；企业则购买设备和中间投入。相同总收入增长在不同主体间分配，会改变产品需求，进而影响哪些行业值得扩张。政治经济学的分配问题因此与生产结构直接相连。

这一机制可以被其他力量抵消。公共投资、外需和新产品需求可能支持增长，家庭也可能基于预期自愿储蓄。若要判断收入传递是否构成约束，需要结合消费、储蓄、债务和服务资料。观察到消费偏弱本身，不足以证明只有分配一种原因。

进口也不应只被当作从本国产出中扣除的坏事。它提供原料、设备、技术与消费选择，能够提高国内生产效率。净出口对GDP的核算贡献不代表每一笔进口的福利损失。研究全球联系，应同时看生产能力、居民所得、可用商品与风险。

% ARGUMENT P219
\phantomsection\label{arg:P219}
GDP支出法中扣除进口，是为了从包含进口成分的消费、投资和政府支出中排除境外生产，避免将其计入国内产出；它不是给进口贴上负面福利标签。进口一台提高效率的设备，可以在当期被扣除其境外生产价值，同时支持以后更多国内生产。两件事并不矛盾。

进口消费品也可能扩大选择并降低生活成本，但受竞争的国内企业和劳动者可能需要调整。因此，福利比较需要看消费者收益、生产者变化、就业调整与财政成本，不能仅用净出口贡献正负替代。相同进口增加，对不同阶层可能具有不同净效果。

对于技术依赖，关键是可替代性而非进口总额。大量可替代普通投入与少量难以替代关键设备，形成的风险不同。将进口减少一概称为自主能力提高，可能掩盖生产收缩；将进口增加一概称为依赖加深，又可能遗漏多元化和设备更新。应回到具体能力与供应条件。

\section{贸易冲击由谁承担}
关税、运输中断或技术限制的经济负担，可能由海外客户、国内出口商、员工、供应商和财政共同承担。法定缴税者只是付款环节，价格与合同调整决定最终分担。汇率变化也会同时影响出口收入、进口成本和外币负债，不能只给出单向利好判断。

% ARGUMENT P220
\phantomsection\label{arg:P220}
关税或运输费用增加以后，合同双方会重新分担成本。海外需求较不敏感、国内产品难以替代时，出口商可能保留价格；替代竞争较强时，则可能让利。国内企业还可能调整工资、采购和利润。法定税款在哪个边境由谁支付，并不能穷尽经济负担。

汇率变化也有双向作用。出口收入以外币计价、成本主要为人民币的企业，可能与大量进口材料或持有外币债务的企业受到不同影响。套期保值和合同期限还会改变传导速度。不能只依据汇率一个方向变化，就断言整个制造业同步受益或受损。

风险在供应链中的落点，取决于谁可以修改订单和支付期限。大买方取消采购可能把库存留给小供应商，小供应商又可能减少工时或拖欠劳动报酬。需要沿着合同追踪最终承担者，并区分合理风险分担与单方面违约。单个企业利润保持稳定，可能是效率改善，也可能是调整成本已经传到别处。

如果订单突然下降，企业可以减少利润、缩短工时、裁员、寻找新市场或延迟付款；不同选择把风险移向不同主体。公共支持可以帮助有竞争力的企业转型，却需要区分临时冲击和长期需求消失。保留一切现有产能，可能阻碍资源进入更有价值的用途。

% ARGUMENT P221
\phantomsection\label{arg:P221}
订单冲击首先是现金流问题，随后才可能变成资产和就业问题。企业原有工资、租金和债务按期到来，但销售回款减少；若有储备和可持续订单，可以跨过短期缺口，若市场永久缩小，则继续借款可能扩大最终损失。公共支持需要判断这一区别，但判断只能基于可核查的经营信息，不能保证完全准确。

对工人，缩短工时与裁员产生不同后果；对供应商，延迟付款可能比订单减少更危险，因为前者涉及已经垫付的成本。对家庭，技能培训如果没有生活支持和可达岗位，也未必可行。将调整政策只设计为给企业贷款，会遗漏这些成本承担者。

可以把支持与转型计划、支付纪律和信息披露相联系，同时为没有持续前景的业务设计退出。保护必要生活和可转换能力，并不要求保护全部股东历史收益。若支持长期维持无法销售的产能，资源可能被占用；若停止过急，又可能损坏原本可恢复的供应网络。政策应根据订单、净现金流和替代用途动态调整。

本章的中国经验不能压缩为“外向”或“内向”二选一。更有解释力的问题是，国际协作创造了哪些真实能力，这些能力的收益如何进入大陆家庭生活，以及外部冲击到来时是否有可信的分担机制。开放与风险管理是需要共同设计的制度问题。

% ARGUMENT P222
\phantomsection\label{arg:P222}
本章将开放理解为生产协作关系，将韧性理解为冲击下维持必要功能的能力。两者可以互补：多样化国际供应和客户可以降低单一依赖，国内公共服务与技能体系可以支持转换。把开放与自主设成绝对对立，会使实际可行的组合从讨论中消失。

对各阶层而言，贸易收益需要经过不同通道兑现。资本所有者关注利润与市场，专业人员关注技术与职业路径，产业工人关注工资和工作连续性，小供应商关注回款，居民关注价格与公共资源。统一的贸易总额不能替代这些结果，也不能由某一群体获益推断全部关系改善。

由此可以进入危机问题：如果生产能力、融资、市场需求与家庭再生产在时间和结构上错开，局部经营困难可能通过这些关系相互放大。下一章将把前面各章的传导机制联结起来，同时区分核算约束、行为假设与已经观察到的事实。

\chapter{积累、失衡与危机}
\section{积累能够扩大能力，也可能积累错误}
资本积累是把当期资源用于未来生产。它包括设备、基础设施、知识和组织能力，也可能表现为对既有资产的购买。后者改变所有者，不必然增加全社会生产能力。讨论投资不能把购置旧资产、形成新资产和资产重估混为一谈。

% ARGUMENT P223
\phantomsection\label{arg:P223}
购买一台新机器与购买一家已有企业的股票，可以都被日常语言称为投资，但对社会生产账户的含义不同。机器生产动用了当期劳动和材料，并形成未来服务能力；股票转手主要改变既有收益权的持有人。股票交易可以改善融资与治理条件，间接支持以后生产，但不能把成交额直接加进新增产出。

资产价格上涨还可能在没有新增实物的情况下提高账面财富。它可以通过抵押和预期影响消费投资，却不是同额新增社会资源。如果所有持有人同时试图兑现，价格与买方支付能力又会发生变化。因此，账面资本增值、可实现现金收入和真实生产能力应当分别记录。

积累也包括维护。忽视设备、管网、技能和环境修复，短期可以留下更多可支配现金，长期却可能损耗已有能力。社会的净积累不能只看新建数量，还须扣除损坏与退化。对依赖公共设施和劳动收入的家庭，维护不足可能首先表现为通勤、健康和工作中断，而不立即表现为显眼的资产损失。

积累的收益取决于用途、需求和维护。新建厂房可能扩大产量，老旧管网维修可能减少损耗，基础研究可能很久以后才显示价值。用短期营业收入评价所有投资，会低估部分公共和长期项目；把长期目标作为所有低回报项目的理由，又会使错误难以纠正。

% ARGUMENT P224
\phantomsection\label{arg:P224}
不同投资需要不同观察期。基础研究可能没有短期销售，公共卫生能力可能只在冲击时显现，但这不意味着无法评价。可以观察中间能力、使用对象、替代方案和维护状态，再根据长期目标安排复核。相反，如果一个项目声称创造即时商业回报，就应接受相应现金流检验，不能在表现不佳后无限改称长期公共项目。

时间错配会改变损失分担。项目服务寿命很长而融资期限很短，可能在收益尚未形成时遭遇偿付压力；短期收入看似充足但维护被推迟，又可能将费用留给以后使用者。期限安排不仅是金融技术，还决定哪一代纳税人、居民或劳动者承担成本。

合适的评价因此要同时说明目标、时间、维护和退出。对不可逆投入，前期论证尤其重要；对可以分阶段试验的项目，则应利用小规模学习减少误判。把全部资金一次投入，可能使后续决策被已经发生的成本绑住。理性调整应比较从现在开始的增量收益与成本，而不是只为了证明过去决定正确而继续扩张。

经典积累分析关注竞争如何推动扩张和集中，以及劳动者处境如何随之改变。本文保留这一动态视角，但不把集中和危机视为在所有制度下按同一速度发生的机械规律。竞争政策、公共投资、金融约束和保障制度能够改变路径，其效果也需要观察。\cite{Marx1}

% ARGUMENT P225
\phantomsection\label{arg:P225}
竞争推动积累，是因为企业希望降低成本、扩大市场或避免被淘汰。成功者可能继续扩张，失败者可能退出，资源和控制因而集中。但是，这一过程并非只有一条方向：新技术可以降低进入成本，供应链分工可以支持中小企业，竞争制度也可以限制排他行为。需要研究在什么条件下集中持续，在什么条件下新进入能够打破格局。

集中对劳动者的影响也不单一。较大组织可能提供培训、设备和稳定服务，同时可能减少当地雇主选择；分散经营可能增加进入机会，也可能缺少安全投入和保障。不能以规模大小直接确定劳动条件。岗位替代、工资制度和组织问责才是具体机制。

经典积累分析的当代用途，在于要求解释扩张怎样重塑社会关系，而非只追踪资本数量。本文保留这一问题，把公共权限、金融安排和家庭再生产纳入路径条件。若制度使收益与责任更接近、劳动者更有选择，积累的分配后果就可能改变；这种可改变性本身也应当进入理论。

\section{利润率下降不是无条件定律}
在一个简化关系中，利润率可以表示为利润份额乘以资本生产率：
\[
r=\frac{\Pi}{K}=\frac{\Pi}{Y}\frac{Y}{K}.
\]
若设备相对产出增加而利润份额不变，利润率可能下降；但技术提高设备效率、资本品变便宜、利用率改善或分配改变，都可能抵消这种变化。这里的$K$需要说明是按历史成本还是重置价值，$\Pi$也要说明是否扣除折旧、利息与税收。

% ARGUMENT P226
\phantomsection\label{arg:P226}
上述分解是恒等式，不是对利润率变化原因的识别。设一个假想企业利润为20、产出为100、资本为200，利润率为10\%；利润份额为20\%，资本生产率为0.5，两者相乘仍为10\%。这个例子只是展示同一比率可以被拆成分配与资本使用两个部分，没有证明哪一部分决定另一部分。

若需求减少使设备闲置，产出下降可能压低资本生产率；若设备价格下降，按重置价值计算的资本分母也可能下降；若企业利润包含资产出售或重估，利润分子又可能与经营能力不一致。同样观察到利润率变化，背后可以有不同路径。因此，比较必须先统一资本估值、折旧和利润边界。

财务杠杆还使股东回报与总资本经营回报不同。企业借款以后，利息和权益分母同时变化，股东收益率可以提高而生产效率不变。把上市公司股东回报率称为社会资本利润率，会把融资结构带来的变化混入生产解释。本文只使用明确边界下的算例，不用不同口径拼接趋势。

经典利润率下降趋势论提出了重要问题：企业追求节约劳动的技术，怎样反过来影响盈利和积累条件？本文不从这一问题直接推出必然崩溃。即使某段时期利润率下降，也可能来自需求、能源价格、竞争、工资或资产估值，不能只选择符合理论的解释。\cite{Marx3}

% ARGUMENT P227
\phantomsection\label{arg:P227}
节约劳动技术对盈利的影响，取决于企业能否保留成本优势以及市场怎样回应。先采用者可能取得暂时超额收益，竞争者跟进后价格下降，部分收益转给消费者；若设备和维护成本较高，需求不足又使利用率下降，投资回报可能不及预期。这里存在真实动态张力，却没有一个无需条件的单向结论。

利润率下降也不自动意味着利润总额下降。资本规模扩大时，较低比率可能对应更高总额；但较高利润总额也未必保障支付，因为应收款、债务到期和再投资安排会改变现金。危机分析因此要把比率、总额与现金流分开，不能选择其中一个替代其他两个。

若要检验某一利润率理论，应预先规定资本和利润定义、行业范围以及预期传导。例如它是否预测资本深化以后经营回报持续下降，在控制利用率、产品价格和竞争变化后仍然存在？若只有资产估值解释变化，原机制就需要缩小。可反驳性要求理论允许出现不支持自身的结果。

2026年工业价格的明显分层尤其提示，行业利润变化可能首先来自成本和售价不同步。没有同口径资本存量与净利润序列，本书不画一条声称代表中国长期资本利润率的曲线。少画一条来源不明的线，比给理论配一张看似有说服力的图更有价值。

% ARGUMENT P228
\phantomsection\label{arg:P228}
出厂价格与购进价格的不同变化，会通过投入结构影响行业利润。能源和材料密集企业可能承受较高上游成本，具有定价能力的企业可能部分转嫁，面向价格敏感需求的企业则可能较难转嫁。总PPI与总购进价格指数无法直接给出每家企业的利润变化，因为权重、产品和时间不同。\cite{PPI2026}

劳动成本也不能在成本讨论中自动被视为应当压低的部分。工资既是企业支出，也是家庭收入和需求来源；降低工时或工资可能缓解单家企业现金压力，却通过家庭消费影响其他企业。这个加总问题不意味着企业可以忽略自身支付能力，而是说明宏观协调不能照搬单个主体的节约逻辑。

对2026年的观察，本书将价格分层视为需要继续追踪的事实线索。它支持分别研究投入与售价，却不单独证明全国利润趋势，更不足以验证某种必然危机论。缺少可比序列时，保持判断范围清楚，比把不同月份、不同统计对象的指标拼成完整故事更可靠。

\section{再生产的实现条件}
产品必须找到最终使用者，生产扩张才能转成可持续收入。如果企业不断投资，而家庭消费与公共服务需求不足以吸收产出，外需或继续投资可以暂时支持销售。但未来投资也要面对真实用途，不能靠无限建设维持每一轮收益预期。

% ARGUMENT P229
\phantomsection\label{arg:P229}
再生产需要同时满足实物与支付条件。企业需要投入品、设备和劳动，家庭需要消费品与服务，各部门还需要在适当时间取得支付手段。实物已经生产出来而买方暂时缺少现金，可能形成周转问题；产品根本不符合需求，则可能形成更持久的价值实现问题。两者处理方式不同。

投资需求可以支持生产，但投资的用途仍要回到未来服务。建设一条有用管网即使没有直接销售收入，也可能通过减少损失和改善生活产生回报；反复建设无人使用的设施，则可能只在建设期带来需求，留下维护与债务。判断不能只问支出是否发生，还要问未来谁使用、怎样维护和由谁支付。

马克思第二卷关于资本循环和再生产的讨论提示了部门之间相互实现的要求。本文在此保留协调问题，同时纳入公共服务、银行信用和对外部门，避免把现代经济机械压缩为两个部门的数量比例。核算能够揭示不可能同时成立的账，但具体均衡怎样形成仍需要行为与制度解释。\cite{Marx2}

这不是说消费占比越高一定越好。追赶阶段的基础设施、教育与技术投资能够提高未来生活水平；一些家庭也需要自愿储蓄。问题是投资能否在合理期限内形成有用服务，收入分配能否支持需要的实现，以及风险是否被低估。

% ARGUMENT P230
\phantomsection\label{arg:P230}
高储蓄既可能支持未来投资，也可能反映家庭担忧。区分两者，需要观察储蓄动机、可得保障、年龄结构与投资用途。一个社会为改善基础设施而积累，与家庭因重大支出不确定而压缩基本消费，在总量上都可能表现为消费占比较低，但政策含义不同。

投资同样存在不同期限和社会回报。有用的教育与健康投入可能先增加支出、后提高能力，短期无法用销售回款完整评价；缺乏需求的商业项目则不能仅因被记作投资就视为未来必然有益。将消费与投资设成道德上的高低排序，会妨碍比较具体用途。

更合适的问题是，当期生活需要与未来能力之间是否存在可改进的配置。例如减少低利用率建设，将资源转向可靠照护，可能同时改善当期生活和劳动参与；但转换需要人员与组织，并非把预算数字移动就能立即完成。再平衡必须承认真实资源约束和转换时间。

简单消费不足论与简单供给不足论都可能遗漏结构。有的行业产能过多，托育、养老和优质医疗却仍可能供给不足。资源不能无成本在二者之间转移，技能、地点与融资安排都有约束。调整需要结构变化，而不只是把一个总量按钮调高或调低。

% ARGUMENT P231
\phantomsection\label{arg:P231}
某些商品过剩与某些服务不足可以同时发生，因为资源具有专用性。闲置厂房不能直接变成合格护理人员，建筑工人的技能也不能无成本转成医疗技能。总需求扩张可能先流向已经容易供给的领域，而难以解决长期形成的服务能力缺口。

结构调整因此需要需求、供给和分配配合。向家庭提供支持可以改善支付能力，服务培训和设施改造可以扩大供给，劳动条件改善可以留住从业者。如果只提高支付能力而短期供给固定，价格可能上涨；如果只建设设施而没有可负担需求和人员，利用率可能不足。

阶层差异决定了约束发生在哪里。高收入家庭可能通过市场购买服务，低收入家庭可能用成员时间补足缺口；照护劳动者自身也需要保障与职业发展。政策若只把家庭负担转到低薪服务人员，就没有消除再生产压力。完整协调需要把服务对象和提供者都纳入。

\section{债务动态的边界条件}
以简化的政府债务账户为例，设债务与名义收入之比为$d$，名义有效利率为$i$，名义收入增长率为$g$，当期新增基本赤字占收入之比为$p$。在忽略估值、统计调整与其他存量流量变动的简化情形中，
\[
d_t=\frac{1+i}{1+g}d_{t-1}+p_t.
\]
这解释了为什么债务可持续性与增长、利率和基本收支同时有关。不能把实际增长率与名义利率直接放在同一个式子里，也不能忽略隐性担保兑现和资产处置。

% ARGUMENT P232
\phantomsection\label{arg:P232}
债务比率公式来源于一个简化预算关系：期末债务等于期初债务加利息和不含利息的新增赤字，再除以当期名义收入。分母增长使既有债务相对变小，利息和新增赤字使债务增加。公式本身没有说明政府应该选择多大赤字，也没有说明增长与利率可以被独立控制。

实际账户还需要加入估值与其他存量流量调整。汇率变化可能改变外币债务，担保兑现可能把或有责任变成显性债务，统计合并范围变化也会改变总额。卖出资产用于还债，可以降低毛债务，同时减少未来资产收益；因此，毛债务下降不一定等同于净财政状况同幅改善。

使用有效利率还须注意存量与新增融资的差别。市场新借款利率变化，不会立即作用于全部存量债务；到期结构和合同调整决定传导速度。只用当月利率替代整个债务组合利率，会高估或低估压力。该算式适合澄清变量关系，真实判断仍需具体债务期限与现金流。

以初始债务比率100\%、基本赤字为零的纯示例看，若名义增长5\%、利率3\%，下一期比率约98.1\%；若增长2\%、利率5\%，约102.9\%。这不是2026年中国政府债务预测，也不是对任何真实债务口径的估计，只用于说明相同债务为何在不同条件下产生不同压力。

% ARGUMENT P233
\phantomsection\label{arg:P233}
示例中的增长与利率差异经过多期会累积，但其路径依赖参数不变的假设。名义增长5\%、有效利率3\%时，初始比率经过五期约降到90.8\%；反向组合增长2\%、利率5\%时，约升到115.6\%。附图用同一初始值展示这种复利差异，不能被解读为两个真实政策方案的预测。

如果需要保持债务比率不变，基本收支还要弥补利率与增长带来的差额。在利率高于增长的情形中，可能需要基本盈余；在增长高于利率时，某种程度的基本赤字仍可能与稳定比率相容。但这个代数结果没有告诉我们怎样实现所需收支，也没有保证调整不会影响增长。

例如压缩必要维护能够短期改善基本收支，却可能损害未来收入；提高有效服务支出可能增加短期赤字，却减少长期成本。政策选择必须比较这些反馈。债务管理不能只把一个比率压低，而要保留未来支付与生产能力。

\Chart{debt}{增长与利率改变债务比率的假设路径}{情景图；初始债务／收入为100\%，基本赤字为零，参数固定，忽略存量流量调整。两条曲线由上述恒等式逐期计算；不代表对中国的预测。}

即使增长高于利率，低回报新增项目和或有负债也可能增加压力；即使利率高于增长，可信的财政调整和高质量资产也可能保持偿付。真正的研究需要合并边界与现金流，而不是把一个不等式用作永久乐观或悲观的口号。

% ARGUMENT P234
\phantomsection\label{arg:P234}
偿付压力与总量比率不同。债务规模相同，短期集中到期与长期分散到期面对的现金需求不同；债务以本币还是外币计价，债权人结构和可续借条件也不同。拥有长期高质量资产的主体，仍可能因短期现金缺口出现问题，净资产不足者则可能依靠续借暂时维持支付。

将不同部门合并时，需要抵消内部债权债务，但不能因此忽略机构之间的实际支付约束。公共部门总体拥有资源，不表示每个地方或企业都能即时取得；银行持有公共债权，也不意味着损失不会影响其资本与放贷。总体合并账与分部门现金账应同时使用。

“国家可以协调”是制度能力的描述，不是损失消失的证明。协调可以延长时间、改变承担者并减少恐慌，但最终仍需要真实收入、资产回收或明确让利。将这些来源公开，才能判断风险稳定是否改善了整体结果，还是只是把成本延后并放到不易观察的账户。

\section{三类反馈循环及其竞争解释}
第一类是资产负债表循环。资产价格下降削弱抵押和净值，主体削减投资与消费，进一步影响收入和资产需求。第二类是收入与再生产循环。收入不稳定和照护压力提高预防性储蓄或压缩就业，终端需求与工资形成受到影响。第三类是组织激励循环。扩张收益在本部门兑现，失败成本可以向外转移，新的扩张因而持续。

% ARGUMENT P235
\phantomsection\label{arg:P235}
资产负债表循环的启动，需要资产变化能够影响支出或融资。没有债务、也不准备出售住房的家庭，可能对价格短期变化反应较小；依靠抵押续借的经营者，则可能迅速受到影响。资产价格下降并不自动导致同样幅度消费收缩，关键是负债、抵押约束与收入预期怎样连接。

收入与再生产循环需要区分有能力储蓄与没有能力缓冲的人。收入不确定可能使有余量家庭多留存款，也可能使缺少余量家庭直接减少基本消费或承担高成本借款。两种行为都可能降低当期需求，却对应不同生活损失。以平均储蓄率解释全部家庭，容易把最脆弱者的被迫压缩隐藏起来。

组织激励循环则发生在收益与责任跨越边界时。若管理者在扩张期取得考核收益，而失败成本由以后预算、银行或供应商承担，原有选择可能重复。软预算约束理论讨论的正是预期救助怎样影响事前行为。需要同时观察救助预期、决策权限和实际损失分担，不能仅因某次支持发生便断言所有行为都由救助驱动。\cite{Kornai2003}

三类循环相遇时，地方项目回款不足可能影响银行与供应商，供应商减少工资又影响家庭，家庭支出下降进一步削弱经营。这个链条具有经济上的可能性，但其强度、速度和范围需要数据。机制图的箭头表示待检验传导，不表示每条联系已经取得全国因果估计。

三类循环可以互相加强，也可以被制度切断。公共支出能够稳定收入，银行资本和损失确认可以阻止风险隐藏，社会保障能够降低家庭被迫压缩基本消费的程度。但政策也可能产生新的激励问题，因此需要同时观察短期稳定和长期行为。

% ARGUMENT P236
\phantomsection\label{arg:P236}
切断循环应针对正在约束行为的环节。如果企业有可靠订单却缺少短期周转，及时付款和融资可以避免无谓停工；如果产品长期没有需求，继续贷款只能推迟损失确认；如果家庭因为照护中断退出劳动，仅降低企业借款利率未必改善就业。相同总量刺激，可能因约束不同而效果相差很大。

公共稳定措施能够避免同时收缩，但也需要说明承担方式。补充银行资本可能保护支付与信用功能，同时需要决定原股东和管理责任；地方获得支持可能保住必要服务，同时应约束继续新增无效项目。短期稳定与长期纪律并非只能选择其一，难点在于分别保护功能和处理收益权。

政策还会改变预期。明确、可兑现的基本保障可能减少家庭被迫自保；无条件保护所有投资损失则可能鼓励额外冒险。两种承诺的对象不同。将保护基本生活与保证全部历史回报区分开，有助于既降低灾难性损失，又保留对投资选择的责任。

竞争性解释不能省略。居民减少借款，可能是风险担忧，也可能是人口结构和已完成购房需求；企业投资减少，可能是融资约束，也可能是理性停止低回报项目；出口强劲，可能是竞争力，也可能包含订单时点变化。宏观叙述只有提出可以区分的微观观察，才从解释故事走向检验。

% ARGUMENT P237
\phantomsection\label{arg:P237}
区分竞争解释，需要寻找它们预测不同的观察。如果投资减少主要因为融资约束，融资条件改善后有订单企业应更容易恢复投资；如果主要因为预期回报下降，单纯增加贷款可能效果有限。若两类企业都存在，平均结果就可能掩盖异质性，研究应按订单、杠杆和融资来源分组。

家庭减少借款也有多种原因。已经完成购房的家庭可以正常偿债，担忧未来收入的家庭可能主动减杠杆，无法取得贷款的家庭则可能有未满足融资需求。相同贷款余额下降，不能直接代表相同意愿。需要把申请、批准、偿还和收入预期分别观察。

理论的可信度来自愿意接受不利证据。若预期改善而借贷与消费仍不恢复，需要检查人口、住房与供给条件；若保障扩展没有改变行为，也要检查待遇是否真正兑现以及原约束是否判断错误。不能把所有不符合预测的结果都解释为政策力度不够，否则命题就失去可检验性。

\section{生态约束不是账外背景}
生产依赖能源、水、土地和生态承载。降低单位能耗可以减少成本与污染，但若总产出扩张更快，总资源耗用仍可能上升。单位强度与总量必须同时记录。把污染产业迁往另一地区，可以改善本地指标，却未必减少整体损害。

% ARGUMENT P238
\phantomsection\label{arg:P238}
单位资源强度与总量的关系可以直接复算。设某项产出的单位能耗下降10\%，产量上升20\%，则总能耗变为原来的0.9乘1.2，即1.08倍。这个假设例子说明效率改善可以与总耗用增加同时发生，并不预测中国实际能源变化。判断环境结果必须同时记录强度、规模与能源结构。

生态损害还可能跨越地区和时间。污染排放发生地、最终消费地和承担健康损失的社区未必相同；土地修复或废弃设备处理可能在企业退出以后才发生。若责任只跟随当期营业主体，长期成本就可能留给居民或公共预算。生命周期与责任延续因此是生产关系的一部分。

自然条件不是可以无限被货币替代的投入。有些损害可以修复并估算费用，有些可能不可逆或存在较大不确定性。经济核算可以帮助比较方案，却不能因为缺少可靠货币价格就把损害设为零。遇到难以逆转的约束，需要在决策中明确边界与不确定性。

绿色投资既创造设备与服务，也需要材料、能源、土地和融资。评价应看生命周期、替代了什么、运行利用率及回收安排。装机能力不是实际发电，发电不是最终可用电量，低碳标签也不能替代完整比较。本文未取得可比的2026全年排放账，因此不提供虚构的当前总量。

% ARGUMENT P239
\phantomsection\label{arg:P239}
比较绿色投资，应先规定基准：如果不实施这个项目，需求会由什么技术满足，或是否本来可以减少需求？同一设备在不同电力结构、利用率和寿命下，净环境效果可能不同。只统计新增装机或设备销售，会忽略它究竟替代了什么。

生产、使用和报废阶段也要连接。制造设备需要材料与能源，运行阶段可能节约排放，退役以后还需要回收与处理。为了避免重复计算，同一减排收益不能同时无条件计入设备制造者、运营者和最终使用者的独立贡献。核算边界需要随指标公开。

这种分析不意味着必须等到所有参数完全确定才行动。可以报告合理区间、关键假设和敏感性，优先选择在多种情景下仍有益的方案。缺少2026全年可比排放数据时，本书只讨论机制和评价方法，不用截至某月的产量数据冒充完整生命周期结果。

生态约束还影响分配。能源价格调整可能对低收入家庭形成较大生活负担，关闭污染企业会影响特定地区就业。若转型收益广泛而成本集中，需要补偿与培训机制。补偿应帮助生活和生产转型，而非无期限维持高污染收益权。

% ARGUMENT P240
\phantomsection\label{arg:P240}
能源价格变化对家庭的影响，取决于支出结构和调整能力。收入较低、住宅保温较差或交通替代较少的家庭，可能难以迅速减少必要能源使用；高收入家庭拥有更多选择，也可能消费更多能源。补偿应依据实际负担与基本需要，而不只按平均消费推断。

产业转型的成本往往集中在特定工种和地区。关闭一座污染工厂不仅影响企业股东，也影响工人、供应商与地方预算。支持转型可以包括健康修复、技能、迁移和新活动，但不应把全部补偿交给原企业后假定劳动者自然受益。资金和服务需要沿实际承担者追踪。

公平转型与环境纪律可以相互支持。若受影响家庭知道基本生活和转岗有可靠安排，抵制调整的压力可能降低；若补偿变成永久保留旧收益，转型就可能停滞。因此，应区分对人和社区的支持，与对高污染生产权的无限保证。

\section{中国特征体现为传导方式}
较强的公共组织能力和金融协调，可以延长调整时间、集中资源和维护关键服务；它们也可能使风险长期留在公共或准公共账户中，弱化市场价格对损失的及时反映。分散市场决策能够提供试错，却也可能产生同时收缩和风险转嫁。两者的混合应通过结果研究，而非先验赞美或否定。

% ARGUMENT P241
\phantomsection\label{arg:P241}
公共协调能够改变危机的时间形态。延长债务期限、组织项目接续和集中提供流动性，可以避免短期支付中断毁掉仍有价值的生产网络；但若低回报没有改善，延期只是把损失分摊到更长时间。观察短期违约减少以后，还要检查资产质量、维护、收入和新增融资用途。

市场分散决策也有两面性。独立主体能够尝试不同方向，价格和退出提供反馈；同时，所有人都为自身减少支出，可能使他人收入下降并强化收缩。单个家庭或企业的审慎选择，不能直接推广为整个社会的最优同步动作。宏观制度的任务之一，是协调这种相互依赖。

中国情境中的关键，不是给公共或市场贴上固定优劣，而是观察两种机制怎样组合。公共支持是否保留信息反馈，市场退出是否保护基本功能，责任是否清楚，决定组合效果。制度名称提供背景，具体传导和结果才构成解释。

中国制度下值得重点观察的是地方财政、公共企业、银行和家庭之间的连接。土地收入改变不仅影响开发投资，也影响公共服务；公共债务安排改变不仅影响财政，也影响银行资产；家庭保障改变不仅影响消费，也影响劳动议价。制度特征通过这些具体联系产生解释力。

% ARGUMENT P242
\phantomsection\label{arg:P242}
地方财政与家庭生活之间，可以沿服务、采购和融资三条路径观察。收入变化影响预算安排，预算选择影响学校、医院和设施维护；采购支付影响供应商现金流，进而影响工资和经营；融资变化影响项目接续与银行资产。每条路径都有不同决策者与时间，不应压缩成“财政困难必然直接降低所有人收入”。

银行在其中既是债权人，也是支付与信用组织。容忍回款延迟可能给项目恢复时间，也可能减少用于其他经营者的资金；确认损失可以提高透明度，却可能带来资本补充需求。小企业与普通劳动者未必参与这些决定，却可能通过融资和就业承受结果。风险分析必须把这类间接承担者纳入。

家庭保障则影响反馈强度。收入中断后仍有可用医疗和生活支持，家庭可能较少削减必要消费，也更有条件寻找合适岗位；保障不足者可能接受更不利合同或透支健康。金融稳定与社会再生产由此不是两套互不相干的政策，它们通过家庭选择连接。

本书因此不预测某年必然发生总危机，也不把避免集中违约当作全部风险已经消失。危机可能表现为金融事件，也可能表现为长期低回报、收入机会不足和生活条件被缓慢消耗。区分不同形态，才能提出相应政策和反证条件。

% ARGUMENT P243
\phantomsection\label{arg:P243}
需要区分急性事件与慢性损耗。集中违约、挤兑和支付中断会迅速暴露风险；长期维护不足、收入机会减少、拖欠和照护透支则可能缓慢改变生活。后者不一定符合狭义金融危机定义，但仍是需要观察的经济损失。避免急性事件具有价值，不能因此忽略慢性成本。

同时，不能把任何增长放缓都重新命名为隐藏危机。经济成熟、人口变化或停止低效投资，也可能使增速下降而改善资源配置。要判断是否存在损耗，应观察服务质量、实际收入、健康、企业回款和可持续能力，而不是仅用一个不及预期的增长数字。

这一立场限制了本书的预测野心。本文提出风险传导条件，不给出未经模型与数据支持的崩溃年份，也不作永久安全保证。好的风险分析应当帮助识别哪些联系正在恶化、哪些缓冲有效，并允许新的证据改变判断。

\section{三个候选命题组成什么理论}
公众可兑现收益、家庭再生产能力和风险传递矩阵，分别把产权、需求与金融联系起来。它们的共同主张是：生产能力增加到生活改善之间存在制度传递环节，而这些环节也决定失败成本的落点。这是对既有产权、再生产与金融理论的综合，不是独立发现所有组成思想。

% ARGUMENT P244
\phantomsection\label{arg:P244}
三个候选命题可以分别写成条件式。第一，在公共资产收益为正的条件下，治理、预算与服务规则影响公众能够实际取得多少收益；第二，在生产扩张的条件下，家庭收入、时间和保障影响劳动能力与需求能否持续；第三，在发生损失的条件下，合同、控制与救助安排影响损失落点及下一轮选择。这些命题都不能由“公有”“增长”或“稳定”一个名称直接推出结果。

它们的共同贡献是增加同时核对的要求，而非发明一切组成思想。产权与组织理论已经讨论控制，社会再生产研究已经讨论生活条件，金融研究已经讨论资产负债表与救助预期。本文尝试将它们用于同一中国制度变化，观察一个账户的改善是否伴随另一个账户承担了成本。

如果这种联结不能产生更清楚的区分或更好的预测，就没有理由把它称为新的成熟理论。创新应体现为可操作问题和可检验解释，而不是在原有概念前增加时代或国别修饰。附录给出变量、竞争解释和可能推翻命题的观察，作为后续研究的起点。

可以检验的新增要求，是将同一制度变化在生产、收入和风险三本账中的结果同时追踪。例如公共企业收益规则改革之后，不只观察分红，还观察服务、投资与担保；平台保障扩围后，不只观察参保记录，还观察收入、价格和事故家庭损失。只有这样，才能检验改革是否减少总成本，还是改变了成本所在账户。

% ARGUMENT P245
\phantomsection\label{arg:P245}
联合追踪可以从一项范围明确的改革开始。假设公共服务企业改变分红规则，生产账记录投资、维护与服务数量，权利与收入账记录分红、价格和预算去向，再生产与风险账记录服务可靠性、家庭负担和担保责任。三本账有交叉联系，但不能简单相加为一个金额，否则可能重复计算。

同样，一项平台保障改革应同时观察安全投入与任务安排、费用与报酬、事故后的最终损失。若平台利润下降而家庭事故损失下降更多，可能减少了总成本；若仅把费用转给劳动者而保障未兑现，就需要重新判断。制度效果不是由某一个主体的财务结果单独决定。

识别仍然需要对照。改革对象可能本来就有较好管理或较严重问题，前后变化还可能受到需求与价格冲击。联合账户提高解释完整性，却不会自动解决因果推断。本文把核算完整与识别可信看作两项都需要完成的工作。

若三本账的联合信息不能比单一指标更好解释结果，本文框架的额外价值就有限。若它能识别过去被遗漏的转移和反馈，才值得进一步发展成可估计模型。理论需要接受这一检验，而不是依靠标题的宏大获得成立资格。

% ARGUMENT P246
\phantomsection\label{arg:P246}
框架的评估可以包含三种层次。首先检验是否减少遗漏和重复计算；其次检验新变量是否提高对家庭或企业结果的解释；最后在可获得合适数据时，检验具体制度变化的因果效应。第一层通过不意味着后两层已经证明，理论文字完整也不能替代数据工作。

可能的反例需要认真处理。如果资产名义归属已经足以稳定预测公共收益，新增治理变量没有贡献，那么产权传递框架需要收缩；如果再生产条件改善却不能解释劳动和需求差异，应重新检查机制；如果风险矩阵只是在事后重述已知损失，没有帮助预测行为，其增量价值也有限。

这一要求使理论保持可修订。下一章提出的制度方向，同样只在明确约束和条件下成立。它们不是从宏大理论中自动演绎出的唯一政策，而是根据前文机制形成的可比较方案，需要接受预算、实施和结果检验。


\chapter{制度改进的条件与代价}
\section{改革从约束出发}
改革不是把所有理想目标同时写进一份文件。财政资源、专业人员、组织能力和家庭时间都有限，短期需求支持与长期结构调整也不一定同步。政策比较应先说明要解决的约束，再说明资源来源、受益者、承担者和失败条件。

% ARGUMENT P247
\phantomsection\label{arg:P247}
把约束放在前面，意味着先回答问题发生在哪里。居民买不起已有服务，与当地根本没有服务，需要不同工具；企业有订单却收不到货款，与没有可持续市场，也不能只用同一种贷款安排。政策名称相同，若针对的约束不同，实际效果就可能相反。

改革比较还应建立明确基准。与什么都不改变相比，一项措施可能改善生活；与另一种成本更低的方案相比，它又可能不合算。基准应包括维持现状的真实费用，例如拖欠、等待、设施损坏和家庭时间，不能把现状视为零成本。

本章提出的是实施方向与检验条件，不是已经完成全国预算测算的政策套餐。各项建议之间存在资源竞争，不能把单项收益直接相加。最终组合需要具体地区、人数、服务能力和财政资料支持。

本章采用五项公开标准：能否改善真实生产能力，能否保障基本生活，能否扩大实际机会，能否使风险分担更合理，以及能否维持财政与生态条件。不同读者可以给这些标准不同权重，但不能省略成本。本文没有估算全国改革套餐的精确预算，不把示例金额包装成已经完成的政策测算。

% ARGUMENT P248
\phantomsection\label{arg:P248}
五项标准之间可能发生取舍。例如提高服务可靠性需要备用人员，短期利用率可能降低；改善偏远地区基本服务的单位成本可能高于人口密集地区，但基本机会更均等。只用单一成本效率指标，会把这些目标冲突隐藏起来。公开标准允许不同选择被比较，而不是保证所有目标同时达到最大值。

评价还需报告分布。平均收入改善可能伴随少数家庭失去基本生活，财政总额稳定可能伴随某些地区服务明显下降。对于基本保障，可以设定最低可接受条件；对于超过底线的资源安排，再讨论效率与偏好。底线本身属于需要公开说明的价值选择。

如果一项政策必须依赖“所有人都会因此增加消费”才能成立，就应检验这种行为假设。改善保障的价值不完全取决于消费是否立即增加；相反，刺激支出也不能自动证明长期福利改善。目标与评价变量应保持一致。

政策时序也很重要。若先提高家庭负担再承诺未来保障，短期信心可能恶化；若先提供救助却迟迟不明确退出条件，可能强化冒险预期。过渡安排不是改革之外的附属问题，而是决定改革能否实施的一部分。

% ARGUMENT P249
\phantomsection\label{arg:P249}
时序可以按能力建设、过渡支持与长期规则三层安排。先建立可查询的权益和可用服务，再逐步调整责任与筹资，能够减少承诺无法兑现的落差；但准备阶段也应有期限和可观察进展，避免以条件未成熟为由无限拖延。

对集中受损者，过渡支持需要在冲击发生前尽可能明确。企业退出以后才临时讨论工人工资、社保和安置，容易使最弱一方承担等待成本。相反，预先说明基本生活保护与股东损失责任，可以使调整更可预期。

过渡安排不能永久替代改革。延长偿付、临时补贴和试点例外，都应与信息披露、里程碑或退出条件相联系。时间的价值在于获得学习和调整空间，而不在于让原有问题变得更难观察。

\section{把基本保障与劳动流动连接}
优先方向是提高基本保障的可获得性、可携带性和支付可靠性，包括失业期间的生活支持、跨地医疗使用以及养老权益衔接。其机制是降低劳动者接受新岗位、迁移或拒绝严重不利条件时的基本风险，而不是要求每个人增加借贷和消费。

% ARGUMENT P250
\phantomsection\label{arg:P250}
保障可携带首先需要权益随人记录，随后还需要费用能够在服务使用地结算。只有记录而没有结算，迁移者可能在文本上有资格却仍须垫付；只有结算而缺少服务供给，流入地又可能出现等待。信息、资金与服务能力应同步安排。

优先受益者可以从实际中断风险识别，包括频繁换岗、跨地区就业和不稳定任务劳动者，而不是仅按某个职业名称分配。对稳定受薪者，也要处理家庭成员和照护责任导致的空档。资格规则越依赖复杂证明，最缺少时间和信息的人越可能被遗漏。

改革应允许个人查询、纠错与申诉。系统互联能够降低重复办理，但错误记录如果没有低成本更正渠道，也可能扩大新的排斥。实际可用性需要从办理时间、垫付压力和支付结果来验证。

资源可来自一般税收、保险缴费、既有支出重排和中央对地区间服务成本的分担。每一种来源都有承担者：提高缴费可能影响企业用工或到手工资，税收融资需要评估税基与转嫁，支出重排意味着其他项目减少。不能把“盘活资源”当作没有机会成本的资金来源。

% ARGUMENT P251
\phantomsection\label{arg:P251}
筹资选择应分开短期过渡与长期经常支出。一次性资产处置可以支持系统改造和历史欠账清理，却不适合独自承担永久待遇；经常性服务需要相对稳定的税基、缴费或预算安排。借款可以平滑过渡，但未来偿付仍需明确来源。

中央对地区间成本分担，可以减少流入地独自承担服务人口增加的压力，但计算应跟随实际服务人口和需要，而非只按户籍或建设数量。地区也应承担信息真实和服务质量责任，防止为了获得补助虚报需求。

如果提高企业缴费作为主要来源，需要观察正规用工成本和工资转嫁，特别是利润较薄的小雇主。可以比较不同筹资组合，但不能承诺企业、劳动者和纳税人同时完全不承担新增成本。改革的理由在于共同分担比各自暴露于风险更合适，而非费用凭空消失。

受益者主要是保障不足和流动较多的人群，企业也可能从更稳定的劳动力匹配中受益。实施需要统一可核验的权益记录、跨地区结算以及申诉渠道。过渡成本包括系统衔接、历史欠账和流入地服务扩容。若只有资格扩展却没有服务能力，承诺会变成等待；若负担主要压在最脆弱的小企业，可能抑制正规就业。

% ARGUMENT P252
\phantomsection\label{arg:P252}
实施应把最容易发生断点的流程画出来：离职、迁移、重新就业、受伤、申请和支付。每一步由谁提供记录、谁承担审核责任、超期怎样处理，都应明确。跨部门合作不只是连接数据库，还涉及预算和责任授权。

服务扩容需要按实际需要排序。劳动者集中流入的地区可能首先需要门诊、子女教育与租住支持，老龄地区可能更需要长期照护。统一扩围并不等于每地建设同样设施。资源配置应允许需求差异，同时保持基本权利不因地区财力完全割裂。

失败情景需要预先测试，例如平台退出后记录能否保留，雇主欠费时劳动者能否获得必要待遇，跨地区争议由谁处理。只在正常就业路径上运行顺畅的制度，未必能够在最需要保障的冲击时刻发挥作用。

评价不能只数新增参保人数，应观察实际支付、医疗与失业冲击后的生活损失、转岗时间以及工资和工时。若保障增加却没有改善这些结果，需要检查覆盖、质量和行为反应，不能只继续提高投入总额。

% ARGUMENT P253
\phantomsection\label{arg:P253}
评价可以同时设置过程与结果指标。过程包括办理时长、垫付金额、拒付原因和纠错比例；结果包括重大冲击后的基本消费、就业中断与家庭负债。前者帮助定位问题，后者判断是否真正改善生活，不能只选容易上报的登记数量。

还应观察意外行为变化。企业是否调整合同以回避责任，劳动者是否因为信息不清错失待遇，服务机构是否增加不必要项目，都可能影响净效果。发现问题后应改进规则，而不是用总体覆盖增长掩盖它们。

若实施初期结果尚未改善，需要区分过渡延迟与机制无效；到预定观察期仍无实质变化，则应允许重新分配资源。保障改革也应接受证据约束。

\section{将照护作为生活基础设施}
托育、养老和社区支持适合按实际需求配置，不能复制房地产式大规模建设逻辑。已有场地改造、人员培训、服务质量和时间安排，可能比新建大型设施更重要。公共供给、购买服务与家庭补贴可以组合，但每种方式都需要质量监督。

% ARGUMENT P254
\phantomsection\label{arg:P254}
照护设施的选址需要从生活路径出发。儿童托育应考虑家长工作地点与通勤，老年服务应考虑行动能力、医疗联系和家庭探访。闲置房屋便宜不代表适合改造，远离使用者的低成本场地可能把费用转成家庭交通和等待时间。

服务时段也需要与岗位组合匹配。轮班劳动者、平台劳动者与固定办公时间人员的需要不同；延长服务则需要增加人员与合理排班。把家长的时间约束转成照护人员无偿加班，无法形成可持续供给。

因此，可以先从可达场地、小规模连续服务和专业人员训练开始，再根据使用与质量扩大。建筑面积是投入，可靠照护才是产出。服务对象难以充分表达意见时，质量观察更需要独立性。

资金可由经常性财政收入、合理使用者付费与保险安排共同支持。对基本服务应设可负担边界，防止低收入家庭因费用退出。资本性建设支出与长期人员运营费用应分别预算，避免设施建成后缺少专业人员和日常维护。

% ARGUMENT P255
\phantomsection\label{arg:P255}
用一个明确假设的预算说明长期责任：某服务点一次性改造投入100万元，每年人员、维护与其他运营费用合计60万元。即使改造费用由一次性资金解决，五年运营仍需300万元，且未计价格变化。这个例子不代表任何实际项目成本，只说明建设融资不能替代运营来源。

使用者付费可以提供部分资源并反映需求，但基本服务若完全按支付能力分配，可能排除最需要支持的家庭。补助可以随收入和照护需要调整，也要避免复杂审核造成高门槛。设计应比较实际到手帮助与管理费用。

经常性支出还须覆盖质量。人员工资、培训、保险与休息不足，会提高流动并损害服务连续性。所谓普惠不能仅靠压低从业者报酬实现，应把家庭可负担与劳动者可持续共同纳入预算。

受益者包括需要照护的人、承担无酬劳动的家庭成员及服务劳动者。成本承担者是纳税人、使用者和相关公共预算。过渡中需要培训、资格认定与保障从业待遇，不能依靠低薪高流动劳动支撑“普惠”。质量失败会伤害最难表达意见的儿童、老人和失能者，必须有独立投诉、检查和退出机制。

% ARGUMENT P256
\phantomsection\label{arg:P256}
委托社会机构提供服务时，合同应明确质量、人员、费用、数据与退出责任。仅按接收人数付款，可能鼓励数量而忽略需要更复杂的服务对象；仅按设施规模补助，又可能形成低使用率。可以结合基本运营与实际服务，但指标应避免诱导排斥高需求者。

监督也需考虑受服务者的表达能力。儿童、失能老人和部分残障者难以独立投诉，应有家属参与、必要检查与可追踪纠错，同时保护隐私和尊严。一次满意度调查不足以代替持续质量监测。

机构退出时还需要服务接续。人员工资、预付费、档案和照护转移如果没有安排，经营失败会迅速变成家庭危机。把退出责任写进合同，可以使风险管理前移，而不是每次出现困难后再临时处理。

政策检验要记录照护质量、实际利用、家庭时间和就业选择。若设施利用率低，应区分价格过高、时间不合适、信任不足和地点错误。提高利用率不能靠强迫家庭改变偏好，而应改善服务与需要的匹配。

% ARGUMENT P257
\phantomsection\label{arg:P257}
利用率低有不同原因。价格障碍需要负担支持，时间错配需要调整时段，质量担忧需要改善服务和信任，地点错误则可能要求迁移或停止扩张。单纯降价未必解决质量问题，增加宣传也未必解决距离问题。

家庭时间释放以后，可能用于有酬工作、学习、休息或其他照护。这些用途具有不同但都可能重要的价值。若只把新增就业收入作为收益，就会低估休息、健康和生活自主；若宣称所有释放时间都转成生产，又会高估宏观效果。

政策可以询问家庭是否实现原本希望的安排，并结合服务质量、费用和成员时间记录进行评价。照护改革的目的，是扩大真实选择，而不是替家庭规定唯一生活方式。

\section{公共资产治理从规模转向兑现}
公共企业可以建立经营任务与公共任务分开说明的报告：商业业务的成本与回报，政策性任务的目标、补偿和覆盖，重大担保与风险，以及收益用于再投资、财政和服务的比例。分开报告不意味着机械拆分组织，而是使公众能够判断资源用途。

% ARGUMENT P258
\phantomsection\label{arg:P258}
公共企业的任务分账，应先区分商业活动与明确公共义务，再说明共同成本怎样分配。若同一网络同时服务高收益地区和低收益地区，不能简单按局部利润要求后者退出；但公共义务也不应成为掩盖所有低效率的理由。需要公开服务目标、成本依据与补偿方式。

报告应能连接预算决策。例如维护不足是否需要增加投入，低价公共服务是否获得相应补偿，盈余上缴是否影响必要更新。只有信息而没有调整权限，透明度可能停留在形式；只有权限而没有信息，又难以问责。

居民与普通职工应有适合其参与的反馈渠道。居民了解服务中断和质量，职工了解实际流程和安全，但他们未必掌握财务整体。不同信息需要汇合，不能用单一满意度或单一利润指标替代全部评价。

增加上缴收益可能支持公共服务，也可能减少企业必要投资。合理比例应根据资本维护、行业风险和真实项目需求确定，不能对所有企业设置同一回收速度。资源来源是可持续净收益，不是依靠额外借款制造表面分红。

% ARGUMENT P259
\phantomsection\label{arg:P259}
可持续净收益要在必要维护、合理投资与风险准备之后判断。高分红如果来自新增借款或减少维护，可能只是把今天的财政收入换成未来负担；过度保留收益如果缺少有效项目，也可能扩大低回报投资。应比较资金留在企业与进入公共预算的真实用途。

行业差异决定合理安排。自然垄断公共服务、竞争性制造业务与金融机构的资本需求不同，不宜采用完全相同的分配速度。还需要明确突发损失以后如何调整，避免短期刚性要求诱导报表操作。

分红规则本身不能保证公众受益。进入预算的收入仍需通过教育、医疗、保障或其他公共用途传递。完整评价必须沿资金去向继续追踪，不能在企业上缴环节就宣布产权收益已经兑现。

受益者取决于财政与服务的最终分配。若利润上缴后主要用于新的低回报项目，产权收益依然没有改善居民生活。实施需要审计、预算公开、服务质量评价以及对管理者的长期责任。风险是把难量化的公共任务削减以提高财务指标，所以评价应同时保留服务与财务两本账。

% ARGUMENT P260
\phantomsection\label{arg:P260}
可以用服务连续性约束财务目标。某项改革提高利润却减少必要维护，短期账面改善可能以未来故障为代价；降低收费若没有稳定补偿，也可能损害服务。价格、质量、投资与债务应作为相互约束的结果共同报告。

管理责任宜与可控制范围相匹配。要求执行者对其无权改变的公共任务全额盈利负责，可能诱导排斥服务；允许管理者以任务复杂为由免于解释可控成本，又会削弱约束。权限、目标与责任需要同时清楚。

对不同阶层，改革可能通过价格、税收、就业和服务产生交叉影响。高收入与低收入家庭使用公共服务的方式可能不同，普通职工也可能承担内部调整。平均公共收益增加以后，仍需检查是否有基本服务群体被遗漏。

候选试点可以选择任务较清晰、数据较完整的公共服务企业，比较改革前后维护、价格、质量与债务。若透明度增加却没有实质问责和资源调整，改革可能只增加报告负担。政策的成效在结果，而不在新指标数量。

% ARGUMENT P261
\phantomsection\label{arg:P261}
试点可以从任务较单纯、历史资料可比的单位开始，但推广时必须承认样本选择。管理基础较好的单位更容易成功，资源薄弱地区可能需要额外能力建设。将试点结果直接乘以全国规模，会忽略这种差异。

对照应比较相近任务、需求和成本条件，避免把行业景气或自然条件误算为治理效果。短期可观察维护、支付和服务，中期观察债务与使用者负担，长期再评估能力。评价计划应在实施前公开核心指标，减少事后挑选。

若报告数量增加却没有改善纠错、预算调整和服务，应该简化形式负担并检查责任机制。透明制度的价值在于帮助决策与问责，不在于形成更多难以使用的表格。

\section{地方财政与债务调整}
短期可通过合理置换和跨期安排缓解集中偿付，长期需要让稳定支出责任对应稳定收入，并降低对土地一次性收入的依赖。中央与地方的分担应考虑服务人口、地区财力和历史债务，不宜只奖励当期扩张。

% ARGUMENT P262
\phantomsection\label{arg:P262}
债务置换可以改变利率与到期结构，但不自动创造项目收益。它适合缓解集中偿付，争取核查与调整时间；如果置换以后继续形成相同低回报承诺，压力仍会累积。因此，短期融资安排应与支出责任、项目选择和信息披露一起改变。

稳定收入与稳定支出的对应，也需要考虑服务人口迁移。某地承担更多常住人口的教育、医疗和基础设施，却无法取得相应财力，可能依赖不稳定收入。转移支付应尽可能对应真实服务需要，同时保留地方节约与质量责任。

减少土地收入依赖不意味着可以瞬间填补所有缺口。既有债务与公共承诺需要过渡预算，税源调整需要征管能力。政策应明确先保护哪些基本服务、哪些项目放缓和哪些资产可以合理处置。

债务重组的资源来自未来财政能力、债权人让利、资产处置与必要的上级支持，不能把所有负担隐藏在银行展期里。债权人承担损失可能影响金融资本，需要同步说明银行资本补充及对支付体系的保护。资产处置则要防止在低迷期不透明低价转让。

% ARGUMENT P263
\phantomsection\label{arg:P263}
重组首先需要确认同一损失，避免不同主体各自把它转出账外。银行让利可能减少其资本，财政补充资本又形成公共支出，资产处置可能影响未来收益。把这些步骤合并看，才能判断总体资源损失与分担；单看每一步“风险降低”，容易重复计算改善。

支付体系与原有投资收益应分别保护。维持基本金融服务具有公共价值，却不必保证所有股东和债权人取得原约定回报。具体责任需要依法依约处理，本书只提出经济评价原则，不给个案债务安排作法律判断。

处置还应减少信息不对称。公开估值依据、竞价和关联关系，可以降低低价转移公共资产的风险；保留资产等待更好价格，也要比较维护、融资与市场风险。两种路径都不是天然无成本。

受益者包括依赖基本服务的居民、有经营能力但受拖欠影响的供应商和能够恢复运转的项目。原有股东、债权人、土地相关利益与纳税人可能分别承担成本。改革应优先保护必要公共功能和已经发生的合理劳动收入，同时清理没有前景的投资承诺。

% ARGUMENT P264
\phantomsection\label{arg:P264}
保护对象应围绕基本功能和已经发生的合理权利识别。供应商已交付产品、劳动者已完成工作，与股东对未来收益的期待，经济性质不同；必要公共服务的连续性，也不同于维持原项目全部规模。把这些对象混为一谈，容易使救助资金先流向最有谈判能力的人。

小企业可能兼具债权人与雇主身份。清理欠款能够帮助其支付工资，但需要核验资金是否真正到达劳动与经营环节。对没有持续经营前景的企业，也应考虑已完成劳动的支付和有序退出，不能要求劳动者以等待企业复苏无限垫付成本。

公共支持应记录最终受益。名义上救助项目，实际可能保护债权人；名义上支持企业，实际可能维持居民必要服务。两者都可以有理由，但理由与承担方式需要公开，不宜仅用一个笼统的稳定目标替代。

过渡风险是地方在收入收缩时削减维护、拖欠采购或增加不稳定收费。应建立基本服务底线、欠款透明与新增项目约束。若救助之后仍持续形成同类债务，说明激励没有改变；若约束过急导致有用项目停摆和就业集中受损，也需要调整节奏。

% ARGUMENT P265
\phantomsection\label{arg:P265}
财政调整中的隐性压缩常发生在维护和付款。推迟修缮、延迟采购结算或增加临时收费，短期不一定显著改变正式支出指标，却会增加居民与企业成本。基本服务底线应因此包含质量与及时支付，而非只保持预算名义数额。

新增项目约束需要穿透不同组织形式。把支出移到公共企业或其他承诺安排，不应使真实公共责任从观察中消失；同时，也不能把所有企业经营风险无条件并入财政。合并范围应依据具体控制、担保和责任关系，并披露不确定性。

政策节奏要与地区承受能力相适应。过快停止必要项目可能损坏已形成能力，过慢确认损失则可能继续占用资源。可以根据项目完成度、可用价值与未来现金流分组处理，而不是只按是否负债作统一判断。

\section{住房政策回到居住与交付}
住房政策可以优先处理交付、基本居住和城市存量改造，而不是保证全部资产持有者的历史买入价格。对未完成项目，应评估后续建设需要和项目可回收价值，明确购房者、金融机构、开发商及公共部门的责任。

% ARGUMENT P266
\phantomsection\label{arg:P266}
处理未交付住房时，应先核实工程状态、剩余投入、已收资金、权利关系与可实现价值。已经接近完工的项目，与位置不佳、需求不足且资金缺口很大的项目，不能使用同一方案。统一填补所有缺口可能消耗过多公共资源，完全放任也可能使购房家庭和周边生活长期受损。

目标是恢复合理居住功能，并在明确规则下分担既有损失。原开发股东、债权人、购房者与公共部门掌握的信息和控制不同，责任安排应反映这种差异。具体法律顺位需要按适用制度和个案判断，经济分析不能凭一句保交付替代。

项目完成以后仍需质量与维护。只以交付钥匙统计完成，可能遗漏基础设施、消防和长期维修问题。居住效果应成为验收的一部分，使投入真正对应家庭可使用的住房。

资金可以来自项目资产处置、合规融资、股东承担与有条件财政支持。任何公共投入都应说明是否取得相应权利和未来回收。保交付不应变成对集团所有旧债无差别担保，否则必要居住功能与不合理风险承担会再次混在一起。

% ARGUMENT P267
\phantomsection\label{arg:P267}
公共投入可以采取有条件支持，并明确是否取得项目权益、未来回收或可用住房。若只投入现金而不保留可追踪权利，最终收益可能重新流向原有控制者；若要求不现实的高回报，又可能妨碍必要交付。条件应根据项目性质而定。

新增资金还应避免被集团其他负担吸收。项目层面记录资金用途、支付对象和工程进度，可以帮助区分完成居住功能所需费用与其他历史债务。账户隔离不是损失消失，而是让剩余投入优先用于明确目标。

购房家庭的过渡居住与现金压力也需被观察。工程恢复需要时间，等待期间租金、还款和工作安排可能同时发生。是否提供支持，应依据实际需要和资源约束，而不能只在最终交付日评价家庭处境。

保障租赁和存量收购需要依据人口与就业需求。受益者应是有真实居住需要的家庭，成本包括购买、改造、管理和长期维修。低入住率、远离工作地点、质量差和不透明定价都是失败条件。改造后实际居住效果比收购套数更接近政策目标。

% ARGUMENT P268
\phantomsection\label{arg:P268}
存量收购是否划算，需要比较收购价格、改造、空置、运营和维修，与新建或直接补助租住等替代方案。低购入价格未必意味着低全周期成本，结构不适合、交通不便或长期空置都可能抵消价格优势。

需求核验应结合就业、人口与家庭支付能力。某地区有很多未售住房，不等于当地存在同等规模的可持续租赁需求；有居住需要的家庭，也未必能承担改造后的租金和通勤。将供给与具体使用者匹配，比完成收购数量更重要。

退出和调整同样必要。若实际需求持续不足，应允许改变用途、暂停收购或处置，而不是为维持指标继续投入。公共住房的价值在使用与稳定生活，不在把全部市场库存转到公共资产表。

税费与土地收入结构调整则应考虑现金流不足但拥有自住房的家庭、小业主和租户的差异。渐进实施、清晰估值和必要缓缴可以减少集中冲击，但也会增加行政复杂性。没有可靠征管和过渡方案时，不宜用理论上理想的税种代替实施研究。

% ARGUMENT P269
\phantomsection\label{arg:P269}
资产相关税费讨论需要把净资产与现金收入区分。拥有自住房但现金收入较低的退休家庭，与持续取得多项资产收益者，支付条件不同。必要缓缴或分期可以减少被迫出售，但会推迟财政收入并增加管理成本，应公开这些取舍。

租户受到的影响取决于转嫁条件。税费改变可能影响租金、供应和维护，不能事先承诺房东或租户独自承担。需要结合地区供需和合同观察，必要时配合居住支持与透明租赁规则。

改革还需避免重复依赖不稳定估值。市场快速变化时，评估方法、复核与申诉非常重要。本文不将某项尚未完成地方测算的税种当作普遍财政缺口的现成答案，重点是建立可执行的筹资与过渡安排。

\section{技术、竞争与劳动改善相互配合}
产业支持应优先补足共性研发、检测、基础软件、技术扩散和职业培训等外溢较强的环节。资金来自公开预算或有明确风险边界的投资工具，失败成本需预先说明。只按设备、专利和产能发放支持，可能鼓励可见投入而非有效创新。

% ARGUMENT P270
\phantomsection\label{arg:P270}
支持技术扩散往往不同于支持少数领先企业。共性检测、维护培训和基础工具可以帮助更多小企业采用已有技术，效果可能表现为质量稳定与成本下降，而非新增专利数量。评价应区分前沿探索与扩散应用，允许它们使用不同指标。

公开资助需要说明选择与退出。可以采用阶段拨付、独立技术评价和利益冲突披露，同时允许真正不确定的研究失败。失败成本应由事先约定的风险承担者承受，而不是通过拖欠工资和供应商款项转给没有选择权的人。

产业收益最终还需经过就业、工资、公共收入和价格传递。资助企业销售增长不能独自证明公众获得净收益，应同时比较补助成本、挤出效应和没有获得支持的竞争者处境。

竞争政策可以针对排他规则、合同不透明和不合理退出成本，保留规模带来的协调收益。平台互操作可能降低锁定，也会增加安全、隐私和兼容成本，需要分层接口与责任安排。不能因开放有收益就假定所有数据都应无条件共享。

% ARGUMENT P271
\phantomsection\label{arg:P271}
竞争措施可以从具体退出障碍入手，例如迁移商品与评价记录的合理方式、合同变更通知和争议处理。互操作程度可以分层，而不必要求所有系统完全开放。基础信息可携带与敏感信息共享具有不同风险，需要分别安排。

开放成本还需要分担。接口维护、安全测试与兼容需要资源，如果全部由小经营者承担，形式开放未必带来实际进入。平台与使用者责任应清楚，同时避免要求公开可直接被用于欺诈的细节。

效果应观察替代渠道是否真正增加、费用和服务是否改善，而不是只看接口数量。若开放导致风险上升却没有改善竞争，需调整设计；若有明确收益，也应核验是否覆盖依赖程度较高的小商家和劳动者。

劳动保护方面，可提高工资与结算透明度，强化重大算法决定的申诉，推进平台工作风险的可携带保障。成本可能通过利润、收费和价格传递，应持续观察。若保障规则导致平台将任务进一步外包并逃避责任，就需要按实际控制而非单一合同名称判断责任链。

% ARGUMENT P272
\phantomsection\label{arg:P272}
按实际控制追踪责任，可以避免外包链条把义务变得模糊。谁决定任务标准、时限和报酬，谁掌握事故记录，谁有能力纠正危险安排，都是需要调查的事实。它们可以为责任讨论提供依据，但不能在缺少个案材料时直接替代法律认定。

结算透明需要让劳动者知道单价、扣费、等待和处罚怎样进入最终报酬。规则变更应留出理解和调整时间，重大决定应能够申诉。对文化程度、语言或数字技能不同的人，信息必须实际可用，不能仅在很长的协议中形式公开。

保障费用如何分担，应持续跟踪收入、价格与服务变化。制度若使平台把任务进一步转移到更难监督的承包商，就需要修订责任链，而非只要求劳动者自行识别复杂组织结构。

技术培训应以实际转岗和收入结果评价，不能只看结业证书。培训内容与本地岗位脱节，或者学员因照护和交通无法参加，都会降低效果。必要时将培训、生活支持与企业真实需求共同安排，避免要求最脆弱的人先自行承担全部转型成本。

% ARGUMENT P273
\phantomsection\label{arg:P273}
培训的经济约束常在课程之外。失业者可能缺少交通费和生活来源，照护者可能无法在固定时段参加，已有岗位者可能没有可用于学习的时间。把培训作为唯一支持工具，可能让资源较充足者更容易受益。

课程应与真实任务和可达岗位连接，同时保留可迁移基础能力，避免只为一家企业培养高度专用技能。企业参与可以提高相关性，也需要防止把公共培训变成替其支付全部内部培训费用。

评价应追踪就业质量与持续性，并比较未参加或尚未参加者的合理对照。结业以后短期找到工作，不一定意味着长期收益；较长时间没有岗位改善，也应允许停止低效课程。培训数量只是投入与过程。

\section{生态转型的补偿与纪律}
环境定价和标准可以使生产者承担更多社会成本，但价格传递可能影响居民基本生活。对低收入家庭的补偿宜针对生活需要，而不是继续补贴高耗能使用。对受冲击地区的支持宜用于修复、技能与新活动，避免把旧产能永久保留下来。

% ARGUMENT P274
\phantomsection\label{arg:P274}
环境价格提高以后，对低收入家庭可以采用与基本需要相关的支持，使其维持生活而仍保留节能激励。若按实际耗能越多补贴越多，可能削弱原有价格信号。具体补偿应考虑住宅条件、地区气候和可替代交通，而非只按统一人数推算。

企业改造支持应结合可验证的环境结果和期限，避免长期为已经缺乏前景的产能承担全部成本。对受影响工人和地区，则可以提供技能、健康修复与公共服务支持，使转型收益与集中负担相协调。

这种安排并不保证没有冲突。财政有限、环境目标与就业过渡可能存在时间张力，应公开排序与补偿边界。隐藏冲突只会使实际成本由缺少谈判能力的人承担。

资源可来自环境相关收入、一般预算和受益主体投入；这些收入本身会随排放减少而变化，不能被当作永久增长的财源。基础设施项目还需要全生命周期维护预算。若转型主要由地方债融资而收益长期不明确，就可能把生态目标与新的财政脆弱性绑在一起。

% ARGUMENT P275
\phantomsection\label{arg:P275}
环境相关收入具有内生性：如果政策成功减少排放，相关收入可能下降。因此，不宜用持续增长的环境收入支持全部永久性服务承诺。可以用于过渡和修复，并为长期支出安排更稳定来源。

生态项目融资也应记录运营、维护与退役。建设期贷款容易形成可见投资，后续维护却可能缺少收入。若项目提供公共环境收益，应明确由谁长期支付，而不是假定未来使用者付费必然覆盖。

地方财政压力较大时，更应比较分阶段、维护改造和需求管理等方案。绿色目标本身不能使资金成本与低利用率消失。经济评价要保护环境目标，同时防止以该目标为名重复低效建设。

受益是健康、生态与长期生产条件，成本则常集中于某些企业、工种和地区。应提前公开关闭、改造和补偿安排，减少突然冲击。检验不能只看本地排放降低，还需防止污染外迁、投入品排放遗漏以及低利用率设备重复建设。

% ARGUMENT P276
\phantomsection\label{arg:P276}
地区转型应尽早公布合理安排，使劳动者和家庭能够调整教育、住房与就业计划。突然关闭可能增加集中损失，无限延期又会累积健康与环境成本。过渡期限需要结合危害程度和替代条件，而非统一延长。

污染外迁可以通过供应链与地区合作追踪。一个地方减少排放而增加外购高污染投入，可能只是改变生产地点；反之，采用更清洁技术形成真实减少，则应在一致边界下确认。生产地与消费地的责任讨论需要同一套数据。

对补偿的核查，应关注实际到达家庭、修复工程和新能力的资源。资金拨给原企业并不自动意味着工人得到支持，地方获得专项资金也不自动意味着环境改善。终点结果应当能够被观察。

\section{实施顺序与退出标准}
在数据和执行能力不足时，政策应从可核验的基本功能开始：及时支付、服务可用、权利可携带、风险披露和责任明确。它们未必产生最醒目的短期总量，却能改善其他改革发挥作用的条件。大规模扩张应在约束、维护与退出安排明确后推进。

% ARGUMENT P277
\phantomsection\label{arg:P277}
实施顺序可以优先处理阻碍其他政策发挥作用的环节。拖欠未解决时，新增贷款可能先被用于填补旧回款；权益记录不连续时，扩围可能难以兑现；服务人员不足时，新增设施可能闲置。这些前提工作不一定带来显眼数字，却决定后续投入能否有效。

优先级还应结合紧迫性与可逆性。基本生活即将中断的家庭需要及时支持，长期不可逆投资则应留出充分论证。不能把所有项目按同一种审批速度处理，也不能因追求谨慎而拖延已经清楚的基本支付责任。

可以按约束、资源、责任、结果四项建立公开实施清单，保持简洁并允许修订。管理工具应帮助行动，不应成为增加基层填表负担的另一层目标。

试点不是可以无限续期的免责名称。开始时应确定目标群体、预算、观察窗口、对照方法和停止条件。若出现预算失控、利益被少数主体截留或基本服务质量恶化，应允许修改和退出。也要防止只选择条件最好的地区试点，却把结果无条件推广到资源更弱地区。

% ARGUMENT P278
\phantomsection\label{arg:P278}
试点前应说明核心问题和评价窗口，并记录为何选择这些地区或群体。若只选择财力最强、管理最好的地区，得到的是有利条件下的效果；推广前还需检验资源较弱地区能否实施。失败案例同样有价值，应记录发生在哪个环节。

停止条件可以包括基本质量持续不达标、费用明显超出预算、利益被截留或目标群体难以使用。触发以后应有调整、暂停或退出机制，同时保护已经依赖服务的人。项目退出不应突然切断必要生活支持。

评价者的独立性也重要。实施部门掌握信息，但可能更愿意展示成功；外部评价有比较优势，却需要理解实际流程。两者合作并公开方法，有助于减少仅凭指标完成情况判断的偏差。

政策效果可能滞后，短期没有变化不总等于失败；但“长期才见效”必须配合中间指标，例如人员到位、服务使用和融资结构改善。没有任何可观察进展的长期承诺，无法接受公共问责。

% ARGUMENT P279
\phantomsection\label{arg:P279}
长期效果可以通过事先明确的中间环节接受检查。若希望培训改善就业，先应看到有效参与与技能变化；若希望服务减轻照护，先应看到可靠供给、实际使用与时间安排改变。中间环节缺失，不能仅凭长期目标继续推断最终收益。

也要防止指标替代目标。机构为了提高利用率可能优先接收需求较简单者，为了提高就业率可能选择本来容易就业者。应同时观察进入与退出样本，报告被排除者情况，使评估覆盖原本要帮助的人。

证据不充分时，可以保留不确定性并限定规模，而不是在确定成功与完全取消之间二选一。持续学习需要预算纪律，也需要允许合理修正。

本章的改革方向不是一套脱离资源约束的愿望清单。它要求把扩大生产能力与收益兑现联系，把风险稳定与损失责任联系，把劳动选择与生活保障联系。每一项安排都应能够回答同一个问题：它减少了多少真实成本，又把剩余成本放到了谁的账户上。

% ARGUMENT P280
\phantomsection\label{arg:P280}
本章各项安排可以相互支持，但不能自动相加。照护改善可能支持劳动参与，保障可携带可能改善岗位匹配，债务处理可能恢复公共服务；它们也会竞争财政、专业人员和组织能力。真正的实施方案必须在具体约束下排序，而不是复制全部方向后声称已经形成完整预算。

对阶层关系的改善，应从可行选择判断：劳动者能否拒绝严重不利条件，小经营者能否及时得到已完成交易的款项，普通家庭能否承受疾病和失业，公共资源配置者与企业控制者是否承担相应责任。这些结果比抽象宣布某一阶层上升更可检验。

政策最终需要接受与理论相同的约束。若预定机制没有发生，应调整方案；若收益建立在遗漏其他群体成本之上，应重新核算。制度改进的可信度，来自真实兑现和可纠错，而不来自承诺范围越来越大。

\chapter*{结语：一套允许被修正的政治经济学}
\addcontentsline{toc}{chapter}{结语：一套允许被修正的政治经济学}
\section*{把同一件事放回完整关系中}
设想一个纯粹用于说明方法的县域产业园。地方建设道路和厂房，银行提供贷款，企业采购机器、招聘工人，家庭期待工资改善。若产品有需求，企业提高效率并及时支付，地方获得可持续收入，家庭获得稳定生活，积累就完成了从投资到公共与私人收益的传递。

% ARGUMENT P281
\phantomsection\label{arg:P281}
继续这个假设例子，还应区分园区内的人。地方决策者决定公共资源用途，企业控制者决定投资与用工，工程人员组织生产，工人以劳动取得工资，小商户围绕园区形成经营，家庭成员提供照护。项目成功可能使这些位置都受益，但每一项收益都经过合同、预算和组织传递，不是厂房建成以后自然同时出现。

评价时可以按时间分开建设、运营和维护。建设期的岗位与采购不能代替长期经营回报，长期公共价值也不能掩盖建设中已经发生的拖欠。完整叙述把不同阶段放在同一过程里，既保留形成能力的可能，也允许及时发现错配。

如果订单不足，企业账面仍可能依靠应收款确认收入，地方仍可能按已完成投资报告增长，银行通过展期保持贷款正常，家庭却开始面对工资迟发和照护支出。四个部门各自使用一个看似合理的指标，合起来却可能掩盖同一个现金流缺口。此时问题不是“有没有增长”这一句判断，而是哪些生产已经形成可用财富，哪些收入尚未实现，哪些风险正在转移。

% ARGUMENT P282
\phantomsection\label{arg:P282}
现金缺口如果沿付款链传递，强弱关系就会影响顺序。能够要求预付款的供应者可能较早收回资金，缺少替代客户的小企业可能继续垫付，普通工人和家庭则可能以等待、借款或减少支出承受结果。每个主体都有自己的预算约束，但这些约束之间并不对称。

研究的任务不是仅给某个主体归责，而是查明谁能够改变决定、谁取得了收益、谁承担了未预期成本。责任可以复杂，却不能因为复杂而从分析中消失。将这些联系公开，也有助于设计减少总损失的处置，而非让风险继续传向最不容易表达的一方。

反过来，不能因为项目出现困难，就否认全部公共投资和产业政策。道路可能服务其他企业，厂房可以改造，部分技术与劳动技能可以保留。正确处理需要把可持续能力与已经发生的损失分开，将必要服务与原有收益要求分开。整体分析应帮助这种区分，而不是把一切好处或坏处归给一个抽象主体。

% ARGUMENT P283
\phantomsection\label{arg:P283}
是否继续项目，应比较从当前起的可用价值与新增成本。已经建成的道路可以服务其他活动，专用设备可能缺少替代用途，培训形成的技能又可能转向其他企业。把它们分别评估，可以避免因为局部失败否定全部能力，也避免为了保护过去投入继续扩大损失。

对劳动者和居民，保留能力应包括健康、技能和服务连续性。资产重组如果只关注土地和机器，而让人员长期失去收入与训练，可能损坏更难重建的部分。因此，经济调整中的人不是资产处置完成以后的附带问题，而是生产能力本身的组成。

这个例子说明本书三本账的实际用途。生产账防止把资产价格当成新增财富，权利与收入账说明谁能够决定与索取，再生产与风险账说明谁维持生活和吸收失败。三本账不能直接相加，却必须相互核对。

% ARGUMENT P284
\phantomsection\label{arg:P284}
三本账的核对不是创造一个包罗一切的总分。生产账可以记录实物与服务，权利与收入账可以记录决定权限和现金，再生产与风险账可以记录时间、保障与损失；其中有些指标能够货币化，有些适合保持原单位。强行全部相加，反而会掩盖无法互相替代的基本需要。

同一工资支付会出现在企业成本和家庭收入中，这是一项转移关系，不能在社会净收益中重复计数；一次事故补偿改变损失分担，但并不使健康损害消失。账户之间的连接，帮助我们说明哪些是资源增加，哪些是权利变化，哪些只是成本换了承担者。

\section*{哪些经典命题被保留}
本文保留商品背后存在社会关系、劳动合同连接交换与控制、资本需要经过生产和销售完成循环、积累会改变组织与社会位置，以及制度具有历史性的研究方向。它们能够持续提出有价值的问题，而不要求每项现实价格都符合一条单一公式。

% ARGUMENT P285
\phantomsection\label{arg:P285}
保留经典问题意识，不要求保留每个时代形成的全部答案。商品背后的协作与权力、劳动合同中的控制、生产与销售的时间差、积累对社会位置的重塑，都仍可以提出具体研究问题。它们是否解释某个中国现象，需要在制度和资料中逐项判断。

对原典也应公平。马克思关于价值与生产价格、利润与竞争的讨论，比“每个价格严格等于劳动时间”的口号复杂。本书批评的是不能处理条件与中介的强版本解释，不把被简化的口号当成原著全部内容。理论更新应建立在准确重述之后，而不是先把旧理论压缩成容易推翻的对象。

本文修订劳动价值作为普遍价格决定论的强版本，把劳动投入、市场价格和社会贡献分开；修订非工资收入的同质化解释，区分必要维护、创新、风险、中介服务与租金；修订单一阶级标签，引入控制、收益、保障和风险轴线；修订必然危机的机械推断，加入部门账户、现金流和制度反馈。

% ARGUMENT P286
\phantomsection\label{arg:P286}
这些修订改变的是推论的范围。劳动投入仍是研究生产和生活的重要维度，但市场价格还受稀缺、竞争、权利与需求影响；剩余分配仍需分析控制关系，但会计余额不能未经分解就等同于某一种社会贡献或损害；阶级仍是关系概念，但其现实位置需要同时观察控制、保障与家庭条件。

危机理论也从单一路径转向条件反馈。债务、需求、利润和生态压力可以相互放大，也可能被有效制度缓冲。允许不同路径，并非取消解释，而是要求更明确地说明何时、通过什么环节、对哪些人发生作用。只有这样，理论才可能指导观察并接受反证。

尚不能判断的，是这些修订在中国不同地区和行业的量化权重，以及三组候选命题相对于竞争模型的额外预测能力。没有新微观数据和严格识别，本书不能声称已经解决这些问题。明确未解决之处，是理论工作的组成部分。

% ARGUMENT P287
\phantomsection\label{arg:P287}
本书的经验工作限于可取得的公开统计、制度文件、文献与复算。没有受限家庭微观样本，也没有新的调查、访谈或因果回归；一些文献只核验到摘要和公开章节，不能承担超出其阅读层级的细节推论。这些边界在附件中保留，以便后续研究继续推进。

下一步最有价值的工作，是选择一项范围明确的制度变化，匹配组织、家庭与公共账户，检验本文是否增加了解释。若新的资料推翻某项判断，应修改文字和机制，而不是只增加例外。研究专著的完成意味着形成了可讨论、可追踪的论证，不意味着有关问题已经被最终解决。

\section*{中国特色应当体现为什么}
中国特色不应只是把普遍概念换成中国词语，也不应被当作不接受比较和反证的理由。它应体现为具体制度关系：公共产权与财政传递、中央地方组织、土地权利与住房、户籍与人口流动、家庭照护、银行信用以及制造业全球协作。这些关系既能创造能力，也会形成特定约束。

% ARGUMENT P288
\phantomsection\label{arg:P288}
中国特色应当落在可以比较的制度连接上。例如公共资本收益怎样进入预算与服务，地方支出责任怎样连接人口流动，土地与住房怎样影响家庭资产和融资，银行协调怎样改变风险时间，亲属照护怎样支持劳动。这些联系具有具体历史，又可以与其他制度安排比较。

比较不是寻找一个外部标准后机械打分。不同地区和国家的资源、人口与发展阶段不同，应先识别可比问题，再观察不同组织方式怎样处理约束。不能因存在差异而拒绝一切比较，也不能因使用相同指标就假定条件完全相同。

在不同历史阶段，同一制度可能发挥不同作用。原先有助于迅速积累的安排，可能在需求和人口结构改变后增加维护与风险成本；原先看似主要属于福利的服务，在新的技术与就业结构下也可能成为提高效率的条件。政治经济学的任务，是解释这种变化，而不是永久固定赞美或批评。

% ARGUMENT P289
\phantomsection\label{arg:P289}
制度作用会随环境改变，因为它处理的主要约束也会改变。基础设施不足时，快速组织建设可能形成较高收益；存量规模扩大以后，维护、需求和债务管理可能更重要。家庭能够以亲属时间承担照护的条件，也会随人口迁移、就业和健康改变。

这要求评价保持历史性。同一安排在过去有效，不保证以后仍以同样方式有效；今天出现问题，也不证明过去全部投入没有价值。政策学习需要解释转折发生在哪里，而不是把一个阶段的判断永久投射到全部历史。

本书的理论综合把“资本增殖”放在更完整的制度循环中：资源被组织为生产能力，权利决定收益传递，家庭与公共服务维持再生产，风险分担影响下一轮选择。这个循环没有预设必然成功或失败。它要求持续核查事实，也要求政策公开代价。

% ARGUMENT P290
\phantomsection\label{arg:P290}
这个循环可以被具体行动改变。及时付款改善企业与家庭周转，服务可靠减少被迫退出劳动，透明规则改善交易与协商，明确损失责任改变下一轮投资。它们未必消除所有冲突，却可以减少因信息、控制与责任分离造成的额外成本。

因此，理论综合的实践用途是帮助看见中间环节。增长怎样成为收入，收入怎样支持生活，失败怎样影响下一轮选择，都需要继续追踪。若读者能够据此识别某个政策承诺省略了谁的成本，或某项统计结论跳过了什么口径，框架便已经提供了可检验的价值。

读者可以从不同价值立场提出不同改革排序，但至少应共享几项约束：不能把预测写成实绩，不能把账面变化重复计算，不能把一组相关关系直接叫作因果，也不能让没有议价能力的群体承担的成本从分析中消失。新的政治经济学如果有意义，应首先提高人们识别这些问题的能力。

% ARGUMENT P291
\phantomsection\label{arg:P291}
对社会结构的描述，最终应当允许真实的人保留差异与选择。职业、收入、年龄或居住地都不能替个人决定其全部利益与立场；但个人差异也不应使持续存在的资源与权力不对称变得不可讨论。尊重两者，才能同时避免标签化和去结构化。

本书将结论保持开放：让生产形成可用能力，让收益经过清楚规则进入生活，让控制与责任能够对照，让遭遇失败的人仍有重新开始的基本条件。至于在特定约束下怎样排序、怎样筹资、怎样修正，需要继续依靠公开证据与实际结果。新的政治经济学应当使这种讨论更准确，而非用新的标题结束讨论。


\appendix
\chapter{概念、核算与理论修订表}
\section{本书使用的概念边界}
\begin{longtable}{p{29mm}p{57mm}p{57mm}}
\toprule 概念 & 本书含义 & 常见误用\\\midrule\endhead
社会财富 & 满足需要的物品、服务、能力及其持续条件 & 用资产价格或GDP单独穷尽全部福利\\
劳动投入 & 直接及产业链间接劳动；需说明工时和技能调整 & 用工资总额直接代替劳动时间\\
市场价格 & 具体条件下实现的货币交换比例 & 视为社会贡献或道德应得的直接量度\\
增加值 & 总产出扣除中间投入 & 把产业链每次销售全部相加\\
会计利润 & 按会计规则确认收入与费用后的余额 & 等同现金流或全部社会剩余\\
经济租金 & 超过维持相应资源供给所需回报的部分；依赖反事实 & 任何正利润都称租金\\
收入 & 一定期间形成的流量；具体口径另注 & 将借款本金当收入\\
净财富 & 一定时点资产减负债；估值有不确定性 & 把未实现升值当作当期生产\\
实际控制 & 对资源、合同和重大决策的实际影响能力 & 单凭股权登记推断所有权限\\
公众可兑现收益 & 公共资产经治理与分配转化为实际服务和现金利益 & 把公共总资产平均分给居民作为可用钱财\\
再生产能力 & 维持健康、照护、教育与持续生活的资源和时间条件 & 只数劳动力人数，不看维持条件\\
风险承担 & 冲击后的最终经济损失归属 & 把名义债务、减值和救助重复相加\\
金融脆弱性 & 收入、到期承诺、抵押与融资渠道的组合状态 & 单凭高债务比率判定必然违约\\
技术暴露度 & 任务可能被某项技术影响的程度 & 等同确定失业比例\\
\bottomrule
\end{longtable}

同一个术语在马克思文本、国民核算和现代公司财务中可能有不同定义。本书在跨体系讨论时优先说明对应关系，不把词语相同当作指标相同。例如劳动者获得的工资性收入、企业支付的劳动报酬和国民收入中的劳动份额，包含范围并不一致。资本也可以指物质生产资料、金融索取权或一套社会关系，使用时应说明分析层次。

\section{核心理论的处理结果}
\begin{longtable}{p{27mm}p{61mm}p{55mm}}
\toprule 主题 & 保留、修订或暂不判断 & 检验与适用边界\\\midrule\endhead
商品与价值 & 保留使用与交换、社会劳动配置的区分；不接受未经检验的劳动价格严格比例 & 可复制产品、质量调整与投入产出表；稀缺位置和垄断另行解释\\
劳动与劳动力 & 保留合同交换与实际劳动控制的区别 & 岗位替代、保障、协商和退出条件不同，结果不同\\
剩余价值 & 保留剩余形成及控制问题；不把会计利润直接等同经典剩余 & 先统一增加值、折旧、混合收入与融资口径\\
资本积累 & 保留竞争、扩张与结构变化；区分新能力、旧资产交易和重估 & 需求、维护、技术、退出与公共任务影响结果\\
货币与信用 & 将银行创造存款与金融约束纳入循环 & 贷款创造购买力不等于创造真实资源；关注清算和偿付\\
循环与再生产 & 保留生产、销售与周转；加入家庭、公共服务和生态条件 & 不同账户不能重复相加；时间与资源约束需独立记录\\
利润与利息 & 分解组织、创新、风险、中介成本与租金；普遍权重暂不判断 & 需要竞争、成本及风险资料，不能只看收入名称\\
地租 & 保留稀缺条件与权利索取；加入规划、公共投资、土地与财政 & 地方和区位差异显著；使用功能与资产重估分开\\
阶级结构 & 保留关系性；采用劳动依赖、控制、资产收益、风险四轴 & 是待验证分类框架，不宣称已测出各类全国人数\\
技术变化 & 替代效应与新任务、需求及议价共同决定结果 & 暴露度、采用、效率和就业依次检验，不能跳步\\
利润率趋势 & 保留资本强度与盈利关系问题；放弃无条件机械推断 & 资本价格、利用率、分配与需求可改变方向\\
危机 & 用部门错配、现金流与反馈说明脆弱性 & 不能据模型宣称某年必然总危机；需要反证与时点信息\\
\bottomrule
\end{longtable}

“放弃”针对强命题的普遍适用，不意味着否认任何特定样本中的相关性。“保留”针对研究方向，也不意味着所有历史论证自动正确。“暂不判断”针对当前证据不能识别的量化关系，不能改写为“原则上已经证明，只差数据”。这些用语构成全书理论立场的边界。

\section{三个候选命题的识别设计}
命题一是产权与收益兑现之间存在可观察的传递差异。理想研究以公共企业收益分配或服务规则的可识别变化为处理，比较相近地区的服务质量、居民获得、企业投资和或有负债。需要检查改革是否与其他财政支持同时发生，是否只选择经营较好的企业，以及受益群体是否构成变化。若联合指标不比资产规模和地区收入更有解释力，就应降低命题的独立地位。

命题二是家庭再生产条件影响生产扩张的需求承接与劳动配置。理想研究利用保障或照护供给的分期实施，跟踪家庭就业、时间、消费和冲击后恢复。竞争解释包括人口、收入趋势、住房与偏好。若实际服务得到改善而上述结果仍无系统变化，或变化全部由独立收入冲击解释，就应承认命题在该范围内支持不足。

命题三是决策权与最终损失责任的分离影响事前选择。研究需要在冲击之前记录担保、救助预期与项目特征，在冲击之后记录真实损失及承担者，避免用事后得到救助倒推事前一定预期救助。若可转移风险的主体并未增加冒险，可能是问责、声誉或监管约束发挥作用；这些观察属于边界条件，而非可以忽略的例外。

三项研究都应预先确定主要结果、比较组和停止规则，报告失败与无显著变化。观察研究不能完全替代随机实验，随机试点也不能自动外推到全国。本文提出的是研究路线，附件没有任何伪装成已运行的回归表。

\chapter{数据口径、复算与核查记录}
\section{资料截止与版本规则}
资料发布截止为北京时间2026年10月14日，扩充版的再次下载与复核日期为10月15日；复核日期不扩大资料发布窗口。本书采用本次可公开取得并核验的相应期别：价格为9月，货物贸易为前三季度，居民收入为上半年，财政和所核实金融报告为前八个月，人口与产业年度结构为2025年，时间利用调查为2024年。各指标发布日期不同，没有把它们统称为一张“2026全年”截面。

截稿当天仍可能有分时发布，搜索索引也可能滞后。金融报告未获得央行主站的同文链接，使用标注人民银行来源的新华社完整转载，并记录这一出处层级；海关当日数字使用新华社发布，未声称已下载海关全套原表。这些限制不靠增加转载网站数量消除。后续如获得原表，应逐项核对而不静默改数。

IMF二月评估和世界银行七月报告用于说明当时的分析与竞争性解释，没有把其增长预测当成2026年实绩，也没有把旧预测称作十月最新预测。BIS用于债务统计定义；ILO用于职业暴露方法。WID国别表为2024估计并有内部一致性问题。国外资料的使用目的不同，不能把它们排列成一组声称相互独立的现状验证。

年度初步核算会修订，家庭调查存在抽样误差和覆盖不足，国际估计还可能调整顶端分布与国民账户。附件保留来源、发布期、观察期、单位、范围及处理方法；重新取得更新数据时应增加版本，而非把原观察期替换为下载日期。

\section{关键口径冲突怎样处理}
第一，国民收入与居民可支配收入并不相同。WID税前国民收入的分配估计不能与住户调查的可支配收入均值直接拼接，不能从两者差额推断“隐藏收入”总额。需要先处理部门、税收、转移和未分配收益。

第二，社会融资规模存量不是纯债务。它包含股票融资，其他分类也需要对应工具与部门。银行贷款统计中的企事业单位包含机关团体，其增量不能直接等同生产性企业新增投资。BIS与国内统计的部门定义须逐项匹配。\cite{BISCredit}

第三，时间利用的参与者均值不等于全部居民均值。如果参与率为$q$，参与者平均时间为$h$，在未参与者时间记零且口径一致时，全体均值为$qh$。男女参与率和活动分类不同，不能仅用参与者时间相减就推断全体人口时间差。调查还涉及日记日与加权规则。\cite{TimeUseSample}

第四，名义零售、投资与贸易增速不能直接和实际工业增长比较差额。价格、覆盖范围和核算边界不同。本文只用并列观察提出机制问题，不据此估计一个精确的全国供需缺口。

第五，公共预算本级收入与含转移支付后的可用财力不同；出让收入与净收益不同；债务置换与债务减免不同。财政分析必须说明合并边界，不能以一张全国收入支出表替代全部地方现金流。

\section{可复算的例子}
第二章两部门例子的劳动投入为0.625和0.5，价格为约0.657895和0.565789。劳动比例1.25与价格比例约1.162791不同。示例只证明所列假设下不必然成比例，不估计现实中国价格与劳动的相关程度。

第六章的置换例子为100单位本金乘以利率从5\%降至3\%的差，得到当期利息减少2单位。计算没有消除本金，没有计入担保价值，也没有预测未来利率。

第十一章两条债务曲线使用相同初值与零基本赤字，按固定名义增长和利率计算五期。可以通过改变参数检验敏感性，但不能把曲线名称改为乐观或悲观“中国预测”。缺少初始真实债务口径和行为反应，它们只是教学情景。

第八章平均数相对中位数的差为$(22981/19036-1)\times100\%$，约20.7\%；贸易差额为$23.20-16.74=6.46$万亿元。它们分别是分布的有限描述量和货物差额，不能转换成基尼系数或经常账户余额。

\section{发现的问题及处置}
WID中国国别表的底部组均值、份额与组间倍数不能在所标注的人口分组下同时满足加权恒等式。核查保留官方PDF及页面图像。扩充修订版撤下精确份额比较图，原列数值仅在来源审计中照录，不进入精确分配结论、阶层规模或政策测算；未擅自修改原表。待发布方提供一致版本或口径说明后再评估采用范围。

CHFS数据申请范围不等于最新完成调查范围，本文不声称取得尚未公开申请的样本。旧调查只作为历史证据，不由“2026版”书名获得新的时效性。\cite{CHFS}

若某文献仅能读取公开摘要和书目信息，评议限定于其明示机制和研究范围，不声称重做全文估计。来源索引逐条记录阅读层级。本书没有进行穷尽式系统文献综述，也未证明其概念组合在全球研究中首次出现；“理论综合”是有意保持的贡献定位。

图表共12幅，其中8幅为发布数据图，3幅为机制或风险分担示意，1幅为债务假设曲线。产业、价格、工资、信贷、财政、照护、收入与贸易图分别标注期别和分母；社会位置图、公共收益传递图及风险分担图均不使用伪造的经验系数。WID原份额图已撤下，相关原列数值仅保留作来源审计。

\section{附件与独立编译}
主文件内嵌全部图表数值、绘制指令和参考文献，无须外部图片、数据文件或BibTeX数据库即可编译。研究附件保存来源索引、关键数据台账、情景数据、计算检查、文献评议、争议记录及复算脚本。附件用于复核与后续修订，不是主文件编译依赖。

脚本只完成已声明的算术、图表生成和结构检查。它们不下载受限调查，不调用隐藏数据，也不产生因果估计。编译状态、导出状态和视觉检查单独记录，不能以源码存在代替PDF交付，也不能以PDF生成代替研究结论得到证明。

\chapter{扩充版的识别、复算与段落审计}
\section{原稿段落与扩写对应}
扩充前的正文保存为研究底稿，并建立291个实质性段落编号。编号包含导论、十二章和结语，个别段落同时含章或节标题；图表定义、参考文献与技术附录不进入该编号范围。扩充版逐一保留对应关系，修订的原段另记原文校验值与处理说明，不通过复制旧图表增加正文篇幅。

附件的段落对应表记录原段、所属章节、扩写锚点、行号与新增汉字数。行号随修订变化，锚点以arg:P001至arg:P291形式保留在源码中。它们用于追踪论点，不表示每个原段都必须套用同一个写作模板。正文按实际问题分别展开概念、历史、机制、证据、阶层后果和竞争解释。

篇幅采用保守汉字计数：只统计正文叙述中的汉字，不计标点、数字、英文、注释、标签、引用键、绘图代码、封面、参考文献或附录；正文内的图注和表格另列而不并入严格正文计数。该数用于验收10至15万字范围，不把带格式源码长度当作正文长度。

\section{WID原表的一致性复算}
设某组人口占比为$q_j$，组平均收入为$\bar y_j$，全体均值为$\bar y$，同口径收入份额应满足
\[
s_j=q_j\frac{\bar y_j}{\bar y}.
\]
原表为《世界不平等报告2026》国别表PDF第9页、印刷页166的中国表。按其原列数字：
\[
0.5\times\frac{1988}{14396}\times100\%\simeq6.90\%,\qquad
\frac{43.3/0.1}{13.7/0.5}\simeq15.80.
\]
它们分别与原列底部50\%收入份额13.7\%及组间倍数31.4不一致。财富列的同类复算为
\[
0.5\times\frac{5447}{86462}\times100\%\simeq3.15\%,
\]
与表列6.3\%不同。此处只审计原列数字；6.90\%、15.80和3.15\%都不是本书提供的中国分配替代估计。若人口单位、排序或其他口径另有解释，仍需发布方说明。扩充版因此暂停使用该页作精确比较，原始页面与复算结果随附件保存。\cite{WID2026}

\section{技术采用、需求与总工时的边界}
设同一质量任务的单位工时为$h$，任务量为$Q$，总工时为$L=hQ$。在任务构成和质量不变的简化条件下，若单位工时变为原来的$a$倍，则
\[
\frac{L_1}{L_0}=a\frac{Q_1}{Q_0}.
\]
第九章十小时变六小时的示例中，$a=0.6$，任务量需要增加到$1/0.6\simeq1.67$倍，才维持相同总工时。公式没有包含新任务、质量、审核、维护与培训，也没有从总工时直接推出人数或工资。完整采用评价应将这些变化另行记录，避免把局部节省当作全流程净收益。

\section{损失转移与真实损失的区分}
若同一项已确认损失为$\ell$，初始由各主体承担的金额为$b_i$，且$\sum_i b_i=\ell$，一次纯粹分担调整可写为
\[
b_i'=b_i+\sum_j T_{ji}-\sum_j T_{ij},\qquad
\sum_i b_i'=\ell.
\]
这里$T_{ij}$表示由$i$向$j$转出的同一损失承担额。只要没有新增真实损失或回收，分担调整不改变总和。财政补偿债权人以后，不能把原债权损失与同额财政支出当作两项独立社会资源损失相加；但融资、行政、停工与健康等附加成本可能确实发生，必须单独识别。

纯粹转移与减少损失也不同。恢复有价值项目的运行可能增加回收，及时支付可能避免供应商破产，安全投入可能减少事故；这些变化会改变$\ell$本身。风险矩阵要同时记录总损失、承担份额和产生变化的机制，不能只用一张转移表证明政策具有净收益。

\section{候选命题的最低检验要求}
\begin{longtable}{p{.18\textwidth}p{.27\textwidth}p{.23\textwidth}p{.20\textwidth}}
\toprule 命题 & 所需观测 & 竞争解释 & 支持不足的情形\\\midrule\endhead
公共收益传递 & 相近任务下的治理、预算去向、维护、服务与担保 & 行业景气、地区财力、改革对象选择 & 新变量未提供资产规模之外的信息\\
再生产与能力 & 实际服务、家庭时间、就业、支出和冲击恢复 & 人口、住房、收入趋势及偏好 & 服务兑现后预期中间环节仍未出现\\
控制与损失责任 & 冲击前权限与预期，冲击后回收及承担者 & 问责、声誉、监管与不同项目风险 & 责任分离未增加相关冒险或转移\\
平台依赖与风险 & 替代渠道、净费率、申诉、退出与家庭损失 & 质量、欺诈、需求及商家构成 & 高依赖条件下仍有有效协商与广泛传递\\
\bottomrule
\end{longtable}

这些要求构成后续研究设计，没有在本书中被伪装成已经完成的实验。比较前后变化需要处理共同冲击和样本退出，观察相关需要处理反向因果，随机试点也需要处理外推边界。明确可观察的失败条件，才能使理论、政策和资料在后续修订中真正互相约束。

\section{新增来源的阅读层级}
阶层部分结合已取得的《中国社会各阶层分析（新版）》工作底稿，但该底稿不作为独立统计验证来源。国内分层研究回顾用于比较分类框架；Wright公开章节用于核对中间位置的概念，不据此估计中国阶层规模。AI客服与增加值贸易的正式刊物入口本次受访问限制，采用同一作者的NBER公开工作论文摘要核对机制和版本，未混用不同版本效果量。职业伤害保障采用吉林公开实施说明支持省级事实，全国实际覆盖未从该省材料外推。

资料在2026年10月15日重新读取，不代表使用该日以后新发表材料。受验证页面、权限或缺少微观数据限制的环节均保留在来源与争议表中。所附脚本执行确定性复算、映射、排版核查和内置编译调用，没有自动形成新的经验估计。

% BIBLIOGRAPHY_BEGIN
\backmatter

\renewcommand{\bibname}{参考文献与数据来源}

\cleardoublepage\phantomsection

\addcontentsline{toc}{chapter}{参考文献与数据来源}

\begin{thebibliography}{99}

\small\sloppy

\bibitem{Marx1} 马克思. 资本论，第一卷. 1867. 商品、劳动过程、剩余价值、积累诸章。\href{\detokenize{https://www.marxists.org/chinese/marx/capital/index.htm}}{原始来源}。访问：2026-10-14。

\bibitem{Marx2} 马克思. 资本论，第二卷. 1885. 资本循环及社会总资本再生产。\href{\detokenize{https://www.marxists.org/chinese/marx/capital/marxist.org-chinese-marx-capital-vol2-01.htm}}{原始来源}。访问：2026-10-14。

\bibitem{Marx3} 马克思. 资本论，第三卷. 1894. 利润、利息、地租、利润率变化。\href{\detokenize{https://www.marxists.org/chinese/marx/capital/marxist.org-chinese-marx-capital-vol3-01.htm}}{原始来源}。访问：2026-10-14。

\bibitem{NBS2025} 国家统计局. 中华人民共和国2025年国民经济和社会发展统计公报. 2026. 一、综合；人口表；工业、投资等分节。\href{\detokenize{https://www.stats.gov.cn/sj/zxfb/202602/t20260228_1962662.html}}{原始来源}。访问：2026-10-14。

\bibitem{Income2026} 国家统计局. 2026年上半年居民收入和消费支出情况. 2026. 一、居民收入；二、居民消费；表1。\href{\detokenize{https://www.stats.gov.cn/sj/zxfbhjd/202607/t20260715_1964129.html}}{原始来源}。访问：2026-10-14。

\bibitem{IMF2026} International Monetary Fund. People’s Republic of China: 2025 Article IV Consultation, Country Report 26/044. 2026. 摘要及宏观政策评估。\href{\detokenize{https://www.imf.org/en/publications/cr/issues/2026/02/17/peoples-republic-of-china-2025-article-iv-consultation-press-release-staff-report-and-574028}}{原始来源}。访问：2026-10-14。

\bibitem{WB2026} World Bank. China Economic Update, July 2026: Rebalancing Growth. 2026. 发布说明与报告概要。\href{\detokenize{https://www.worldbank.org/en/news/press-release/2026/07/07/rebalancing-growth-china-economic-update}}{原始来源}。访问：2026-10-14。

\bibitem{CPI2026} 国家统计局. 2026年9月份居民消费价格同比上涨0.9\%. 2026. 正文首段与分类表。\href{\detokenize{https://www.stats.gov.cn/sj/zxfbhjd/202610/t20261014_1965477.html}}{原始来源}。访问：2026-10-14。

\bibitem{PPI2026} 国家统计局. 2026年9月份工业生产者出厂价格同比上涨5.0\%. 2026. 首段、一、同比变化及分类表。\href{\detokenize{https://www.stats.gov.cn/sj/zxfbhjd/202610/t20261014_1965476.html}}{原始来源}。访问：2026-10-14。

\bibitem{Zhou2007} 周黎安. 中国地方官员的晋升锦标赛模式研究. 2007. 经济研究，42(7):36–50。\href{\detokenize{https://ir.pku.edu.cn/handle/20.500.11897/73869}}{原始来源}。访问：2026-10-14。

\bibitem{Xiong2018} Wei Xiong. The Mandarin Model of Growth. 2018. NBER Working Paper 25296。\href{\detokenize{https://www.nber.org/papers/w25296}}{原始来源}。访问：2026-10-14。

\bibitem{CompanyLaw} 全国人大常委会；市场监管总局公布. 中华人民共和国公司法（2023年修订）. 2023. 第七章，168—170条。\href{\detokenize{https://www.samr.gov.cn/zw/zfxxgk/fdzdgknr/fgs/art/2023/art_067c072db6ef4679a2e0180996be4cf8.html}}{原始来源}。访问：2026-10-14。

\bibitem{Kornai2003} János Kornai, Eric Maskin, Gérard Roland. Understanding the Soft Budget Constraint. 2003. Journal of Economic Literature, 41(4):1095–1136。\href{\detokenize{https://pubs.aeaweb.org/doi/abs/10.1257/002205103771799999}}{原始来源}。访问：2026-10-14。

\bibitem{Gotha} 马克思. 哥达纲领批判. 1875. 第一部分。\href{\detokenize{https://www.marxists.org/archive/marx/works/1875/gotha/ch01.htm}}{原始来源}。访问：2026-10-14。

\bibitem{Migrant2025} 国家统计局. 2025年农民工监测调查报告. 2026. 一、规模；二、基本特征；三、就业；调查说明。\href{\detokenize{https://www.stats.gov.cn/sj/zxfb/202604/t20260430_1963472.html}}{原始来源}。访问：2026-10-14。

\bibitem{Acemoglu2019} Daron Acemoglu, Pascual Restrepo. Automation and New Tasks: How Technology Displaces and Reinstates Labor. 2019. Journal of Economic Perspectives, 33(2):3–30。\href{\detokenize{https://pubs.aeaweb.org/doi/abs/10.1257/jep.33.2.3}}{原始来源}。访问：2026-10-14。

\bibitem{Grossman1986} Sanford J. Grossman, Oliver D. Hart. The Costs and Benefits of Ownership: A Theory of Vertical and Lateral Integration. 1986. Journal of Political Economy, 94(4):691–719。\href{\detokenize{https://www.journals.uchicago.edu/doi/10.1086/261404}}{原始来源}。访问：2026-10-14。

\bibitem{Song2011} Zheng Song, Kjetil Storesletten, Fabrizio Zilibotti. Growing Like China. 2011. American Economic Review, 101(1):196–233。\href{\detokenize{https://www.aeaweb.org/articles?id=10.1257/aer.101.1.196}}{原始来源}。访问：2026-10-14。

\bibitem{Hsieh2009} Chang-Tai Hsieh, Peter J. Klenow. Misallocation and Manufacturing TFP in China and India. 2009. Quarterly Journal of Economics, 124(4):1403–1448。\href{\detokenize{https://academic.oup.com/qje/article-abstract/124/4/1403/1917179}}{原始来源}。访问：2026-10-14。

\bibitem{BOE2014} Michael McLeay, Amar Radia, Ryland Thomas. Money creation in the modern economy. 2014. Quarterly Bulletin 2014 Q1。\href{\detokenize{https://www.bankofengland.co.uk/quarterly-bulletin/2014/q1/money-creation-in-the-modern-economy}}{原始来源}。访问：2026-10-14。

\bibitem{Money2026} 中国人民银行. 2026年8月金融统计数据报告（新华网完整转载）. 2026. 一、社融存量；五、人民币贷款；注释。\href{\detokenize{https://www3.xinhuanet.com/fortune/20260914/06447ad39fda407a8da93deb2a607b64/c.html}}{原始来源}。访问：2026-10-14。

\bibitem{BISCredit} Bank for International Settlements. Credit to the non-financial sector: documentation. 2026. 部门、工具与估值定义。\href{\detokenize{https://www.bis.org/statistics/totcredit/credpriv_doc.pdf}}{原始来源}。访问：2026-10-14。

\bibitem{Minsky1992} Hyman P. Minsky. The Financial Instability Hypothesis. 1992. Levy Economics Institute Working Paper 74。\href{\detokenize{https://www.levyinstitute.org/publications/the-financial-instability-hypothesis/}}{原始来源}。访问：2026-10-14。

\bibitem{LandLaw} 全国人大常委会. 中华人民共和国土地管理法（2019年修正）. 2019. 全国人大常委会公报2019年第5号；第九条。\href{\detokenize{https://wb.flk.npc.gov.cn/flfg/PDF/2476cfe9607a43e4bab90d1b10624dbc.pdf}}{原始来源}。访问：2026-10-14。

\bibitem{Fiscal2026} 财政部国库司. 2026年1—8月财政收支情况. 2026. 一、一般公共预算；二、政府性基金预算。\href{\detokenize{https://gks.mof.gov.cn/tongjishuju/202609/t20260918_3997709.htm}}{原始来源}。访问：2026-10-14。

\bibitem{Wen2021} 温铁军（Wen Tiejun）. Ten Crises: The Political Economy of China’s Development (1949–2020). 2021. 导论与城乡成本转移章节。\href{\detokenize{https://link.springer.com/book/10.1007/978-981-16-0455-3}}{原始来源}。访问：2026-10-14。

\bibitem{Fraser2016} Nancy Fraser. Contradictions of Capital and Care. 2016. New Left Review, 100:99–117。\href{\detokenize{https://newleftreview.org/issues/ii100/articles/nancy-fraser-contradictions-of-capital-and-care}}{原始来源}。访问：2026-10-14。

\bibitem{TimeUse2024} 国家统计局. 第三次全国时间利用调查公报（第三号）. 2024. 二、无酬劳动；不同性别参与者时间。\href{\detokenize{https://www.stats.gov.cn/zt_18555/zthd/lhfw/2025/2025_qgsjlydc/202501/t20250127_1958534.html}}{原始来源}。访问：2026-10-14。

\bibitem{PYZ2019} Thomas Piketty, Li Yang, Gabriel Zucman. Capital Accumulation, Private Property, and Rising Inequality in China, 1978–2015. 2019. American Economic Review, 109(7):2469–2496。\href{\detokenize{https://www.aeaweb.org/articles?id=10.1257/aer.20170973}}{原始来源}。访问：2026-10-14。

\bibitem{WID2026} World Inequality Lab. World Inequality Report 2026, Country sheets: China. 2026. PDF第9页，印刷页166，Table 1与图注。\href{\detokenize{https://wir2026.wid.world/www-site/uploads/2026/09/WIR26_Country_Sheets.pdf}}{原始来源}。访问：2026-10-14。

\bibitem{CHFS} 西南财经大学中国家庭金融调查与研究中心. 数据申请说明. 2026. 公开申请范围、内部申请条件。\href{\detokenize{https://chfs.swufe.edu.cn/sjzx/sjsq.htm}}{原始来源}。访问：2026-10-14。

\bibitem{Rochet2003} Jean-Charles Rochet, Jean Tirole. Platform Competition in Two-Sided Markets. 2003. Journal of the European Economic Association, 1(4):990–1029。\href{\detokenize{https://academic.oup.com/jeea/article-abstract/1/4/990/2280902}}{原始来源}。访问：2026-10-14。

\bibitem{ILO2025} Paweł Gmyrek 等. Generative AI and Jobs: A Refined Global Index of Occupational Exposure. 2025. ILO Working Paper 140；DOI 10.54394/HETP0387。\href{\detokenize{https://researchrepository.ilo.org/esploro/outputs/encyclopediaEntry/Generative-AI-and-jobs-a-refined/995653520102676}}{原始来源}。访问：2026-10-14。

\bibitem{Injury2026} 人力资源和社会保障部. 新就业形态人员职业伤害保障试点扩围. 2026. 2026年7月1日发布。\href{\detokenize{https://www.mohrss.gov.cn/wap/xw/rsxw/202607/t20260701_579347.html}}{原始来源}。访问：2026-10-14。

\bibitem{Current2026} 国家统计局. 2026年8月份国民经济运行数据. 2026. 一、工业；三、零售；四、投资；六、就业；统计附注。\href{\detokenize{https://www.stats.gov.cn/sj/zxfbhjd/202609/t20260915_1965307.html}}{原始来源}。访问：2026-10-14。

\bibitem{Trade2026} 海关总署数据；新华社发布. 前三季度我国货物贸易进出口数据. 2026. 首段总值及进出口分项。\href{\detokenize{https://www.xinhuanet.com/20261014/60d6cef6c94e4afb9234dddf3bbea5c0/c.html}}{原始来源}。访问：2026-10-14。

\bibitem{TimeUseSample} 国家统计局. 第三次全国时间利用调查公报（第一号）. 2024. 调查对象、抽样与时间日记说明。\href{\detokenize{https://www.stats.gov.cn/sj/zxfb/202410/t20241031_1957217.html}}{原始来源}。访问：2026-10-14。

\bibitem{MarxPrice} 马克思. 资本论第三卷，第九章：一般利润率的形成和商品价值转化为生产价格. 1894. 第九章；公开原文。\href{\detokenize{https://www.marxists.org/archive/marx/works/1894-c3/ch09.htm}}{原始来源}。复核：2026-10-15。
\bibitem{Card2022} David Card. Who Set Your Wage?. 2022. American Economic Review 112(4):1075–1090；原刊摘要与公开稿入口。\href{\detokenize{https://www.aeaweb.org/articles?id=10.1257/aer.112.4.1075}}{原始来源}。复核：2026-10-15。
\bibitem{Wright1985} Erik Olin Wright. Classes: The Biography of a Concept—Contradictory Class Locations. 1985. 第二章；作者高校主页公开章节。\href{\detokenize{https://www.sscc.wisc.edu/soc/faculty/pages/wright/Published%20writing/classes-2.pdf}}{原始来源}。复核：2026-10-15。
\bibitem{Chamon2010} Marcos Chamon, Kai Liu, Eswar Prasad. Income Uncertainty and Household Savings in China. 2010/2013. NBER WP 16565；正式版JDE 105:164–177；1989—2006年家庭资料。\href{\detokenize{https://www.nber.org/papers/w16565}}{原始来源}。复核：2026-10-15。
\bibitem{Bai2016} Chong-En Bai, Chang-Tai Hsieh, Zheng Michael Song. The Long Shadow of a Fiscal Expansion. 2016. NBER WP 22801；摘要所述历史机制，不外推为2026年债务测算。\href{\detokenize{https://www.nber.org/papers/w22801}}{原始来源}。复核：2026-10-15。
\bibitem{Migration2023} Fanghua Li, Chenyang Ji, Moshe Buchinsky. Social Insurance and Migration: Evidence from a Nation-Wide Institutional Reform in China. 2023. NBER WP 31819；摘要及研究设计；局部效应。\href{\detokenize{https://www.nber.org/system/files/working_papers/w31819/w31819.pdf}}{原始来源}。复核：2026-10-15。
\bibitem{Housing1998} 国务院. 关于进一步深化城镇住房制度改革加快住房建设的通知. 1998. 国发〔1998〕23号；第二部分；政府公开历史原文。\href{\detokenize{https://www.beijing.gov.cn/zhengce/zhengcefagui/qtwj/201309/t20130924_776668.html}}{原始来源}。复核：2026-10-15。
\bibitem{LiuTian2018} 刘欣、田丰. 社会结构研究40年：中国社会学研究者的探索. 2018. 江苏社会科学2018年第4期；公开重刊，采用学术框架比较。\href{\detokenize{https://www.cssn.cn/skgz/202209/t20220916_5504623.shtml}}{原始来源}。复核：2026-10-15。
\bibitem{MinWage2014} Yi Huang, Prakash Loungani, Gewei Wang. Minimum Wages and Firm Employment: Evidence from China. 2014. IMF WP 14/184；研究机构摘要及作者公开说明；PDF直连受限。\href{\detokenize{https://www.imf.org/external/pubs/ft/wp/2014/wp14184.pdf}}{原始来源}。复核：2026-10-15。
\bibitem{ValueTrade2014} Robert Koopman, Zhi Wang, Shang-Jin Wei. Tracing Value-Added and Double Counting in Gross Exports. 2014. American Economic Review 104(2):459–494；核算框架。\href{\detokenize{https://pubs.aeaweb.org/doi/abs/10.1257/aer.104.2.459}}{原始来源}。复核：2026-10-15。
\bibitem{AIWork2025} Erik Brynjolfsson, Danielle Li, Lindsey Raymond. Generative AI at Work. 2025. Quarterly Journal of Economics 140(2):889–942；特定客服企业研究。\href{\detokenize{https://academic.oup.com/qje/article/140/2/889/7990658}}{原始来源}。复核：2026-10-15。
\bibitem{TradeAdjust2016} David Autor, David Dorn, Gordon Hanson. The China Shock: Learning from Labor Market Adjustment to Large Changes in Trade. 2016. NBER WP 21906；美国劳动市场调整经验，仅用于比较机制。\href{\detokenize{https://www.nber.org/papers/w21906}}{原始来源}。复核：2026-10-15。
\bibitem{Tax1994} 国务院. 关于实行分税制财政管理体制的决定. 1993. 国发〔1993〕85号；1994年起实施；财政部公开原文。\href{\detokenize{https://yss.mof.gov.cn/zhuantilanmu/cztzgg/zcgz/200806/t20080627_54328.htm}}{原始来源}。复核：2026-10-15。
\bibitem{InjuryJilin2026} 吉林省社会保险事业管理局. 吉林省新就业形态从业人员职业伤害保障试点工作全面启动. 2026. 2026-06-30；7月1日起实施；只能直接证明吉林实施。\href{\detokenize{https://jlsi.jl.gov.cn/sbdt/202606/t20260630_9656905.html}}{原始来源}。复核：2026-10-15。
\bibitem{AIWorkNBER} Erik Brynjolfsson, Danielle Li, Lindsey R. Raymond. Generative AI at Work. 2023. NBER Working Paper 31161；2023年11月修订摘要；正式出版版本为QJE 2025，未混用不同版本的效果量。\href{\detokenize{https://www.nber.org/papers/w31161}}{原始来源}。复核：2026-10-15。
\bibitem{ValueTradeNBER} Robert Koopman, Zhi Wang, Shang-Jin Wei. Tracing Value-added and Double Counting in Gross Exports. 2012/2013. NBER WP 18579；2013年2月修订摘要；核算分解框架，正式版AER 2014。\href{\detokenize{https://www.nber.org/papers/w18579}}{原始来源}。复核：2026-10-15。
\end{thebibliography}
% BIBLIOGRAPHY_END

% Edition 1.0: standalone source, evidence cutoff 2026-10-14.
\end{document}
